\documentclass[11pt,a4paper]{article}
\usepackage[T1]{fontenc}
\usepackage[utf8]{inputenc}
\usepackage{lmodern,amsmath,amssymb,amsthm,mathtools,microtype,mathrsfs}
\usepackage[margin=25mm,top=23mm,bottom=25mm]{geometry}
\usepackage{booktabs,array,longtable,enumitem,xcolor,fancyhdr}
\usepackage[colorlinks=true,linkcolor=blue!45!black,urlcolor=blue!45!black,citecolor=blue!45!black]{hyperref}
\hypersetup{pdftitle={Reading the Alpoge--Claude RH and Dirichlet-GRH proof: working version 0.48},pdfauthor={ChatGPT},pdfsubject={Partial human-readable reconstruction with explicit arithmetic inputs},pdfkeywords={Riemann hypothesis, Lean, theta lifts, prime sums, Arakelov slopes}}
\setlength{\parindent}{0pt}
\setlength{\parskip}{5pt plus 1pt}
\setlist{topsep=3pt,itemsep=3pt,leftmargin=*}
\setlength{\headheight}{14pt}
\pagestyle{fancy}\fancyhf{}
\fancyhead[L]{\small RH and Dirichlet GRH: a reconstruction}
\fancyhead[R]{\small Working version 0.48}
\fancyfoot[C]{\thepage}
\renewcommand{\headrulewidth}{0.3pt}
\numberwithin{equation}{section}
\newtheorem{theorem}{Theorem}[section]
\newtheorem{proposition}[theorem]{Proposition}
\newtheorem{lemma}[theorem]{Lemma}
\newtheorem{corollary}[theorem]{Corollary}
\theoremstyle{remark}\newtheorem{remark}[theorem]{Remark}
\newcommand{\Rea}{\operatorname{Re}}
\newcommand{\HS}{\mathrm{HS}}
\newcommand{\tr}{\operatorname{tr}}
\newcommand{\codim}{\operatorname{codim}}
\newcommand{\covol}{\operatorname{covol}}
\newcommand{\OK}{\mathcal O_K}
\newcommand{\src}[2]{\href{https://github.com/CocoSgt/riemann-hypothesis-lean/blob/7cbf7a0c32fca1badf05931879a37c3f8a2b03b4/src/#1}{\texttt{#2}}}
\newcommand{\status}[1]{\textbf{Status:} #1}
\begin{document}
\begin{center}
{\LARGE\bfseries RH and Dirichlet GRH: new reconstruction chapters}\\[7pt]
{\large Version 0.48 --- 23 November 2026}\\[5pt]
{\large One-character estimates, exact measures, and common compatibility}
\end{center}
\textbf{Scope.} This edition contains the current audit and the new
chapters. It remains a partial human reconstruction. Earlier named
claims are in the cumulative manuscript and preserved component files.
The underlying formal proof is accepted as correct; the exact remaining
human reconstruction boundaries are stated below.
\setcounter{tocdepth}{1}
\tableofcontents
\clearpage
\begingroup\emergencystretch=3em
\section{Version 0.48: one-character estimates, exact measures, and common compatibility}
\label{sec:dependency-audit-v48}

\subsection{Current result and the completed selected implication}

The pinned formal source is accepted as fully checked and correct for
reconstruction. The task is to write out its mathematical argument,
allowing simpler proofs of the same original statements. Version 0.48
adds 89 numbered class-B claims to the 2851 preserved predecessor
claims, for 2940 in all. The claim graph records these human proofs;
it is not a mechanically extracted formal dependency graph.

The selected \texttt{TASK\_\allowbreak ONE\_\allowbreak CHAR\_\allowbreak PROOF} is now complete.
Its exact contract is the implication
\[
 \mathrm{InpZ}\ \wedge\
 \mathrm{AutomorphyUnramifiedCoefficients}
 \quad\Longrightarrow\quad\mathrm{OneChar}.
\]
The already reconstructed \(\mathrm{InpZ}\) record supplies the first
argument. The second argument remains the separate open producer
\texttt{AUT\_\allowbreak UNRAM\_\allowbreak COEF}: for every original unramified-coefficient
datum and fixed \(\iota\), one marked \(g_1\) must supply
\(\mathrm{FinA}\), \(\mathrm{AutB}\) and \(\mathrm{LemW}\)
together. Completing the implication does not construct that datum.

The selected sufficient route retains 121 nodes and 161 edges.
Its whole-obligation register decreases from 26 to 25:
18 class C and 7 class B. There are 18 shared RH/Dirichlet-GRH
obligations, 7 additional Dirichlet-GRH obligations, and no
RH-only obligations in this route. A whole obligation can contain
many substantial producers; these counts are an inventory, not a
percentage of completion.

\subsection{The four new chapters}

\paragraph{OCH1--OCH30: the original one-character estimate.}
The local argument works over every original finite Artin coefficient
algebra. Positive finite cohomology, ordinary flags, obstruction
lifting and an explicit nilpotent twisting argument identify the
actual crystalline component. Its twelve-variable chart has weight
multiplicities \((1,2,3,4,2)\). The original normalized global
deformation ring then gives an integral graded model with exactly
the prescribed numbers and weights of generators and equations.

The point, derivation and lifting fields are proved for that same
model. Strict-locality gives boundedness and adic continuity, while
finite-module and Taylor-division arguments control full formal
series and arbitrary dual-number maps. The actual Selmer-to-cotangent
injection and the finite algebra estimate yield the original trace
bound. Passage to an unramified coefficient ring leaves the fixed
\(\iota\), character values and trace scalar unchanged, and full
coefficient independence transports the result back.

More precisely, under the original \(GD\) and trace-nonvanishing
premises, \(V\!\operatorname{len}\) is finite and there is an actual
finite set \(X\) of representations, irredundant under
\(\mathrm{RepIso}\), for which
\[
 V_n(\Omega,c,e,\psi)
 \le 2\sum_{\rho\in X}
   v_\ell\!\left(\operatorname{tr}\rho(\sigma)-t\right).
\]
Every member has the original \(\mathrm{Adm3}(\iota,e,N)\) or
\(\mathrm{Adm2}(\iota,e,N|\operatorname{disc}K|)\) property.
No additional lower bound on \(\ell\) or \(e\) is imposed.
Isomorphism-irredundancy does not assert that every three-dimensional
member of the second branch is irreducible.

\paragraph{AMD1--AMD17: the full measured-data constructor.}
Canonical tensor splitting gives the exact adelic bridge, rational
points and compact subgroup conditions. The actual closed groups
and semidirect product \(S=N\rtimes H\) give the precise Haar
product identity, including the case of noncompact \(H\).
A Tonelli comparison promotes the source's weak Haar-on-set condition
without adding a hidden regularity hypothesis.

This proves all original \(\mathrm{MGExists}\),
\(\mathrm{MWCore}\) and \(\mathrm{MExists}\) fields on the same
bridge and word frames. It also constructs the whole prescribed-measure
\(\mathrm{MeasW}\), retaining the given infinity measures, global
\(N/H\) measures and every local \(S\) measure. An explicit compact
fundamental domain for the rational unipotent group permits covolume
one while preserving the prescribed infinity, \(H\) and local
\(S\) measures. The later Bessel vector, orbital/fiber identities,
Siegel--Weil ties, full transport and arithmetic count remain distinct.

\paragraph{PUN1--PUN24: principal models and actual moduli points.}
The quadratic norm and Hermitian arguments retain the fixed
Hermitian space and scalar, every finite place, and nonprincipal
projective lattices. They supply every arithmetic \(\mathrm{LatModel}\)
field and the complete \(\mathrm{ModelsOf}/\mathrm{ModelBase}\)
construction on a supplied \(\mathrm{FamilyUnif}\).
The actual center lattice is finitely generated and full in every
rank and component family. The precise different, cusp ideals and
half-trace factor give the entire original \(A=1\)
\(\mathrm{ExpoMatch}\), independent of the cusp basis.

Explicit periods and a canonical primitive Appell--Humbert line
produce the original individual moduli points. The proof includes
the negative Weil sign, every positive torsion level, both markings,
every \(\mathbf C\)-to-\(k\) scalar extension, and
\(\mathrm{KTriv}\) over all \(\mathbf C\)-algebras, including
nonreduced algebras. The original point-equality criterion and
marked-genus nonemptiness follow. The finite holomorphic
\(\mathrm{FamilyUnif}\), its theta normalization and local lifting
are additional original fields. Thus \(\mathrm{N1c}\) remains open.

\paragraph{GCF1--GCF18: exact common compatibility.}
The finite-adelic Hilbert--90 calculation gives the actual finer
quotient \(\mathrm{KDq}\), its congruence and rational maps, and
the relative determinant law on the supplied frames. Full column
extension preserves the given compatible group. The normalized CM
period remains in the existing coefficient field. These constructions
give all \(\mathrm{GroupCompatibilityAll}\) fields on one tuple.

A single original true-volume \(\mathrm{GJStage}\) is selected
before every later \(W,V,S,N\), proving the full
\(\mathrm{GIs2}\) stage theorem. Division transports the exact
comparison map, columns and compatibility tuple. On the selected
D3 class, \(\mathrm{ChosenE}\) already gives
\(\mathrm{ChosenTT}\) and \(\mathrm{ChosenBad}\), hence
\(\mathrm{TopsF}\) and the needed \(\mathrm{Adm2F}\).
The eventual common \(\mathrm{WPt}/\mathrm{TSel}\) realization,
including the given comparison and kernel data, remains a producer.
No whole global-local-gain closure follows merely from the stage
construction or division transport.

\subsection{Foundations, preservation and independent review}

All 89 numbered claims are class-B deductions. The named classical
foundations remain separate inputs, including Hasse--Minkowski,
Appell--Humbert, algebraization/GAGA and the universal dual Picard
property. Their exact statements and uses are recorded in the
component foundation registers; they are not newly invented claim
nodes. In particular, the dual universal property is used on all
schemes, which is needed for the nonreduced coefficient-ring clause.

Every new claim has an independent mathematical and interface review
by an agent other than its author. Review manifests bind the exact
proof file and each reviewed block. Cited source bytes are compared
with the pinned git blobs and included in the archive. All predecessor
claim records and the complete historical proof body are preserved;
mutable predecessor metadata is frozen in
\texttt{history/\allowbreak version\_\allowbreak 0\_\allowbreak 47}.

Draft TeX compilation checks syntax, references and missing characters
in both readers. It is not PDF visual review or formal replay.
This remains a partial human reconstruction of RH and unrestricted
Dirichlet GRH. The following version-0.48 register is the current
scope authority; historical conditional statements retain their
printed hypotheses even where a later chapter supplies them.

\subsection{The complete current selected register}

\paragraph{1. GLOBAL\_\allowbreak DENOMINATOR — C — Actual LinkN local table and binding producer}

Closed(Des $\mathcal C$ 3) LinkN: the actual etaSum is bounded with exact dZ0 and same-host vvTerm, subject only to the literal ShaFinOdd and summability guards.

\emph{Shared objects and choice order.} one FamP; one class $\mathcal A$ and its target-compatible subclass; one coefficient field per family datum; one x and i at each application; original selected rows, words and levels; thresholds chosen before the common family; actual D.nd; actual etaJ and Aj; sha\_\allowbreak eq

\emph{Source.} \href{https://github.com/CocoSgt/riemann-hypothesis-lean/blob/7cbf7a0c32fca1badf05931879a37c3f8a2b03b4/src/Formalism/SharpRestriction/GlobalTheoremInputs.lean#L45-L48}{LinkN}.

\emph{Source.} \href{https://github.com/CocoSgt/riemann-hypothesis-lean/blob/7cbf7a0c32fca1badf05931879a37c3f8a2b03b4/src/Formalism/SharpRestriction/GlobalTheoremInputs.lean#L40-L44}{LinkN\_\allowbreak at}.

\emph{Source.} \href{https://github.com/CocoSgt/riemann-hypothesis-lean/blob/7cbf7a0c32fca1badf05931879a37c3f8a2b03b4/src/Development/AutomorphicArithmetic/CharacterBoxes/DivinNOfHypV.lean#L117-L134}{divinN\_\allowbreak of\_\allowbreak hypV}.



\paragraph{2. GLOBAL\_\allowbreak RETAINED\_\allowbreak DEFECT — C — Original retained-row defect bound LinkIRR}

Rows.DelpNonneg and delp$\le$exp(-x\textasciicircum{}(ez/7)) uniformly for every late actual row datum. CURRENT VERSION0.47 QUALIFICATION: Finish the same-D/common-family adelic setting after ASG: normalized nonzero Bessel W and decay, measured-field block, original Siegel-Weil rescaling and measure/character ties, Bessel uniqueness/local ties, then the original AutTie/RowDictTie/LabTie. RowsFinDim is now proved universally by RFD. Supply the actual CountIdempotentsClass theorem or its three sufficient arithmetic cylinder counts, with one simultaneous family threshold; reuse the completed ARO structure. Completed new original producers: ARO1-ARO30, RFD1-RFD9, ASG1-ASG18. Any earlier narrative naming these completed producers as outstanding is historical; this qualification is the current scope authority. CURRENT VERSION0.48 QUALIFICATION: Finish the same-D/common-family adelic setting after ASG and AMD: normalized nonzero Bessel W and decay, original orbital/fiber and Siegel-Weil measure/character ties, Bessel uniqueness/local ties, and the full original transport theorem, then AutTie/RowDictTie/LabTie. The complete MExists measured-data block and unipotent covolume normalization are now supplied, and universal RowsFinDim remains available. Supply the actual CountIdempotentsClass theorem or its three sufficient arithmetic cylinder counts with one simultaneous family threshold. Completed current producers or conditional consumers: ARO1-ARO30, RFD1-RFD9, ASG1-ASG18, AMD1-AMD17. Earlier frontier descriptions are historical; this qualification is the current scope authority. Original whole-contract premises and quantifiers are retained.

\emph{Shared objects and choice order.} one FamP; one class $\mathcal A$ and its target-compatible subclass; one coefficient field per family datum; one x and i at each application; original selected rows, words and levels; thresholds chosen before the common family; AllRows A; LabelCountGoodInjective; AllJoinA; CountAllRows

\emph{New completed inputs.} ARO1-ARO30, RFD1-RFD9, ASG1-ASG18, AMD1-AMD17

\emph{Current remaining subfrontier.} Finish the same-D/common-family adelic setting after ASG and AMD: normalized nonzero Bessel W and decay, original orbital/fiber and Siegel-Weil measure/character ties, Bessel uniqueness/local ties, and the full original transport theorem, then AutTie/RowDictTie/LabTie. The complete MExists measured-data block and unipotent covolume normalization are now supplied, and universal RowsFinDim remains available. Supply the actual CountIdempotentsClass theorem or its three sufficient arithmetic cylinder counts with one simultaneous family threshold.

\emph{Source.} \href{https://github.com/CocoSgt/riemann-hypothesis-lean/blob/7cbf7a0c32fca1badf05931879a37c3f8a2b03b4/src/Formalism/SharpRestriction/GlobalTheoremInputs.lean#L55-L58}{LinkIRR}.

\emph{Source.} \href{https://github.com/CocoSgt/riemann-hypothesis-lean/blob/7cbf7a0c32fca1badf05931879a37c3f8a2b03b4/src/Development/AbelianVarietyTheory/Proofs/GlobalInterpolation/LinkIRROfJoinEx.lean#L17-L38}{linkIRR\_\allowbreak of\_\allowbreak join\_\allowbreak ex}.

\emph{Source.} \href{https://github.com/CocoSgt/riemann-hypothesis-lean/blob/7cbf7a0c32fca1badf05931879a37c3f8a2b03b4/src/Development/AbelianVarietyTheory/Proofs/GlobalInterpolation/LinkIRROfJoinEx.lean#L143-L166}{link\_\allowbreak irr}.



\paragraph{3. GLOBAL\_\allowbreak PERIOD\_\allowbreak HEIGHT — C — Full actual period and height comparison LinkDIVt}

Closed conditional TuSize -> DGamSize -> (LinkEta -> DIVinN -> LinkDIV), with every original LinkDIV\_\allowbreak at clause.

\emph{Shared objects and choice order.} one FamP; one class $\mathcal A$ and its target-compatible subclass; one coefficient field per family datum; one x and i at each application; original selected rows, words and levels; thresholds chosen before the common family; same RR target; same Bu/besS; same Nu and Bk

\emph{Source.} \href{https://github.com/CocoSgt/riemann-hypothesis-lean/blob/7cbf7a0c32fca1badf05931879a37c3f8a2b03b4/src/Formalism/SharpRestriction/GlobalTheoremInputs.lean#L77-L90}{LinkDIV\_\allowbreak at}.

\emph{Source.} \href{https://github.com/CocoSgt/riemann-hypothesis-lean/blob/7cbf7a0c32fca1badf05931879a37c3f8a2b03b4/src/Formalism/SharpRestriction/GlobalTheoremInputs.lean#L94-L96}{LinkDIVt}.

\emph{Source.} \href{https://github.com/CocoSgt/riemann-hypothesis-lean/blob/7cbf7a0c32fca1badf05931879a37c3f8a2b03b4/src/Development/AbelianVarietyTheory/Proofs/GlobalInterpolation/LinkBrest.lean#L7-L25}{GlobalTheory.Top.link\_\allowbreak brest}.

\emph{Source.} \href{https://github.com/CocoSgt/riemann-hypothesis-lean/blob/7cbf7a0c32fca1badf05931879a37c3f8a2b03b4/src/Development/AbelianVarietyTheory/Proofs/GlobalInterpolation/WireSel.lean#L11-L41}{GlobalTheory.Top.wire\_\allowbreak sel}.



\paragraph{4. GLOBAL\_\allowbreak LOCAL\_\allowbreak GAIN — C — Original global LinkG installation}

One uniform LinkG with actual local lattice lengths, served-orbit coverage and MeanKappaBound; delta(x)->0 before varying actual indices. CURRENT VERSION0.47 QUALIFICATION: Construct the original common GroupCompatibilityFamily and selected tops/Adm2F for global local gain. LGS now supplies full auxiliary admission, actual first-growth and mean inputs; these should not be reopened. Continue the distinct denominator, period/height, normalized spectral and second-layer compatible-selection producers with the original common objects and thresholds. Completed new original producers: GLG1-GLG18, LGS1-LGS13. Any earlier narrative naming these completed producers as outstanding is historical; this qualification is the current scope authority. CURRENT VERSION0.48 QUALIFICATION: Construct the actual eventual common WPt/TSel realization and its GroupCompatibilityFamily on the given Extra comparison and kernel data. GCF supplies the full single GIs2 stage, period normalization, relative determinant and division transport; selected D3 tops and Adm2F already follow from ChosenE. LGS supplies local first-growth, mean and full auxiliary admission. Continue the separate global denominator, period/height, normalized spectral and second-layer compatible-selection producers with the original common objects and thresholds. Completed current producers or conditional consumers: GLG1-GLG18, LGS1-LGS13, GCF1-GCF18. Earlier frontier descriptions are historical; this qualification is the current scope authority. Original whole-contract premises and quantifiers are retained.

\emph{Shared objects and choice order.} one FamP; one class $\mathcal A$ and its target-compatible subclass; one coefficient field per family datum; one x and i at each application; original selected rows, words and levels; thresholds chosen before the common family; actual G.Inst; actual GK; actual lenMN; same orbit representatives

\emph{New completed inputs.} GLG1-GLG18, LGS1-LGS13, GCF1-GCF18

\emph{Current remaining subfrontier.} Construct the actual eventual common WPt/TSel realization and its GroupCompatibilityFamily on the given Extra comparison and kernel data. GCF supplies the full single GIs2 stage, period normalization, relative determinant and division transport; selected D3 tops and Adm2F already follow from ChosenE. LGS supplies local first-growth, mean and full auxiliary admission. Continue the separate global denominator, period/height, normalized spectral and second-layer compatible-selection producers with the original common objects and thresholds.

\emph{Source.} \href{https://github.com/CocoSgt/riemann-hypothesis-lean/blob/7cbf7a0c32fca1badf05931879a37c3f8a2b03b4/src/Formalism/SharpRestriction/GlobalTheoremInputs.lean#L122-L137}{LinkG\_\allowbreak at}.

\emph{Source.} \href{https://github.com/CocoSgt/riemann-hypothesis-lean/blob/7cbf7a0c32fca1badf05931879a37c3f8a2b03b4/src/Formalism/SharpRestriction/GlobalTheoremInputs.lean#L138-L141}{LinkG}.

\emph{Source.} \href{https://github.com/CocoSgt/riemann-hypothesis-lean/blob/7cbf7a0c32fca1badf05931879a37c3f8a2b03b4/src/Development/AbelianVarietyTheory/Proofs/GlobalInterpolation/LinkG.lean#L14-L27}{GlobalTheory.Top.link\_\allowbreak g}.



\paragraph{5. GLOBAL\_\allowbreak SPECTRAL — C — Original normalized kernel/Gram estimates LinkS}

Closed TuSize -> LinkS with exact Gamma normalization, off-diagonal Hilbert-Schmidt bound, diagonal error and Dall/Dstar estimate, errors exp(-x\textasciicircum{}(epsilon/8)). CURRENT VERSION0.48 QUALIFICATION: Complete the original same-family spectral/Bessel and Siegel-Weil inputs with all measure, character, transport and common-object ties. The complete MExists measured block and prescribed-measure covolume normalization are supplied; these are not the full normalized spectral theorem. Completed current producers or conditional consumers: ASG1-ASG18, AMD1-AMD17, GCF1-GCF18 at their stated scopes. Earlier frontier descriptions are historical; this qualification is the current scope authority. Original whole-contract premises and quantifiers are retained.

\emph{Shared objects and choice order.} one FamP; one class $\mathcal A$ and its target-compatible subclass; one coefficient field per family datum; one x and i at each application; original selected rows, words and levels; thresholds chosen before the common family; actual KD; same S.Cu spaces; GamOf; Dstar; Dall

\emph{New completed inputs.} ASG1-ASG18, AMD1-AMD17, GCF1-GCF18 at their stated scopes

\emph{Current remaining subfrontier.} Complete the original same-family spectral/Bessel and Siegel-Weil inputs with all measure, character, transport and common-object ties. The complete MExists measured block and prescribed-measure covolume normalization are supplied; these are not the full normalized spectral theorem.

\emph{Source.} \href{https://github.com/CocoSgt/riemann-hypothesis-lean/blob/7cbf7a0c32fca1badf05931879a37c3f8a2b03b4/src/Formalism/SharpRestriction/GlobalTheoremInputs.lean#L145-L149}{LinkS\_\allowbreak at}.

\emph{Source.} \href{https://github.com/CocoSgt/riemann-hypothesis-lean/blob/7cbf7a0c32fca1badf05931879a37c3f8a2b03b4/src/Formalism/SharpRestriction/GlobalTheoremInputs.lean#L152-L157}{LinkS}.

\emph{Source.} \href{https://github.com/CocoSgt/riemann-hypothesis-lean/blob/7cbf7a0c32fca1badf05931879a37c3f8a2b03b4/src/Development/AbelianVarietyTheory/Proofs/GlobalInterpolation/LinkAlt.lean#L16-L50}{GlobalTheory.Top.link\_\allowbreak alt}.



\paragraph{6. AUT\_\allowbreak UNRAM\_\allowbreak COEF — C — Actual unramified-coefficient automorphy}

For each actual Dat with coefficient ramification one and fixed iota, choose one marked g1 supplying FinA, AutB and LemW together. CURRENT VERSION0.48 QUALIFICATION: Construct the actual AutomorphyUnramifiedCoefficients package: for every original Dat with coefficient ramification one and fixed iota, choose one marked g1 that supplies FinA, AutB and LemW together. The original InpZ record, full RigII, CNR3, finite PropGK consumer and now the entire OneChar implication are supplied. Continue RH Cases I/V and the independent additional Dirichlet-GRH CaseIII/promodularity, SeedIV, tensor-factor and rank-four-chain tasks. Completed current producers or conditional consumers: OCH1-OCH30 (the full consumer implication), KIZ original InpZ, KTR1-KTR13, RNC16, PGK1-PGK6. Earlier frontier descriptions are historical; this qualification is the current scope authority. Original whole-contract premises and quantifiers are retained.

\emph{New completed inputs.} OCH1-OCH30 (the full consumer implication), KIZ original InpZ, KTR1-KTR13, RNC16, PGK1-PGK6

\emph{Current remaining subfrontier.} Construct the actual AutomorphyUnramifiedCoefficients package: for every original Dat with coefficient ramification one and fixed iota, choose one marked g1 that supplies FinA, AutB and LemW together. The original InpZ record, full RigII, CNR3, finite PropGK consumer and now the entire OneChar implication are supplied. Continue RH Cases I/V and the independent additional Dirichlet-GRH CaseIII/promodularity, SeedIV, tensor-factor and rank-four-chain tasks.

\emph{Source.} \href{https://github.com/CocoSgt/riemann-hypothesis-lean/blob/7cbf7a0c32fca1badf05931879a37c3f8a2b03b4/src/Development/AutomorphicArithmetic/AutomorphicVarieties/AutomorphicCases/AutomorphyUnramifiedCoefficients.lean#L24-L27}{AutomorphyUnramifiedCoefficients}.

\emph{Source.} \href{https://github.com/CocoSgt/riemann-hypothesis-lean/blob/7cbf7a0c32fca1badf05931879a37c3f8a2b03b4/src/Development/AutomorphicArithmetic/AutomorphicVarieties/AutomorphicCases/AutomorphyUnramifiedCoefficients.lean#L12-L13}{FinA}.

\emph{Source.} \href{https://github.com/CocoSgt/riemann-hypothesis-lean/blob/7cbf7a0c32fca1badf05931879a37c3f8a2b03b4/src/Development/AutomorphicArithmetic/AutomorphicVarieties/AutomorphicCases/AutomorphyUnramifiedCoefficients.lean#L15-L17}{AutB}.

\emph{Source.} \href{https://github.com/CocoSgt/riemann-hypothesis-lean/blob/7cbf7a0c32fca1badf05931879a37c3f8a2b03b4/src/Development/AutomorphicArithmetic/AutomorphicVarieties/AutomorphicCases/AutomorphyUnramifiedCoefficients.lean#L19-L22}{LemW}.



\paragraph{7. TASK\_\allowbreak CASE3\_\allowbreak AUT\_\allowbreak JOIN — B — Complete the original Case-III automorphy/image join}

Install the actual Case-III ProModAlt, SeedIVAlt and shared PMExAlt with the inherited CohomologyHyp/Ferrero-Washington/HFit/SeedII facts into the original whole Case-III automorphy/image conclusion.

\emph{Source.} \href{https://github.com/CocoSgt/riemann-hypothesis-lean/blob/7cbf7a0c32fca1badf05931879a37c3f8a2b03b4/src/Development/AutomorphicArithmetic/AutomorphicVarieties/ModularityArgument/ThresholdArguments/DenseThresholdCases/PromodularityCaseThree.lean#L43-L44}{AutomorphyRankFourDense}.

\emph{Source.} \href{https://github.com/CocoSgt/riemann-hypothesis-lean/blob/7cbf7a0c32fca1badf05931879a37c3f8a2b03b4/src/Development/AutomorphicArithmetic/AutomorphicVarieties/ModularityArgument/ThresholdArguments/DenseThresholdCases/CuFinal.lean#L44-L44}{main4\_\allowbreak alt}.



\paragraph{8. TASK\_\allowbreak CASE3\_\allowbreak PROMOD\_\allowbreak JOIN — C — Reconstruct the actual Case-III chain-to-promodularity proof}

Use the original tensor-factor tower, ChainRankFour, ChainDesc4, seeded rank-four chain reduction, actual PMExAlt and SeedIVAlt to construct ProModAlt on every original Ctx3 star point at the same fixed iota, preserving all finite-coefficient, inertia and support maps.

\emph{Source.} \href{https://github.com/CocoSgt/riemann-hypothesis-lean/blob/7cbf7a0c32fca1badf05931879a37c3f8a2b03b4/src/Development/AutomorphicArithmetic/AutomorphicVarieties/ModularityArgument/ThresholdArguments/DenseThresholdCases/PromodularityCaseThreeOfLemmas.lean#L131-L153}{UpperTheory.V.Cu.promodularity\_\allowbreak case\_\allowbreak three\_\allowbreak of\_\allowbreak chain\_\allowbreak reduction}.

\emph{Source.} \href{https://github.com/CocoSgt/riemann-hypothesis-lean/blob/7cbf7a0c32fca1badf05931879a37c3f8a2b03b4/src/Development/AutomorphicArithmetic/AutomorphicVarieties/ModularityArgument/ThresholdArguments/DenseThresholdCases/PromodularityCaseThree.lean#L40-L42}{PromodularityCaseThreeDense}.



\paragraph{9. SEED\_\allowbreak IV\_\allowbreak ALT — C — Actual Case-III seed with prechosen residual datum}

Base Ctx3/Four tower, iota and Amb -> w,s before every later Dies tower -> compatible finite-coefficient seed with exact residual support/root entries.

\emph{Source.} \href{https://github.com/CocoSgt/riemann-hypothesis-lean/blob/7cbf7a0c32fca1badf05931879a37c3f8a2b03b4/src/Development/AutomorphicArithmetic/AutomorphicVarieties/ModularityArgument/ThresholdArguments/DenseThresholdCases/PromodularityCaseThree.lean#L19-L25}{SeedIVAlt}.



\paragraph{10. TENSOR\_\allowbreak FACTOR\_\allowbreak TOWER — B — Actual tensor-factor tower and characters}

Original finite-field, tame, factor and character constraints; source-level construction is sizeable and has not been reconstructed here.

\emph{Source.} \href{https://github.com/CocoSgt/riemann-hypothesis-lean/blob/7cbf7a0c32fca1badf05931879a37c3f8a2b03b4/src/Development/AutomorphicArithmetic/AutomorphicVarieties/ModularityArgument/ThresholdArguments/DenseThresholdCases/PromodularityCaseThreeOfLemmas.lean#L131-L153}{UpperTheory.V.Cu.promodularity\_\allowbreak case\_\allowbreak three\_\allowbreak of\_\allowbreak chain\_\allowbreak reduction}.



\paragraph{11. CHAIN\_\allowbreak RANK\_\allowbreak FOUR — C — Actual SomeIn-to-AllIn propagation}

Same source Space and PM ideal, residual point Good4, original bound 30(C+1)\textasciicircum{}2 <= d; geometry must preserve all identities.

\emph{Source.} \href{https://github.com/CocoSgt/riemann-hypothesis-lean/blob/7cbf7a0c32fca1badf05931879a37c3f8a2b03b4/src/Development/AutomorphicArithmetic/AutomorphicVarieties/ModularityArgument/ThresholdArguments/DenseThresholdCases/PromodularityCaseThree.lean#L26-L32}{ChainRankFour}.



\paragraph{12. CHAIN\_\allowbreak DESC4 — B — Original residual chain descent}

Finite residual descent retaining source Chain4 compatibility.

\emph{Source.} \href{https://github.com/CocoSgt/riemann-hypothesis-lean/blob/7cbf7a0c32fca1badf05931879a37c3f8a2b03b4/src/Development/AutomorphicArithmetic/AutomorphicVarieties/ModularityArgument/ThresholdArguments/DenseThresholdCases/PromodularityCaseThreeOfLemmas.lean#L131-L153}{UpperTheory.V.Cu.promodularity\_\allowbreak case\_\allowbreak three\_\allowbreak of\_\allowbreak chain\_\allowbreak reduction}.



\paragraph{13. RANK4\_\allowbreak CHAIN\_\allowbreak REDUCTION — C — Seeded reduction to the original residual chain}

Select T1 before every later Deep tower; retain one compatible seed and Chain4 for every original residual lift over Q.

\emph{Source.} \href{https://github.com/CocoSgt/riemann-hypothesis-lean/blob/7cbf7a0c32fca1badf05931879a37c3f8a2b03b4/src/Development/AutomorphicArithmetic/AutomorphicVarieties/ModularityArgument/ThresholdArguments/DenseThresholdCases/PromodularityCaseThree.lean#L33-L39}{RankFourChainReductionDense}.

\emph{Source.} \href{https://github.com/CocoSgt/riemann-hypothesis-lean/blob/7cbf7a0c32fca1badf05931879a37c3f8a2b03b4/src/Development/AutomorphicArithmetic/AutomorphicVarieties/ModularityArgument/ThresholdArguments/DenseThresholdCases/CuPChainIV.lean#L16-L82}{UpperTheory.V.Cu.chainIV\_\allowbreak alt}.



\paragraph{14. TASK\_\allowbreak COUNT\_\allowbreak V\_\allowbreak INSTALL — B — Complete the original CountV finite assembly}

From exactly the listed source host/cohomology inputs, complete countV\_\allowbreak of\_\allowbreak finiteness\_\allowbreak duality\_\allowbreak euler\_\allowbreak host, including original host/frame formation, normalized length, the CaseV guard and R1<=2\textasciicircum{}c2.

\emph{Source.} \href{https://github.com/CocoSgt/riemann-hypothesis-lean/blob/7cbf7a0c32fca1badf05931879a37c3f8a2b03b4/src/Development/AutomorphicArithmetic/AutomorphicVarieties/HeightAndVanishing/Counts.lean#L69-L69}{CountV}.

\emph{Source.} \href{https://github.com/CocoSgt/riemann-hypothesis-lean/blob/7cbf7a0c32fca1badf05931879a37c3f8a2b03b4/src/Development/AutomorphicArithmetic/AutomorphicVarieties/ModularityArgument/ThresholdArguments/DenseThresholdCases/CountVOf.lean#L33-L60}{UpperTheory.V.UnitaryThreeInertia.countV\_\allowbreak of\_\allowbreak finiteness\_\allowbreak duality\_\allowbreak euler\_\allowbreak host}.



\paragraph{15. THM\_\allowbreak S6\_\allowbreak S — C — Original unitary-three ordinary support and automorphy}

Original host/frame and fixed iota -> one marked g1 -> every universal R,u -> weight-zero finiteness and every Type-I BPrime\_\allowbreak s point.

\emph{Source.} \href{https://github.com/CocoSgt/riemann-hypothesis-lean/blob/7cbf7a0c32fca1badf05931879a37c3f8a2b03b4/src/Development/AutomorphicArithmetic/AutomorphicVarieties/ModularityArgument/ThresholdArguments/DenseThresholdCases/ThmS6CoreS.lean#L17-L26}{ThmS6\_\allowbreak s}.

\emph{Source.} \href{https://github.com/CocoSgt/riemann-hypothesis-lean/blob/7cbf7a0c32fca1badf05931879a37c3f8a2b03b4/src/Development/AutomorphicArithmetic/AutomorphicVarieties/ModularityArgument/ThresholdArguments/DenseThresholdCases/ThmS6CoreS.lean#L12-L16}{BPrime\_\allowbreak s}.



\paragraph{16. HOST\_\allowbreak FRAME\_\allowbreak LENGTH — B — Original Case-V host/frame and normalized-length installation}

Actual HostData, frame, local-condition equality, abstract count and normalized original Selmer length must all refer to the same objects.

\emph{Source.} \href{https://github.com/CocoSgt/riemann-hypothesis-lean/blob/7cbf7a0c32fca1badf05931879a37c3f8a2b03b4/src/Development/AutomorphicArithmetic/AutomorphicVarieties/ModularityArgument/ThresholdArguments/DenseThresholdCases/CountVOf.lean#L33-L60}{UpperTheory.V.UnitaryThreeInertia.countV\_\allowbreak of\_\allowbreak finiteness\_\allowbreak duality\_\allowbreak euler\_\allowbreak host}.



\paragraph{17. HOST\_\allowbreak NUMBER\_\allowbreak S — C — Uniform original-host rank-three catalogue}

One exponent before all closures, coefficient rings, levels, characters, frames, embeddings and marked points; one separating trace function controls all later finite pairwise nonisomorphic collections.

\emph{Source.} \href{https://github.com/CocoSgt/riemann-hypothesis-lean/blob/7cbf7a0c32fca1badf05931879a37c3f8a2b03b4/src/Development/AutomorphicArithmetic/AutomorphicVarieties/ModularityArgument/ThresholdArguments/DenseThresholdCases/HostNumberS.lean#L18-L28}{HostNumber\_\allowbreak s}.



\paragraph{18. TASK\_\allowbreak HOST\_\allowbreak VAN\_\allowbreak E\_\allowbreak INSTALL — B — Transport K/duality to literal HostVanE}

Use the already supplied LemmaK/local-global duality and finiteness on the actual CountV host/frame and exact original coefficient modules.

\emph{Source.} \href{https://github.com/CocoSgt/riemann-hypothesis-lean/blob/7cbf7a0c32fca1badf05931879a37c3f8a2b03b4/src/Development/AutomorphicArithmetic/AutomorphicVarieties/UnitaryThreeDeformations/HostNumber.lean#L9-L12}{HostVanE}.

\emph{Source.} \href{https://github.com/CocoSgt/riemann-hypothesis-lean/blob/7cbf7a0c32fca1badf05931879a37c3f8a2b03b4/src/Development/AutomorphicArithmetic/AutomorphicVarieties/LevelCohomology/SoHostVanE.lean#L11-L18}{UpperTheory.V.X.So.hostVanE}.



\paragraph{19. TASK\_\allowbreak HOST\_\allowbreak XPP\_\allowbreak INSTALL — B — Transport CorX to literal HostXpp}

Preserve Psi\textasciicircum{}-2 eta, the prime-to-ell quotient, actual local conditions and the original CountV host while installing BR7.

\emph{Source.} \href{https://github.com/CocoSgt/riemann-hypothesis-lean/blob/7cbf7a0c32fca1badf05931879a37c3f8a2b03b4/src/Development/AutomorphicArithmetic/AutomorphicVarieties/UnitaryThreeDeformations/HostNumber.lean#L14-L17}{HostXpp}.

\emph{Source.} \href{https://github.com/CocoSgt/riemann-hypothesis-lean/blob/7cbf7a0c32fca1badf05931879a37c3f8a2b03b4/src/Development/AutomorphicArithmetic/AutomorphicVarieties/LevelCohomology/Xpp.lean#L11-L84}{UpperTheory.V.X.xpp}.



\paragraph{20. DESIGN\_\allowbreak CHOICES\_\allowbreak COMMON — C — Construct the original common design family and bundle}

For the one target HostExW, construct the original compatible DesignChoices on Cdes=WithTw(WithSm Cprime), where Cprime=WithChN(A intersect target class), A is the one Built construction class, and cT(eps,y)=cTmin(y). Retain every original field, tuple, word, metric and common bundle duty. DSZ10 and RMP11 give size/rank properties of class members, not member existence. The original RecExP/RecAtG and common object/row route remains coupled. Construct ONE original common design family R and its recursive/geometric half with original Inputs/Pins/Holds/ExAt/InGeo and required CoreP. PKT now reconstructs full exists\_\allowbreak pkT, rowPk\_\allowbreak nonempty and exists\_\allowbreak moved with Copy/Coh on those original given inputs and a supplied fitting later stage: one catalogue, both ChT, complete LfT, actual integral lattices and all metric/measure/period constants are retained; one numerical record precedes every eta and V/S/N/G. RDR supplies one Side0 before all readers, ReadPk/pins and full ReadL with exact o.eta=q.eta. The initial actual geometry, Fits/division/central conditions, primitive rational tests and eta agreement must still be produced. Inherited CSR supplies finite columns, one ReadC before later twists, period/evaluation and scalar/error transports, RowSideDetermined and same-R Bundle/BTw/BWit/BunU consumers. RPI supplies exact RowPairInt for actual pairs and finite translates with original mu\_\allowbreak b/fd0; first literal DivDef3 at the same updated R,F,Y,Z,C,sEE follows from its retained original readings, Pins/Holds, stored nonzero IsItW chain, five real-group predicates and group tie. Remaining common fields include host/frame, MultEq nonzero-part support, that original nonzero chain/group tie, eigensystem/carrier/parts and all-embedding logLam, HD/d3b, NEb, actual rational tests and base rows. Retain every simultaneous family hypothesis and eventual threshold. Stronger second DivRows3 remains separate; universal quotient-height CuspRapidDecay and actual common member existence are not inferred from the comparison constructors. Version 0.41 supplies the finite phase-correct rational theta route: exact coefficient extension, split representatives, square-law/PhiIs, Pins from CoversAt, class-sum/point-identity transfer to OneMapOverL, actual-good-projector RatG and both-level LiftTestB. Its original higher covers, all-tuple class-sum uniqueness including phase rescaling, one PointIdentCore before O, base ClassSumsSpan and geometric LatFull remain explicit producers. MLat is available from the original LatFull input. This sharpens the rational-test work without claiming actual common geometry, Fits or eta agreement. Version0.46 proves finite generation of the original whole integral cusp lattice by actual integral-frame evaluation. It also proves the exact identity-cover, analytic unfolding, original-field coefficient descent, global clearing-scalar and both-level lattice-transport implications. Full spanning retains the supplied original IsNab connection and the five explicit principal-chart, level-system, raw-Gaussian, short-rule/dictionary and ramified-jet producers. The higher tuple covers, actual common class sums, PointIdentCore and ClassSumsSpan remain separate. CURRENT VERSION0.47 QUALIFICATION: Construct one original common geometric family and the higher CoversAt/all-tuple class sums and uniqueness with original phase rescalings, one PointIdentCore with common lambda,z, and base ClassSumsSpan. CNC now supplies the full original IsNab; PCL supplies the conditional trace-dual principal-lattice and quantitative normalization. Actual uniformization/HasPrinc and the five principal-chart, finite-level-system, Gaussian/reproduction, short-rule/simultaneous-dictionary and ramified-jet producer groups remain. Keep all original both-level, coefficient-field, Fits and eta agreements. Completed new original producers: ICL1-ICL21, CNC1-CNC16, PCL1-PCL12. Any earlier narrative naming these completed producers as outstanding is historical; this qualification is the current scope authority. CURRENT VERSION0.48 QUALIFICATION: Construct the original finite holomorphic FamilyUnif/N1c, normalized theta families and local lifting, then the common principal charts and finite-level theta systems. PUN now supplies the exact fixed-space models, all-family center lattices, A=1 exponent matching and individual primitive points with the full point-equality criterion; CNC supplies the full canonical connection and PCL the quantitative trace-dual normalization. Complete higher CoversAt and all-tuple class sums/uniqueness with the original phase rescalings, one PointIdentCore with common lambda,z, and base ClassSumsSpan. The Gaussian/reproduction, short-rule/simultaneous-dictionary and ramified-jet producers, and all original both-level, coefficient-field, Fits and eta agreements remain explicit. Completed current producers or conditional consumers: ICL1-ICL21, CNC1-CNC16, PCL1-PCL12, PUN1-PUN24, GCF1-GCF18 at their stated scopes. Earlier frontier descriptions are historical; this qualification is the current scope authority. Original whole-contract premises and quantifiers are retained.

\emph{New completed inputs.} ICL1-ICL21, CNC1-CNC16, PCL1-PCL12, PUN1-PUN24, GCF1-GCF18 at their stated scopes

\emph{Current remaining subfrontier.} Construct the original finite holomorphic FamilyUnif/N1c, normalized theta families and local lifting, then the common principal charts and finite-level theta systems. PUN now supplies the exact fixed-space models, all-family center lattices, A=1 exponent matching and individual primitive points with the full point-equality criterion; CNC supplies the full canonical connection and PCL the quantitative trace-dual normalization. Complete higher CoversAt and all-tuple class sums/uniqueness with the original phase rescalings, one PointIdentCore with common lambda,z, and base ClassSumsSpan. The Gaussian/reproduction, short-rule/simultaneous-dictionary and ramified-jet producers, and all original both-level, coefficient-field, Fits and eta agreements remain explicit.



\paragraph{21. LOCAL\_\allowbreak NBUILT\_\allowbreak COMMON — C — Actual local specialization on the common selected family}

Construct, with the required closed-class property on the same Des(C,3), one eventual family of actual NBuilt certificates for the original Fm.D(x,i), retaining the certificate before forgetting to TSel. HLP now reconstructs full HeckeLift0, all-index commutation, HeckeStd0 and HeckeIntTw from the original MtInput and complex Hecke facts, in the full ambient GenAt range k<=nV. Every maximal ideal is an actual supported residual character class; literal MtSub is the O-span of projections of the original MtLat generators, is finite free, and localizes onto the exact NCov.part. Literal Psi has its exact range and off-space projection error, both original span assertions hold, and ActEx/ActHK are constructed on the same lattice under original semisimplicity. Original MjLat relation quotients, universal coefficient-map relation fields, selected tame generators, coefficient embeddings and all identity transports give the complete nonempty\_\allowbreak spec constructor on its full printed input list. A single compositum combines the four original coefficient requirements with big fixed beforehand. Whole ndSpec\_\allowbreak exists still needs actual common geometric/choice MtInput, complex HeckeStdPlain and the no-Jordan producer, original F3 bridges with valuation/GoodCoef/iq, original F4 common-overfield reading of the entire coefficient field, and exact arithmetic character/reciprocity/admissibility suppliers. Original AutSpec, Galois specialization and the remaining NBuilt fields on that same installation also remain. Do not replace the actual projected lattice or maps by abstract objects with the same dimensions.



\paragraph{22. SECOND\_\allowbreak LAYER\_\allowbreak COMPATIBLE\_\allowbreak SELECTION — C — Original second-layer selector and class duties}

For the same A, supply Duty on A and a Closed common original TSel family whose selected records are the original true-column rec7 records (or have the entire dependent tuple transported), with their real witness and original coefficient-field root data. Existence/compatibility remains arithmetic; merely abstractly isomorphic fields or independently selected columns do not suffice. CSR21–22 now construct the same-R Bundle/BTw/BWit/BunU outputs from the exact selected realization/row-side packet and NEb, with the original EvP/ParEv quantifier order. Actual common selection, the five stated realization packets and NEb remain upstream; none is replaced by an independently chosen compatible record. Version 0.38 closes the selected analytic extension and reconstructs the first updated DivDef3 consumer; it leaves the original common selection, geometric readings, rational tests, base rows, host/parts/NEb and exact retained DI premises upstream, with unchanged EvP/ParEv thresholds. Version 0.39 supplies the original given-data reader constructors and exact RowPairInt/first DivDef3 on original premises. It leaves one actual common geometric/comparison/test member, eta agreement, compatible selection, rational tests, base rows, host/parts/NEb and original EvP/ParEv thresholds upstream. Version 0.40 also supplies original comparison/RowPk/fitting-stage Moved algorithms and exact common data transport. The actual initial geometric member, fitting tests, eta agreement and compatible second-layer selector with original eventual thresholds remain upstream. Version 0.41 supplies the finite phase-correct rational theta route: exact coefficient extension, split representatives, square-law/PhiIs, Pins from CoversAt, class-sum/point-identity transfer to OneMapOverL, actual-good-projector RatG and both-level LiftTestB. Its original higher covers, all-tuple class-sum uniqueness including phase rescaling, one PointIdentCore before O, base ClassSumsSpan and geometric LatFull remain explicit producers. MLat is available from the original LatFull input. This sharpens the rational-test work without claiming actual common geometry, Fits or eta agreement.



\paragraph{23. BASE\_\allowbreak THETA\_\allowbreak READINGS — C — Actual base reading on the common bundle}

For every family in the original class, eventually on the SAME selected Bundle and for each i<m, supply the actual R,F,Read,ReadL, weak models/Weil data and phase-one pins. Identify their LiftTie.theta, mtRead, mjRead and ThDef3 basis with that Bundle.P.A3, P.Zs, P.S0 and its stored maps. Exact field, embeddings, metrics, measures, row family and weak choices are retained. This remains one part of the coupled BTH12 producer. PKT now supplies the complete original comparison packet, RowPk and fitting-stage Moved with Copy/Coh on original given geometry, Inputs/Pins/Holds/ExAt/InGeo and Fits. Its ChT/LfT/Reps, actual lattice carriers, original Haar and norm constants, period maps, representatives and all-place RrT are constructed together; its one numerical record precedes every eta/V/S/N/G. RDR gives the original level-zero side, moved-Read and ReadPk/pin constructors on those objects, with one Side0 before every reader and all twelve Read ties. Full ReadL retains the exact o.eta(j,v)=q.eta(j,v) guard for every 1<=j<=m and place v; arbitrary Over does not imply it. Inherited CSR fixes one ReadC before all TwArgs, transports both metrics and all thirteen reading maps, and gives actual period/evaluation with kappa=(c*nP)\textasciicircum{}(-1/2). RPI supplies exact row-pair integrability, so first updated DivDef3 follows from the retained original Pins/Holds, nonzero chain, five real-group predicates and group tie. Initial theta extension remains complete. Remaining work is actual initial common geometry and its original predicates, Fits/primitive rational tests, eta agreement and outer family/eventual/same-Bundle duties, reFace compatibility, MultEq support, original chain/group data, actual base rows and any stronger DivRows3 theta square. No new compatible reading, family threshold or Haar normalization is substituted. Version 0.41 supplies the finite phase-correct rational theta route: exact coefficient extension, split representatives, square-law/PhiIs, Pins from CoversAt, class-sum/point-identity transfer to OneMapOverL, actual-good-projector RatG and both-level LiftTestB. Its original higher covers, all-tuple class-sum uniqueness including phase rescaling, one PointIdentCore before O, base ClassSumsSpan and geometric LatFull remain explicit producers. MLat is available from the original LatFull input. This sharpens the rational-test work without claiming actual common geometry, Fits or eta agreement.



\paragraph{24. BASE\_\allowbreak THETA\_\allowbreak RATIONAL\_\allowbreak TESTS — C — Simultaneous original base rational tests}

On the SAME original Bundle/reading/theta map, for every i<m supply one finite separating complex-linear test family whose values lie in the fixed embedded coefficient field on theta of EVERY original rational source vector and on EVERY stored target-coordinate basis vector. BTH8/BTH11 use this same LiftTestB for rational adjoint image and matrix coordinates. RTG now reconstructs prescribed-field coefficient extension, original split class representatives/RepsOK, the self-dual split-lattice square law and finite residue PhiIs; every Pins field follows from the actual CoversAt. With the base coefficient record literally fixed, exact class-sum existence/uniqueness including phase-rescaled O, PointIdentCore with common lambda,z before every O, and base ClassSumsSpan imply all-tuple finite lift span and OneMapOverL. The phase is exactly cj(twAt embLL A3.A.v0)(R.phs i). RatG uses the actual good projector eOp(y), and full LiftTestB follows at BOTH original levels. These finite implication proofs are complete. Remaining original producers are the prescribed higher covers, those actual class sums/uniqueness, PointIdentCore, base ClassSumsSpan and the base/common geometric datum; exact MLat is supplied only from original LatFull, whose geometric production remains separate. Finite dimension or an unrelated rational theta model cannot replace these same-object identities. Original outer eventual-family/Bundle quantifiers remain. Version0.46 proves finite generation of the original whole integral cusp lattice by actual integral-frame evaluation. It also proves the exact identity-cover, analytic unfolding, original-field coefficient descent, global clearing-scalar and both-level lattice-transport implications. Full spanning retains the supplied original IsNab connection and the five explicit principal-chart, level-system, raw-Gaussian, short-rule/dictionary and ramified-jet producers. The higher tuple covers, actual common class sums, PointIdentCore and ClassSumsSpan remain separate. CURRENT VERSION0.47 QUALIFICATION: Construct one original common geometric family and the higher CoversAt/all-tuple class sums and uniqueness with original phase rescalings, one PointIdentCore with common lambda,z, and base ClassSumsSpan. CNC now supplies the full original IsNab; PCL supplies the conditional trace-dual principal-lattice and quantitative normalization. Actual uniformization/HasPrinc and the five principal-chart, finite-level-system, Gaussian/reproduction, short-rule/simultaneous-dictionary and ramified-jet producer groups remain. Keep all original both-level, coefficient-field, Fits and eta agreements. Completed new original producers: ICL1-ICL21, CNC1-CNC16, PCL1-PCL12. Any earlier narrative naming these completed producers as outstanding is historical; this qualification is the current scope authority. CURRENT VERSION0.48 QUALIFICATION: Construct the original finite holomorphic FamilyUnif/N1c, normalized theta families and local lifting, then the common principal charts and finite-level theta systems. PUN now supplies the exact fixed-space models, all-family center lattices, A=1 exponent matching and individual primitive points with the full point-equality criterion; CNC supplies the full canonical connection and PCL the quantitative trace-dual normalization. Complete higher CoversAt and all-tuple class sums/uniqueness with the original phase rescalings, one PointIdentCore with common lambda,z, and base ClassSumsSpan. The Gaussian/reproduction, short-rule/simultaneous-dictionary and ramified-jet producers, and all original both-level, coefficient-field, Fits and eta agreements remain explicit. Completed current producers or conditional consumers: ICL1-ICL21, CNC1-CNC16, PCL1-PCL12, PUN1-PUN24. Earlier frontier descriptions are historical; this qualification is the current scope authority. Original whole-contract premises and quantifiers are retained.

\emph{New completed inputs.} ICL1-ICL21, CNC1-CNC16, PCL1-PCL12, PUN1-PUN24

\emph{Current remaining subfrontier.} Construct the original finite holomorphic FamilyUnif/N1c, normalized theta families and local lifting, then the common principal charts and finite-level theta systems. PUN now supplies the exact fixed-space models, all-family center lattices, A=1 exponent matching and individual primitive points with the full point-equality criterion; CNC supplies the full canonical connection and PCL the quantitative trace-dual normalization. Complete higher CoversAt and all-tuple class sums/uniqueness with the original phase rescalings, one PointIdentCore with common lambda,z, and base ClassSumsSpan. The Gaussian/reproduction, short-rule/simultaneous-dictionary and ramified-jet producers, and all original both-level, coefficient-field, Fits and eta agreements remain explicit.



\paragraph{25. SLR\_\allowbreak BASE\_\allowbreak ROWS\_\allowbreak PRODUCER — C — Actual base row identities on that same bundle}

For the SAME selected Bundle, common original fields and readings used by the other three BTH12 inputs, supply actual RowsHold0\_\allowbreak 3 uniformly eventually for every family. Its original primitive row realization, spectrum/type and basis identifications must be retained. BTH supplies a theta consumer only; it does not construct the base row identity. No separate witness may be chosen merely to satisfy this row requirement.



\clearpage
\section{The original one-character trace and length argument}
\label{sec:one-character-v48}

This chapter reconstructs the implication whose source name is
\texttt{one\allowbreak Char}: the exact \texttt{Inp\allowbreak Z} together with
\texttt{Automorphy\allowbreak Unramified\allowbreak Coefficients} implies \texttt{One\allowbreak Char}.
The latter automorphy assertion remains an input.  It supplies one
marked element and three properties of the same deformation problem;
it is not replaced by three unrelated existence assertions.
All source references are at
\texttt{7cbf7a0c32fca1badf05931879a37c3f8a2b03b4}.
The original local period-ring, finite Galois-cohomology, class-field,
and Poitou--Tate foundations are those already recorded in BX1--BX8,
DX1--DX3 and KB8.  No general crystalline deformation-ring theorem is
added as a hypothesis below.

Write \(G=G_{\mathbf Q}\), \(H=G_K\), and
\(\epsilon=\epsilon_\ell\).  The coefficient datum is the original
\(c\), with discrete valuation ring \(\mathcal O\), fraction field
\(E\), and specified inclusion into
\(\overline{\mathbf Q}_\ell\).  The datum \(GD\) includes the
original imaginary quadratic field, \(e\ge1\), odd \(\ell\),
and all its conductor, parity, order, coprimality and residual
inertia conditions.  We preserve that predicate throughout.
Put \(\Psi=\epsilon^e\psi\), let \(\eta\) be the quadratic
character of \(K/\mathbf Q\), and set
\[
 w=(1,0,-1),\qquad \operatorname{wt}(i,j)=w_i-w_j,
 \qquad T(t)=\operatorname{diag}(t,1,t^{-1}).
 \tag{OCH.1}
\]
The disconnected group is
\(\mathcal G=\mathrm{GL}_3\rtimes\{\pm1\}\), where the
nonidentity component acts by
\(\theta(A)=J A^{-t}J\), with \(J\) the reversing matrix.
The base representation has matrix \(T(\Psi(g))\) and component
\(\eta(g)\).  Its adjoint action in that component is
\(-JX^tJ\), not ordinary conjugation alone.
All matrix and coefficient maps below are actual maps, before passage
to isomorphism classes.

\subsection{Positive local cohomology and the specified ordinary component}

% BEGIN OCH1
\paragraph{OCH1. Exactness for positive finite-coefficient representations.}
\label{claim:OCH1}
Let \(L/\mathbf Q_\ell\) and \(E/\mathbf Q_\ell\) be finite.
Consider a de Rham \(E\)-representation with a filtration whose
successive characters, on an open inertia subgroup, are positive
integral powers of \(\epsilon\).  On this category the crystalline
finite-cocycle functor has the following properties: it maps a
surjection onto the finite cocycles of the quotient; and a cocycle
in a subrepresentation is finite if its image is finite.  The second
assertion assumes positivity of both representations.  These
statements allow nilpotent coefficient algebras, viewed as finite
\(E\)-vector spaces.

\emph{Proof.}
Use BX1's actual period rings and BX5's fundamental exact sequence.
For a positive character, pass to a finite extension killing its
finite inertia part.  It is \(\epsilon^a\) times an unramified
character, \(a\ge1\).  The unramified value has absolute value one:
the character is continuous on a compact group, so its image and its
inverse image are bounded.  The eigenvalues of crystalline Frobenius
therefore have nonzero slope \(-a\).  In particular \(1-\varphi\)
is invertible on its crystalline periods.  The same holds for a
successive extension, by an upper triangular matrix calculation;
finite restriction does not introduce a fixed vector.  Positive
weights also give \(\mathrm{Fil}^0D_{\rm dR}=0\).
The fundamental exact sequence consequently identifies
\(H^1_f(L,V)\) with \(D_{\rm dR}(V)\).  De Rham periods are
exact on de Rham representations, as follows by choosing a period
basis and tensoring the short exact sequence with \(B_{\rm dR}\).
Thus the asserted surjectivity and injectivity hold on classes.
For surjectivity on cocycles, lift the desired class, compare its
image with the given cocycle, and lift the vector producing their
coboundary difference.  For injectivity no adjustment is needed.
Scalar restriction, tensor coefficient retractions, and restriction
of periods preserve the original local kernels by KTR1.  This proves
all statements for the supplied coefficient algebras, not merely for
semisimple characters.
% END OCH1

% BEGIN OCH2
\paragraph{OCH2. The only positive-character obstruction is the cyclotomic annihilator.}
\label{claim:OCH2}
Let \(B\) be a finite local \(E\)-algebra with residue \(E\),
and let \(\kappa:G_L\to B^\times\) restrict to \(\epsilon^a\)
on inertia, \(a\ge1\).  Choose \(F\) generating the quotient
by inertia.  Every continuous cocycle can be written
\[
                  z=z_f+z_s,
 \qquad (\kappa(F)-\epsilon(F))z_s(g)=0,
 \tag{OCH.2}
\]
where \(z_f\) is finite.  If \(a\ge2\), one may take
\(z_s=0\).

\emph{Proof.}
Set \(d=\kappa(F)-\epsilon(F)\) if \(a=1\), and \(d=1\)
otherwise.  The ideal generated by
\(\kappa(g)-\epsilon(g)\), for all \(g\), is \((d)\).
For \(a=1\) divide the two characters and use the dense powers of
\(F\); the ideal \((d)\) is closed because \(B\) is finite
 dimensional.  For \(a>1\), inertia contains an element whose
\(\epsilon^{a-1}\)-value differs from one; the resulting nonzero
scalar in \(E\) is a unit.
Let \(W=\operatorname{ann}_B(d)\).  Both \(W\) and \(B\)
with the \(\kappa\)-action are positive representations.  The
local Euler characteristic and Tate duality in BX8 give
\[
 \dim_E Z^1(V)-\dim_E Z^1_f(V)
       =\dim_E H^0(V^*(1))
 \tag{OCH.3}
\]
for either of them; the finite-cocycle dimension here is the
Bloch--Kato dimension formula, with the same actual coboundaries
subtracted on both sides.  The last invariant space for \(B\)
is the dual of \(B/dB\), of dimension \(\dim_E W\).
On \(W\) the action is cyclotomic, so its corresponding dimension
is also \(\dim_E W\).  OCH1 makes the map
\(Z^1(W)/Z^1_f(W)\to Z^1(B)/Z^1_f(B)\) injective;
(OCH.3) makes it surjective.  Subtract a cocycle valued in \(W\)
representing the given quotient class.  This is (OCH.2).
When \(d=1\), \(W=0\).
% END OCH2

% BEGIN OCH3
\paragraph{OCH3. Crystalline triangular lifts over every finite Artin quotient.}
\label{claim:OCH3}
For the original local datum \(D:\texttt{Dat3}\), every upper
triangular crystalline deformation of the base representation lifts
through every surjection of finite local \(E\)-algebras inducing
identity on the residue field.  The component character remains
\(\eta\); its inertia diagonal remains
\((\Psi,1,\Psi^{-1})\).

\emph{Proof.}
Write an upper triangular matrix as a diagonal torus matrix times
\[
 U(x,y,z)=
 \begin{pmatrix}1&x&z+xy/2\\0&1&y\\0&0&1\end{pmatrix}.
 \tag{OCH.4}
\]
The diagonal entries constitute the source's torus cocycle.
After division by their fixed inertia characters, their infinitesimal
changes are unramified.  Indeed each successive square-zero diagonal
variation is a finite cocycle in a weight-zero character module.
The weight-zero Bloch--Kato formula identifies its finite cohomology
with the unramified classes.  The finite local component action is
retained: apply that formula after the finite extension trivializing
it, then take invariants by averaging over the finite quotient.
The same argument respects the disconnected involution on the torus.
Induction through the nilpotent filtration therefore fixes the entire
inertia character, not only its residual value.
Lift their logarithms through the surjection: logarithm and exponential are finite polynomials on the
nilpotent ideals.  On the unramified quotient the prescribed value
at Frobenius determines the character continuously; the involution
on the diagonal is respected by taking its \(+\) and \(-\)
parts, which is permissible since \(2\) is invertible.
This lifts the actual torus, including its disconnected action.

The pair \((x,y)\) is a cocycle in the two-dimensional
superdiagonal module.  Its residual characters are \(\Psi\)
and \(\eta\Psi\); the filtration by powers of the coefficient
maximal ideal makes this module positive.  Crystallinity of the
original triangular representation says exactly that this pair is
finite.  OCH1 lifts that finite cocycle, and it can be adjusted by a
lifted coboundary to have precisely the prescribed image.
The multiplication law of (OCH.4) now prescribes the coboundary of
\(z\).  Its obstruction is in \(H^2\) of the central character
\(\eta\kappa_0/\kappa_2\), whose inertia weight is \(2e\ge2\).
Local duality makes this \(H^2\) zero: a dual invariant would have
inertia weight \(1-2e\ne0\).  Thus an actual continuous \(z\)
exists.  This assertion follows equally from continuous cochains:
a finite-dimensional cohomology class that is zero has a continuous
cochain primitive.

The resulting representation is crystalline.  Over the open subgroup
\(D.HL=\ker\eta\cap\ker\tau\), its first two columns give the
finite extension determined by \(x\).  Twisting that extension by
\(\kappa_2^{-1}\), its last column is an extension with quotient
character of weight \(e\) and kernel character of weight \(2e\).
The kernel has \(H^1=H^1_f\) and \(H^2=0\), by OCH2 and
local duality.  Consequently the last-column cocycle is finite
exactly when its \(y\)-coordinate is finite.  This holds by
construction.  Finite-index descent uses the average of a period
primitive over the finite quotient; the characteristic-zero
coefficient field allows division by its order.
Finally, the difference between the constructed \(z\) and the
specified \(z\) after reduction is a finite central cocycle.
Apply OCH1 once more to lift that difference and add it.  Every
matrix entry now reduces to the original one.  This proves the
literal \texttt{tri\allowbreak Lift}, including its all-surjections quantifier.
% END OCH3

% BEGIN OCH4
\paragraph{OCH4. Every crystalline deformation of the base is ordinary.}
\label{claim:OCH4}
Over a finite local \(E\)-algebra \(B\) with residue \(E\),
a continuous deformation of the specified base representation is
crystalline on \(D.HL\) if and only if it is a conjugate of an
upper triangular crystalline deformation by a lower unipotent matrix
congruent to one.  The crystalline deformation functor is formally
smooth on these finite Artin algebras.

\emph{Proof.}
For dual numbers write the deformation as
\((1+\varepsilon\gamma(g))r_0(g)\).  It is crystalline exactly
when \(\gamma\) is a finite adjoint cocycle: apply the period
criterion to the extension of the base representation by itself.
The negative matrix-weight parts have inertia weights \(-e\)
and \(-2e\).  Their \(H^1_f\) vanishes by the character
Bloch--Kato formula in BX5.  Hence these parts are coboundaries.
Conjugating by \(1+\varepsilon X\), with \(X\) strictly lower
triangular, kills precisely them.  The diagonal and upper entries
remain unrestricted finite cocycles.  This proves the assertion
for dual numbers.

For general \(B\), use induction on its length.  A nonzero last
power of its maximal ideal contains a line \(E\varepsilon\)
annihilated by that ideal.  Apply induction modulo this line.
OCH3 lifts the triangular representative; the three entries of its
lower conjugator lift freely.  The difference between this lift and
the supplied representation is a dual-number deformation.  Its
crystalline condition really does follow from the given two
representations: their amalgam lies in
\(B\times_{B/(\varepsilon)}B\), and its underlying representation
is a subrepresentation of their direct sum.  Crystalline
representations are closed under subrepresentations.  Identifying
this fibre product with \(B\times_E E[\varepsilon]/(\varepsilon^2)\)
turns the difference into exactly the preceding finite cocycle.
Eliminate its lower part and glue its upper part.  This completes
the induction.  Conversely conjugation preserves periods, and
ordinary representatives were required to be crystalline.
Lift an ordinary representative by OCH3 and its lower conjugator
entry by entry to prove formal smoothness.  This is the full
\texttt{ord\allowbreak Eq}/\texttt{cr\allowbreak Lift} argument, including nilpotent
coefficient directions.
% END OCH4

% BEGIN OCH5
\paragraph{OCH5. The saturated ordinary ideal: all its finite-algebra points are crystalline.}
\label{claim:OCH5}
Let \(U\) be the actual local universal ring and let \(I_Y\)
be its original saturated ordinary ideal.  Every map from
\(U\) to a finite \(\mathbf Q_\ell\)-algebra annihilating
\(I_Y\) gives a representation crystalline on \(D.HL\).

\emph{Proof.}
Here is the precise ideal, so that no unnamed crystalline component
is substituted.  Choose the original inertia element \(s_0\in D.HL\)
with infinite cyclotomic value and the original Frobenius \(F\)
of \(D.HL\).  Put \(\psi_i=\Psi^{w_i}\) and
\[
 c_i=\prod_{j\ne i}(\psi_i(s_0)-\psi_j(s_0)),\qquad
 R_i(s)=r(s)-\psi_i(s)I.
\]
The ordinary ideal \(I_{\rm ord}\) imposes the characteristic
polynomials \(\prod_i(X-\psi_i(s))\) on inertia and every ordered
product \(R_0(s)R_1(s')R_2(s'')=0\).
Set \(t_i=\operatorname{tr}(R_{i+1}(s_0)R_{i+2}(s_0)r(F))\),
where the indices are cyclic, and
\[
 d_i=c_{i+1}t_i-\epsilon(F)c_i t_{i+1}\quad(i=0,1),
 \qquad \Delta=\ell d_0d_1,
 \qquad I_Y=I_{\rm ord}:\Delta^\infty.
 \tag{OCH.5}
\]
These are the source's \texttt{cc,tr,dd,del,IY}.

First reduce to a local finite target \(B\).  In
\((U/I_Y)[1/\ell]\), multiplication by \(\Delta\) is
injective, by saturation.  If its map to \(B\) kills a power
\(\mathfrak m^a\) of the inverse image of the maximal ideal,
Artin--Rees for the submodule \(\Delta A\subset A\) gives
\(b\ge a\) such that
\(\Delta z\in\mathfrak m^b\) implies
\(z\in\mathfrak m^a\).  The quotient
\(B'=A/\mathfrak m^b\) is a finite local
\(\mathbf Q_\ell\)-algebra mapping to \(B\), and that map
kills \(\operatorname{ann}_{B'}(\Delta)\).  Its finiteness
follows from the finite residue field extension and the finite
successive quotients of the maximal-ideal filtration.  Thus it is
an actual finite thickening, even when \(\Delta\) vanishes
at the given point.

Over \(B'\), the numbers \(c_i\) are nonzero coefficient-field
scalars, hence units.  The spectral projectors
\(c_i^{-1}R_{i+1}(s_0)R_{i+2}(s_0)\) have rank one.
The ordered-product equations make their partial sums the successive
stable flag: the rightmost factor kills the last quotient, then the
middle factor the next, then the leftmost factor the first.
Multiplying the equations by the projectors proves this for each
inertia element.  This ordered inertia flag is canonical.  Its first
line is the common \(\psi_0\)-eigenline: the projector for
\(s_0\) makes any such line equal to the displayed rank-one
summand.  In the quotient the same argument singles out the
\(\psi_1\)-line, since the three values at \(s_0\) are
distinct.  The characters \(\psi_i\) extend to the
local group, so conjugation carries each of these labelled inertia
submodules to one with the identical label.  Normality of inertia
and uniqueness then preserve both steps of the flag.
In a flag basis the representation on \(D.HL\) is upper triangular, with characters
\(\kappa_i\) whose inertia restrictions are \(\psi_i\).
Its adjacent normalized entries are cocycles of
\(\kappa_0/\kappa_1\) and \(\kappa_1/\kappa_2\).
Moreover \(t_i=c_i\kappa_i(F)\), so \(d_i\) is a unit
multiple of \(\kappa_i(F)/\kappa_{i+1}(F)-\epsilon(F)\).
OCH2 says that the nonfinite part of each adjacent cocycle takes
values in the annihilator of this last element.  The map to \(B\)
kills those annihilators, since it kills the annihilator of their
product \(\Delta\).  Both adjacent cocycles in \(B\) are
therefore finite.  The last-column argument of OCH3 makes the full
rank-three representation crystalline.  Conjugation does not affect
this conclusion.  A general finite algebra is a finite product of
local ones; periods and their spanning criterion respect this
product.  This proves the stated all-target assertion.
% END OCH5

% BEGIN OCH6
\paragraph{OCH6. Crystalline Artin deformations annihilate the same saturated ideal.}
\label{claim:OCH6}
Every finite Artin deformation reducing to the prescribed base and
crystalline on \(D.HL\) annihilates \(I_Y\), including points
where \(\Delta=0\).

\emph{Proof.}
OCH4 gives an ordinary flag.  Its inertia characters are fixed:
a crystalline deformation of a fixed character has unramified
infinitesimal variation, by the weight-zero Bloch--Kato formula.
Thus the characteristic-polynomial and ordered-product equations
vanish, and the map annihilates \(I_{\rm ord}\).
To pass to the saturation, suppose \(\Delta^k a\in I_{\rm ord}\).
Let \(B\) be the finite Artin target, with maximal ideal
\(\mathfrak n\) satisfying \(\mathfrak n^h=0\).
Choose \(M=2kh+1\),
\(C=E[t]/(t^M)\), and \(T=B[t]/(t^M)\).
Evaluation at zero and reduction modulo \(\mathfrak n\) give
\(p_1:T\to B\) and \(p_2:T\to C\); together they surject onto
\(B\times_E C\).
Twist the base torus by the unramified character taking the original
arithmetic Frobenius to \(1+t\).  It is crystalline and reduces
to the base.  Substitution in (OCH.5) gives
\[
                   \Delta=t^j u,\qquad 0\le j\le2,
                       \quad u\in C^\times.
 \tag{OCH.6}
\]
Indeed each adjacent difference is either a unit at \(t=0\)
or has a nonzero linear coefficient; the positive power of
Frobenius used in the source multiplies that coefficient by a
nonzero integer in \(E\).  This also explains why the potentially
critical case \(e=1\) is included.

Amalgamate this twist with the supplied deformation over
\(B\times_E C\).  The amalgam is crystalline as a
subrepresentation of their direct sum.  Lift it to \(T\) by
OCH4, and use the universal property of the actual local ring.
It annihilates \(I_{\rm ord}\), so its images \(P=\Delta^k\)
and \(y=a\) satisfy \(Py=0\); also
\(p_2(P)=t^{jk}u^k\).  Lift the unit \(u^k\) to a unit
\(V\in T\).  Then \(P=t^{jk}V+Q\), with
\(Q\in\mathfrak nT\).  Consequently
\(t^{jk}y=-QV^{-1}y\), and iteration gives
\(t^{jkh}y=0\).  Since \(jkh<M\), the coefficient of
\(t^{jkh}\) in this equality is the constant coefficient of
\(y\).  It vanishes: \(p_1(y)=0\), the required original
image of \(a\).
For a residue field larger than \(E=\operatorname{Frac}\mathcal O\),
first take the subalgebra of elements whose residue lies in \(E\).
The original universal map lands there because its reduction is the
fixed base map.  Apply the preceding argument there and include it
back into the original target.  No common coefficient enlargement
or noncriticality assumption has been imposed.
% END OCH6

% BEGIN OCH7
\paragraph{OCH7. The literal local component has a twelve-variable weighted chart.}
\label{claim:OCH7}
The source's ideal \(I_t\), the intersection of the minimal primes
of \(\sqrt{I_Y}\) through the base, is a single prime.  Its
characteristic-zero completion at the base is a formal power-series
ring in twelve variables.  The variables can be given weights with
multiplicities
\[
 \begin{array}{c|rrrrr}j&-2&-1&0&1&2\\\hline
                   \#X_j&1&2&3&4&2\end{array}.
 \tag{OCH.7}
\]
There is an actual map
\(U\to E[[X_1,\ldots,X_{12}]]\) with kernel \(I_t\),
constant coefficient the specified base map, and every universal
matrix entry homogeneous of its original matrix weight.  It has
both the linear-coordinate surjectivity and the denominator-cleared
cotangent injectivity required by \texttt{Loc\allowbreak Chart}.

\emph{Proof.}
OCH5--OCH6 identify the completed local functor of \(U/I_Y\)
with the crystalline deformation functor, on every finite Artin
algebra.  OCH4 proves formal smoothness.  Choose a basis of its
cotangent space and lift it successively modulo powers of the maximal
ideal.  Formal smoothness lifts the resulting maps at each stage;
the inverse limit is a surjection from a power-series ring.  A
relation of least degree would prevent lifting through the quotient
by that relation and the next maximal-ideal power.  Thus the kernel
is zero.  This proves the power-series assertion directly from the
lifting property.  Completion is faithfully flat on this Noetherian
local ring, so the local ring is a domain; there is exactly one
minimal prime through the base.  Pulling it back to \(U\) gives
\(I_t\), the kernel of the displayed map.  Equivalently one can
clear a single localization denominator \(y\), with nonzero base
value, to obtain \(yI_t\subset I_Y\).

The tangent space is the space of finite adjoint cocycles at \(r_0\).
Its weight \(j\) has matrix dimension \(1,2,3,2,1\).
The framed Bloch--Kato dimension formula is
\(\dim Z^1_f(V)=\dim V+\dim D_{\rm dR}(V)/\mathrm{Fil}^0\):
subtract \(H^0\) when forming coboundaries and add the same
\(H^0\) in the finite-cohomology formula.  The extra term is
zero for \(j\le0\) and is respectively \(2,1\) for
\(j=1,2\).  This gives (OCH.7), including the case where
\(\eta\) is nontrivial unramified locally.

Universal entry deviations span the cotangent space, since the
universal ring is constructed from those entries.  Select twelve
homogeneous deviations forming an \(E\)-basis.  The formal inverse
function argument changes the coordinates to precisely their images.
Conjugation by \(T(u)\) acts on each selected coordinate by
\(u^j\).  Choose an infinite-order \(u\in\mathcal O^\times\).
The equality of formal series under this action forces every monomial
of a matrix entry to have its asserted weight: unequal powers of
\(u\) are distinct.  Thus the map is graded.
If an element has zero constant and linear coefficients, its image
is in the square of the formal maximal ideal.  Faithful flatness,
then localization, puts a multiple by an element \(y\) with
nonzero base value in \((\ker x)^2+I_t\).  Modulo this square,
\(y=x(y)+(y-x(y))\); multiplying once more replaces \(y\)
by a nonzero scalar from \(\mathcal O\).  This proves the
literal denominator-cleared injectivity.  Clearing denominators in
the chosen twelve linear coordinates proves their required
surjectivity.  All assertions concern the original integral prime
\(I_t\), not an ideal chosen only after extending coefficients.
% END OCH7

% BEGIN OCH8
\paragraph{OCH8. Conjugation stability over every original complete local target.}
\label{claim:OCH8}
The ideal \(I_t\) is stable under every near-identity conjugation
in the exact \texttt{Conj\allowbreak Stable} sense.  Together with OCH7 this
proves the full original \texttt{Loc\allowbreak Good} assertion.

\emph{Proof.}
Let \(R=U/I_t\), a domain, and form
\(R[[Y_{ij}]]\).  The matrix \(1+Y\) is invertible and
congruent to one.  Conjugating the universal representation produces
an actual universal map
\(\alpha:U\to R[[Y]]\).  The ordinary equations and the
trace expressions defining \(\Delta\) are invariant under
conjugation.  Moreover \(\Delta\ne0\) in \(R\): it is a
nonzero divisor on \(U/I_Y\), so belongs to no minimal prime
of that ring.  Hence \(\alpha\) kills \(\sqrt{I_Y}\).
Its kernel is prime because its target is a domain.  Evaluation at
\(Y=0\) is the original quotient map, by uniqueness in the
universal property; consequently \(\ker\alpha\subset I_t\).
Minimality of \(I_t\) over \(\sqrt{I_Y}\) forces equality.
Now, for any original complete local target, substitute the entries
of \(h-1\) for the \(Y_{ij}\), and the supplied map for
\(R\).  Adic convergence gives the map on the full power-series
ring.  Its composite with \(\alpha\) is the actual conjugated
representation and annihilates \(I_t\).  This proves the universal
all-target statement, including targets with nilpotents or torsion.
% END OCH8

\subsection{The actual global ring and its graded model}

% BEGIN OCH9
\paragraph{OCH9. The datum, the local universal ring, and the normalized global ring exist.}
\label{claim:OCH9}
Every original \(GD\) datum admits the exact \texttt{Dat} and
\texttt{ULoc}.  For each supplied marked \(g_1\), the original
normalized ring \(S:\texttt{Dat.\allowbreak Ring}\) exists and represents
actual continuous lifts with the prescribed multiplier, local
\(I_t\)-condition, killed inertia, and normalization.

\emph{Proof.}
The quadratic extension gives \(\eta\) with kernel \(G_K\).
Choose an embedding of the algebraic closure into \(\mathbf C\)
to specify \(c\), and the original embedding into
\(\overline{\mathbf Q}_\ell\) to specify the local inclusion
\(\iota_\ell:G_{\mathbf Q_\ell}\to G\).  The equivariant
extension of embeddings identifies its range with the chosen
 decomposition group and its inertia with the chosen inertia.
Since \(\ell\nmid\operatorname{disc}K\), \(\eta\) is
unramified locally.  Put \(\tau=\operatorname{comp}_\ell\psi\)
on local inertia.  Its finite image gives an open kernel, and
\(\Psi|_{I_\ell}=\epsilon^e\tau\).  Parity in \(GD\)
gives \(\Psi(c)=1\).  These choices satisfy the actual
\texttt{Dat3} and \texttt{Dat} fields; none changes the supplied
coefficient datum or character.

Define \(\mathrm{Kill}\) to be the union, at all places away
from \(\ell\), of the kernels of the residual representation
inside inertia.  The prime-to-\(\ell\) order of \(\psi\)
implies that the base representation itself is one on this set:
reduction is injective on its prime-to-\(\ell\) roots of unity,
and the cyclotomic character is unramified there.  Choose a finite
set of places containing \(\ell\), the conductor and the
quadratic discriminant.  The quotient by the closed normal closure
of \(\mathrm{Kill}\) is a quotient of \(G_{\mathbf Q,S}\).
The finite-cohomology hypothesis gives its \(\ell\)-finiteness
condition: for each open subgroup use its permutation module and
Shapiro to embed continuous homomorphisms to \(\mathbf F_\ell\)
in a finite cohomology group.  The local version uses the same
argument with local finite cohomology.

FP2 therefore constructs universal continuous matrix lifts.  For the
disconnected group use its faithful six-dimensional block embedding:
\(\operatorname{diag}(A,\theta A)\) on the identity component,
and the corresponding two off-diagonal blocks on the other
component.  Impose the polynomial block identities and the given
component character, so the resulting maps represent exactly
\(\mathcal G\)-valued lifts.  Quotient by the killed-inertia
entries, by the image of \(I_t\), and by the eight equations
\[
 r(c)_{00}-1,\ r(c)_{11}-1,\ r(c)_{01},\ r(c)_{10},\
 r(c)_{02},\ r(c)_{20},\ Z_{01},\ Z_{10},
 \quad Z=r(g_1)-r(cg_1^{-1}c).
 \tag{OCH.8}
\]
The marked hypotheses are \(c^2=1\), \(\eta(c)=-1\),
\(\Psi(c)=1\), \(\eta(g_1)=1\), and
\(\Psi(g_1)^2-1\in\mathcal O^\times\).
The base point annihilates these equations.  The matrix relation
\(r(c)\theta(r(c))=1\), together with the first six equations
and their residual values, forces \(r(c)=1\); the two remaining
equations are precisely the source normalization.
Quotienting a complete Noetherian local ring by these closed ideals
preserves the required ring category.  Its universal property holds
for all original complete local targets, with uniqueness of the
actual map, proving the assertion.
% END OCH9

% BEGIN OCH10
\paragraph{OCH10. The exact relative presentation and its five weight counts.}
\label{claim:OCH10}
The global box ring before (OCH.8) has a weighted presentation over
\(U/I_t\), with new variables \(z_\nu\) that are actual
matrix-entry deviations.  If \(a_j\) counts the new variables
of weight \(j\), its relation count is
\[
        b_j=a_j+\begin{cases}1&|j|\le1,\\0&|j|>1.\end{cases}
 \tag{OCH.9}
\]
All variables and relations have weights between \(-2\) and \(2\).

\emph{Proof.}
Use the finite set from OCH9 and let
\(N_q=I_q\cap\ker\bar r_0\) at each place away from
\(\ell\).  It is closed and normal in the decomposition group.
Its index in inertia divides the order of the residual image,
which divides \(2|k^\times|\), and hence is prime to \(\ell\).
The local deformation condition of killing \(N_q\) is smooth
in the framed sense used by FP8.  To see the lifting statement,
average over its finite inertia image to conjugate a lifted action
back to its specified finite action.  Invariant matrices are the
image of the same averaging projector; lifts of an invariant vector
can be averaged.  The remaining Frobenius is an intertwiner of two
such finite actions.  Its corrections lie in their invariant
endomorphisms, whose units lift across nilpotent quotients.  This
proves both the small-extension lifting and the corresponding
unramified tangent condition.

Apply FP8's literal presentation theorem and DX3's framed Euler
identity weight by weight.  The possible extra global-dual-invariant
term is zero.  Indeed the \(GD.eta\) condition supplies an element
\(h\) killed by the mod-\(\ell\) cyclotomic character and by
the relevant prime-to-\(\ell\) character component, with
\(\eta(h)=-1\).  Then \(g=ch\) has
\(\eta(g)=1\), \(\bar\Psi(g)=1\), and
\(\bar\epsilon(g)=-1\).  It acts trivially on every residual
adjoint weight and as \(-1\) on its Tate dual.  Since \(\ell\)
is odd there is no invariant.  Complex conjugation acts on matrices
by \(-JX^tJ\).  Its fixed subspace in weight \(j\) has
dimension one for \(j=-1,0,1\), and zero for \(j=-2,2\):
in weights \(\pm1\) it pairs the two entries with a minus
sign; in weight zero it fixes \(\operatorname{diag}(1,0,-1)\);
in weights \(\pm2\) it is minus the identity.  This gives
(OCH.9).  FP8 retains the original local variables and relations,
and chooses each new variable among the actual framed entries;
there is no abstract change to a merely isomorphic deformation
problem.  The six-dimensional faithful embedding has the same
entry-weight differences, so no other weight occurs.
% END OCH10

% BEGIN OCH11
\paragraph{OCH11. Construction of the actual graded model and its numerical shape.}
\label{claim:OCH11}
Every normalized ring \(S\) in OCH9 has the actual graded
presentation, base point and numerical shape required by the original
\texttt{Model} definition.  Its three extension fields
\texttt{pts}, \texttt{der}, and \texttt{lift} are completed
in OCH12--OCH14.  The numerical parameters may be chosen as
\[
 r=a_{-1}+2,\quad s=a_{-2}+1,\quad
 r'=a_1+3,\quad s'=a_2+1,\quad m=a_0+3.
 \tag{OCH.10}
\]
Its ambient smooth chart has weight dimensions
\((s,r,m,r'+1,s'+1)\); its equations have counts
\((s,r+1,m,r',s')\), in weight order \(-2,-1,0,1,2\).
In particular \(r\ge1\).

\emph{Proof.}
Take the ambient complete ring
\(R=(U/I_t)[[z_\nu]]\) before imposing the global relations.
Here is its integral weighted presentation.  Choose finitely many
actual local entry deviations generating the cotangent space of
\(U\) over \(\mathcal O\); successive approximation and
Nakayama give a continuous surjection
\[
 P=\mathcal O[[X_1,\ldots,X_d,z_\nu]]
       \longrightarrow (U/I_t)[[z_\nu]]=R.
\]
Give each \(X_i\) its matrix-entry weight and the new
\(z_\nu\) the weights in OCH10.  If \(J\) is its kernel,
then \(J\) is homogeneous in the precise formal-series sense:
every projection \(f_j\) of an \(f\in J\) onto the
monomials of weight \(j\) also belongs to \(J\).
To prove this integrally, compose the local part with OCH7's
injective map \(U/I_t\hookrightarrow E[[Y_1,\ldots,Y_{12}]]\),
and extend it by the free \(z\)-variables.  This extension is
injective coefficient by coefficient.  Each substituted entry is
homogeneous of its stated weight and has zero constant coefficient,
so substitution commutes with weight projection at every finite
total degree.  If the full image is zero, the image of every
\(f_j\) is zero.  The kernel is precisely \(J\), which
proves the assertion without dividing by any difference of integral
torus eigenvalues.  Completeness and separatedness justify passage
from finite-degree truncations to full series.

Let \(R_j\) be the image of the weight-\(j\) series in
\(P\).  The homogeneous-kernel assertion makes these pieces
independent in finite sums, and multiplication sends
\(R_iR_j\) into \(R_{i+j}\).  Thus
\(A=\bigoplus_j R_j\) is the actual subalgebra of \(R\)
of finite sums of homogeneous series.  Let
\(q:A\to S\) be the induced map.  The relation list \(Q\)
consists of OCH10's relations and (OCH.8), lifted homogeneously to
\(A\).  The eight added weights are
\(0,0,1,-1,2,-2,1,-1\).  Combining them with (OCH.7) and
(OCH.9) gives exactly (OCH.10) and the two asserted count lists.
The base map kills the zero-weight equations as well.

Homogeneous projection of an expression in the complete relation
ideal shows
\(\ker q=(Q)\) in \(A\): each relation list is finite,
and projecting its coefficient series to the complementary weight
retains the required homogeneous expression.  The chart of OCH7,
with the new free coordinates added, gives the required
\texttt{Smooth\allowbreak Pt}; its cotangent injectivity is precisely OCH7's
scalar-cleared injectivity.  The weight-zero algebra is a complete
Noetherian local \(\mathcal O\)-algebra.  For completeness,
its presentation uses the finitely many minimal weight-zero
monomials; every weight-zero monomial is a product of these,
by the componentwise well-ordering of \(\mathbf N^d\).
The same argument in each fixed weight shows that \(A\) is
generated over \(A_0\) by finitely many monomials of weights
\(-2,-1,1,2\).  All nonzero weights vanish at the base.
These are the exact \texttt{Is\allowbreak WRing} fields.

Each universal three-dimensional entry has its homogeneous
preimage in \(A\).  The extra six-dimensional entries create
no new condition: after \(r(c)=1\), inverse-transpose entries
are entries at the conjugated group element.  Truncating a
homogeneous formal series modulo the \(n\)-th maximal-ideal
power expresses its image under \(q\) as an
\(\mathcal O\)-linear combination of monomials of actual
entries of the same weight.  Thus both \texttt{ent} and the
all-\(j,n\) field \texttt{dense} hold.  The remaining three
map properties are proved in OCH12--OCH14, on this same \(A,q\).
% END OCH11

% BEGIN OCH12
\paragraph{OCH12. Finite weight zero and actual local representatives of all graded points.}
\label{claim:OCH12}
If the original \texttt{Fin\allowbreak A} holds, then
\(M2\) holds for the model of OCH11.  More generally, whenever
\(M2\) holds, every algebraic-closure point of \(A/(Q)\)
can be moved by \(T(t)\) to the restriction of an actual
strictly local point of \(S\).

\emph{Proof.}
The image of \(A_0\) in \(S\) is in the closure of the
weight-zero entry algebra, by OCH11's density.  Under
\texttt{Fin\allowbreak A} that algebra is a finite \(\mathcal O\)-module.
It is closed: it is the image of a finite power of the compact ring
\(\mathcal O\) in the Hausdorff complete local ring \(S\).
Thus the image of \(A_0\), a submodule of a finite module, is
finite.  The equality \(\ker q=(Q)\) identifies this image
with the weight-zero ring \((A/(Q))_0\); it is this quotient
ring that is finite over \(\mathcal O\), as required by
\(M2\).  The ambient \(A_0\) need not be finite over
\(\mathcal O\).

Assume now only \(M2\).  Each fixed-weight piece of \(A/(Q)\)
is finite over its weight-zero piece: choose finitely many minimal
monomials of that weight under divisibility by weight-zero monomials.
The quotient weight-zero ring is a finite local
\(\mathcal O\)-algebra.  Every \(\mathcal O\)-linear map
from it to a finite coefficient field is bounded, by choosing finitely
many module generators; hence it is continuous for its adic topology.
The same observation applies to each finite fixed-weight module and
to an \(\mathcal O\)-linear point derivation.  The adic and
uniformizer topologies agree on this finite local quotient, since
its maximal ideal is nilpotent modulo the uniformizer.
Under every algebra map to
\(\overline{\mathbf Q}_\ell\), its elements are integral and
its maximal-ideal elements have absolute value less than one,
since a power of their ideal is contained in the uniformizer ideal.
Choose the finite homogeneous generators of nonzero weights.
For a positive-weight generator of weight \(a\) and a
negative-weight generator of weight \(-b\), the balanced
monomial of powers \(b,a\) is in this maximal ideal.  Whenever
both images are nonzero this gives
\[
              b\,v(\xi x)+a\,v(\xi y)>0.
\]
Therefore the finitely many strict inequalities
\(v(\xi x)+\mathrm{wt}(x)v(t)>0\) admit a common solution;
for generators of only one sign there is only one side to bound.
Choose a rational solution, hence \(t\) in the algebraic closure.
Zero images impose no constraint.  All transformed coordinates now
have absolute value less than one.

Formal evaluation converges in a finite coefficient extension:
the finitely many generator images and a basis of the finite
weight-zero image lie in one such extension.  It agrees with
\(t\cdot\xi\) on all homogeneous series, not just polynomial
coordinates.  In a fixed weight, tails of increasing total degree
are divisible by arbitrarily high powers of the weight-zero maximal
ideal, after extraction of one of finitely many minimal monomials.
Their images tend to zero.  The resulting map annihilates every
complete relation and hence factors through the original \(S\).
Every element of its maximal ideal has image of absolute value
less than one.  This is precisely \texttt{Model.\allowbreak pts}.
% END OCH12

% BEGIN OCH13
\paragraph{OCH13. Extension of every point derivation to the actual complete ring.}
\label{claim:OCH13}
Under \(M2\), let \(f:S\to\overline{\mathbf Q}_\ell\)
be strictly local.  Every point derivation of \(A\) at
\(f\circ q\) annihilating \((Q)\) is the second component
of an actual map
\(S\to\overline{\mathbf Q}_\ell[\varepsilon]/(\varepsilon^2)\)
whose first component is exactly \(f\).

\emph{Proof.}
First the literal strict-local hypothesis already makes the supplied
\(f\) continuous.  If \(u\in S\) had \(|f(u)|>1\),
then \(\ell u^h\in\mathfrak m_S\) for every \(h\),
while \(|f(\ell u^h)|\) would eventually exceed one, a
contradiction.  Thus \(f(S)\) is bounded by one.  A finite
set of generators of \(\mathfrak m_S\) has images bounded
by one number \(q_0<1\), so
\(|f(\mathfrak m_S^h)|\le q_0^h\).  This proves continuity
without an extra hypothesis on the supplied algebra homomorphism.
The map \(a\mapsto f(q(a))+\varepsilon D(a)\) is an algebra
map on \(A\).  Its coordinates have strictly small first
components.  For a formal monomial, its second component is the
finite sum obtained by differentiating one factor.  Along a tail
these terms tend to zero: the first components decay geometrically,
and the possible multiplicity factor is an integer of
\(\ell\)-adic absolute value at most one.  The finitely many
coordinate derivatives lie in one finite extension.
Agreement on arbitrary homogeneous series follows by the fixed-weight
finite-module argument in OCH12, applied also to the derivation.
On a finite \(\mathcal O\)-module every \(\mathcal O\)-linear
map to the finite coefficient field is bounded; the Leibniz rule
makes derivatives of the high powers of the weight-zero maximal
ideal tend to zero as well.  Thus formal evaluation gives a map
on the ambient complete presentation.  It kills its local relations
and the complete ideal generated by \(Q\), because it kills
those relations before completion and multiplication and limits
preserve zero.  Its first component is \(f\): the two maps
agree on the dense polynomial coordinates and are continuous.
This proves the literal \texttt{Model.\allowbreak der}.
% END OCH13

% BEGIN OCH14
\paragraph{OCH14. The integral weight-minus-one lifting field of the model.}
\label{claim:OCH14}
Let \(a\in\mathcal O\setminus\{0\}\).  Suppose a point
derivation \(D:S\to\mathcal O/(a)\) at the base has entry
weight \(-1\), and its local restriction is the reduction of an
integral weight-\(-1\) derivation of \(U\) annihilating
\(I_t\).  Then \(D\circ q\) is the reduction of an integral
weight-\(-1\) point derivation of \(A\).

\emph{Proof.}
Work in the actual ambient ring \((U/I_t)[[z_\nu]]\) of
OCH11, before the global equations.  Use the supplied integral
derivation on its local coefficient ring.  For a new variable of
weight \(-1\), lift its prescribed value from \(\mathcal O/(a)\)
to \(\mathcal O\); put zero on each other variable.
The resulting derivation on formal series is simply its constant
and linear coefficient formula at the base, so it is defined
integrally on the complete ring.  It annihilates the old local
relations by the supplied hypothesis, and restricts to the desired
weight on \(A\).
The reduction agrees with \(D\circ q\) on every coordinate.
For a six-dimensional coordinate this follows by differentiating
\(A^{-1}\) and \(JA^{-t}J\): at the diagonal base its derivative
is a unit scalar multiple of the corresponding three-dimensional
entry derivative with the same weight.  Such coordinates are the
actual ones in FP8.  Equality on coordinates implies equality on
the full complete ring.  Indeed both derivations give local maps
to the complete square-zero ring
\(\mathcal O\oplus\mathcal O/(a)\), hence are adically
continuous, and polynomials in the coordinates are dense.
Restriction proves \texttt{Model.\allowbreak lift}.  Together with
OCH11--OCH13 this completes every field of the original model;
no global equation was required to lift integrally, only its
reduction was required to vanish.
% END OCH14

\subsection{Actual points, rigidity, and admissibility}

% BEGIN OCH15
\paragraph{OCH15. Local conditions of the actual strictly local points.}
\label{claim:OCH15}
Every strictly local point of \(S\) gives the full original
\texttt{Is\allowbreak Pt}.  Its representation is conjugate to the base on
inertia away from \(\ell\), after restriction to \(H\).
At every place above \(\ell\), its \(H\)-restriction has a
stable ordered flag with inertia characters
\(\Psi,1,\Psi^{-1}\), and is potentially crystalline.
The same local crystalline condition holds for every actual
first-order deformation through \(S\).

\emph{Proof.}
A strictly local map is adically continuous.  Choose finitely many
generators of the maximal ideal.  Their images have a common norm
bound \(q<1\), while the whole ring has norm at most one:
reduce an element modulo the maximal ideal and lift its finite
residue value by an integral coefficient.  Thus the image of its
\(n\)-th power has norm at most \(q^n\).
Evaluation of a finite formal presentation then lands in the
complete finite field generated by its finitely many coordinate
images.  This proves continuity, integrality and the strict
residual inequalities in \texttt{Is\allowbreak Pt}.  The other fields follow
from the actual quotient equations in OCH9.

Away from \(\ell\), both the point and the base factor, on
\(I\cap H\), through the same finite prime-to-\(\ell\)
character quotient \(C\).  The averaging matrix
\[
              M=\sum_{g\in C}r_0(g)^{-1}r(g)
\]
intertwines them.  It is \(|C|I\) plus a matrix all of whose
entries have norm less than one, so it is invertible.  This proves
the actual inertia conjugacy, not merely equality of traces.

For the full inertia at \(\ell\), first work in the formal
chart of OCH7 modulo every power of its variable ideal.  OCH5
makes each truncated representation crystalline, and OCH4 makes
it ordinary.  Its diagonal inertia characters equal their fixed
base characters: the quotient character has weight zero and is
crystalline, so its infinitesimal cocycles are unramified; induction
on the nilpotent length proves equality on inertia.  Hence the
characteristic-polynomial and ordered-product equations hold for
\emph{every} inertia element, including those outside \(D.HL\).
Their images vanish at every order of the formal chart, thus vanish
in that chart and belong to \(I_t\).  They therefore hold at
every point of \(U/I_t\).
Choose an inertia element with three distinct \(\Psi^{w_i}\)
values.  Its spectral projectors and the ordered-product equations
produce the flag, exactly as in OCH5.  Any decomposition element
preserves it: conjugating the chosen inertia element gives another
upper triangular matrix with the same distinct ordered diagonal,
so the intertwiner must be upper triangular, by downward solution
of its matrix entries.  OCH5 gives crystallinity on \(D.HL\).
Conjugate decomposition and inertia groups give every other place;
if the conjugating global element has component \(-1\),
\(\theta\) preserves upper triangular matrices and reverses
and inverts their diagonal, leaving \((\Psi,1,\Psi^{-1})\)
in the same order.  Thus no place or flag ordering is lost.
For an arbitrary algebra map to dual numbers with strictly local
first part, its second component is also controlled by finitely many
coordinates; continuity is not an additional hypothesis.  Extend
\(\mathcal O\) finitely to \(\mathcal O'\) containing
the first coordinate values \(a_i\), all strictly small.
The scalar extension of the complete presentation is still complete,
because \(\mathcal O'\) is finite over \(\mathcal O\).
Formal Taylor division in \(\mathcal O'[[X_i]]\) expresses
every series, modulo the square of \((X_i-a_i)_i\), as its
value at \(a\) plus its linear Taylor terms.  These convergent
identities are ring identities.  The supplied map kills that square
and acts by the prescribed embedding on \(\mathcal O'\),
so it is determined by the finitely many values
\(F(X_i)=a_i+\varepsilon b_i\).  All \(a_i,b_i\) lie in
one finite field, and the Taylor formula proves adic continuity.
Apply OCH5 to this finite algebra.  It includes the nilpotent
directions without passing to their semisimplification.
% END OCH15

% BEGIN OCH16
\paragraph{OCH16. Two-sided points and their actual torus orbits.}
\label{claim:OCH16}
At a strictly local point, being neither upper nor lower triangular
on \(H\) forces exactly one of the following source shapes:
\texttt{Type\allowbreak I}, irreducibility on \(H\); or \texttt{Type\allowbreak II},
a middle line and an irreducible outer two-dimensional block.
If two such point representations are isomorphic on \(H\),
they are conjugate by \(T(t)\) on all of \(G\).

\emph{Proof.}
Normalization gives
\[
 r(g_1)-r(cg_1^{-1}c)=\zeta\operatorname{diag}(1,0,-1).
 \tag{OCH.11}
\]
Here \(\zeta\) is congruent, in the strict norm sense, to
\(\Psi(g_1)-\Psi(g_1)^{-1}\), a unit.  Thus \(\zeta\ne0\).
For two points \(\zeta'\ne-\zeta\), because otherwise twice
that fixed unit would have norm less than one, impossible when
\(\ell\ne2\).

An invariant subspace on \(H\) is stable under (OCH.11).
Its three distinct eigenvalues force it to be a coordinate subspace:
apply the three Lagrange interpolation projectors.  A zero entry
\((i,j)\) forced by such a subspace forces also the zero entry
\((2-j,2-i)\), since
\(r(cg^{-1}c)=Jr(g)^tJ\) on \(H\).
Enumerate the six nonempty proper coordinate subsets.
The subsets \(\{0\},\{0,1\}\) force upper triangularity;
\(\{2\},\{1,2\}\) force lower triangularity.
The remaining subsets \(\{1\},\{0,2\}\) force all four
odd-weight entries to vanish.  If the outer block is reducible,
its two distinct marked eigenvalues force a coordinate line and
hence again upper or lower triangularity.  This proves the shapes.

An intertwiner between two points intertwines (OCH.11), so the
nonopposition of \(\zeta,\zeta'\) makes it diagonal,
\(\operatorname{diag}(m_0,m_1,m_2)\), with all \(m_i\ne0\).
Applying the intertwining equation also at \(cg^{-1}c\) gives
\[
 (m_im_{2-i}-m_jm_{2-j})r(g)_{ij}=0.
\]
If \(m_0m_2=m_1^2\), take \(t=m_0/m_1\).
Otherwise all four odd-weight entries vanish, and choose
\(t^2=m_0/m_2\) in the algebraic closure.  In either case
\(r'(g)_{ij}=t^{w_i-w_j}r(g)_{ij}\) on \(H\).
Both representations send \(c\) to the same disconnected element
with matrix one, which commutes with \(T(t)\) in \(\mathcal G\).
Since \(G=H\sqcup Hc\), the equality holds on all of \(G\).
% END OCH16

% BEGIN OCH17
\paragraph{OCH17. The only point having both triangular signs is the base point.}
\label{claim:OCH17}
If an original point is both upper and lower on \(H\), it is
exactly \(r_0\) on \(G\).

\emph{Proof.}
Its three diagonal entries on \(H\) are characters \(a_i\),
and \(a_i/\Psi^{w_i}\) has every value in
\(1+\mathfrak m_{\overline{\mathbf Q}_\ell}\).
By OCH15 the inertia characteristic polynomial is the base one.
Away from \(\ell\), reduction is injective on the finite
prime-to-\(\ell\) values of \(\psi\), so this strict
congruence identifies the correct root:
\(a_i=\Psi^{w_i}\) on inertia.  At \(\ell\), \(GD.gen\)
and the original inertia condition give an inertia element \(s\)
with \(\Psi(s)^2-1\) a unit.  Of the exponents \(1,0,-1\),
only \(w_i\) can be strictly congruent to the \(i\)-th root
at \(s\).  The same is true at an element \(g\) for which
\(\Psi(g)^2-1\) is a unit.  Otherwise apply it to \(gs\)
and divide the character identities at \(gs\) and \(s\).
Thus the same equality holds on all inertia above \(\ell\).

Each ratio is now everywhere unramified over \(K\).  Global
class-field theory makes it factor through the finite ideal class
group.  Its order divides \(h_K\), prime to \(\ell\) by
\(GD.l\_\allowbreak cop\), whereas a root of unity of order prime to
\(\ell\) strictly congruent to one must be one (factor
\(X^n-1\) at \(1\), whose derivative \(n\) is a unit).
All ratios are trivial.  Equality on \(H\), together with
normalization at \(c\), gives equality on \(G\).
% END OCH17

% BEGIN OCH18
\paragraph{OCH18. Type-I points remain irreducible on the full cyclotomic kernel.}
\label{claim:OCH18}
If a point is of type I, its restriction to
\(H\cap\ker\epsilon\), the full infinite cyclotomic kernel,
is irreducible.

\emph{Proof.}
We use an elementary Clifford argument, spelling out its infinite
quotient scope.  The restriction of an irreducible finite-dimensional
representation to a normal subgroup is semisimple: the sum of the
conjugates of a minimal nonzero invariant subspace is the whole
space, and the sum of simple submodules is semisimple.  If the
restriction is reducible, its commuting algebra has dimension
larger than one.  The quotient here is abelian, so conjugation
acts by commuting operators on that algebra.  Its nonzero trace-zero
part is stable.  Over the algebraic closure a commuting family has
a common eigenvector (choose a minimal nonzero common invariant
subspace and an eigenspace inside it).  This yields
\(M\rho(g)=\chi(g)\rho(g)M\), with \(\chi\) trivial on
the normal subgroup.  Its character is nontrivial, since otherwise
Schur's lemma makes the trace-zero matrix scalar and hence zero.
The kernel and image of \(M\) are invariant, so \(M\) is
invertible.  This argument needs no finite-index assumption.

For \(g\) in inertia above \(\ell\) with \(\chi(g)\ne1\),
the trace of \(\rho(g)\) vanishes.  OCH15 gives
\(\Psi(g)+1+\Psi(g)^{-1}=0\), hence \(\Psi(g)^3=1\).
Cyclotomic inertia surjects onto \(\mathbf Z_\ell^\times\):
each finite cyclotomic local extension is totally ramified, and
compactness passes from all finite levels to their inverse limit.
If \(\chi\ne1\), choose an inertia element \(s\) in a
nontrivial \(\chi\)-coset and an inertia element \(u\) with
\(\epsilon(u)=1+\ell\).  If \(\chi(u)\ne1\) apply the
preceding cube relation at \(u\); otherwise apply it at \(su\)
and \(s\) and divide.  In either case \(\Psi(u)^3=1\).
Raise this equality to the finite order of \(\psi\).  It says
\((1+\ell)^{3e\operatorname{ord}\psi}=1\), impossible
in characteristic zero with \(e\ge1\).  Thus no such
self-twist exists, and the restriction is irreducible.
% END OCH18

% BEGIN OCH19
\paragraph{OCH19. The nonpositive base adjoint has no global finite cocycles modulo coboundaries.}
\label{claim:OCH19}
For every \(j\le0\), a global Bloch--Kato cocycle for
\(\operatorname{ad}r_0\), valued in matrix weight \(j\),
is a coboundary, with the exact local conditions of \texttt{rigid0}.

\emph{Proof.}
The involution \(X\mapsto JX^tJ\) preserves each matrix
weight.  Project onto its plus and minus parts using \(1/2\).
The resulting scalar characters are \(\eta\Psi^j\) and
\(\Psi^j\), respectively.  These projections intertwine the
full disconnected action and preserve all local primitives by
KTR1.  If \(j<0\), both are a finite character times a
strictly negative cyclotomic power.  Their finite local condition
at \(\ell\) is zero by the negative-weight Bloch--Kato
formula.  Their away conditions already are inertia coboundaries.
RV7's exact \texttt{so\allowbreak Bar} therefore makes each projected global
cocycle a coboundary, on the supplied closure and coefficients.

For \(j=0\), restrict to \(H\).  The scalar characters
\(1,\eta\) become trivial, and the finite condition at
\(\ell\) is unramified by BX7.  Class-field theory makes an
everywhere unramified additive cocycle factor through the finite
class group, so it is zero.  The remaining quotient is of order
two; averaging its cocycle gives a coboundary.  Reassemble the two
matrix projections, using the same equivariant inclusions.
Finite coefficient descent and individual coefficient retractions
are legitimate by KTR1; no uniform field containing all local
primitives is assumed.  This proves the original \texttt{rigid0}
for every nonpositive integer, including weights whose matrix
space is zero.
% END OCH19

% BEGIN OCH20
\paragraph{OCH20. All normalized tangent cocycles satisfy the original global finite conditions.}
\label{claim:OCH20}
If \(F:S\to\overline{\mathbf Q}_\ell[\varepsilon]/(\varepsilon^2)\)
has strictly local first part \(f\), then
\[
       \gamma(g)=\operatorname{snd}(F(r(g)))\,
                         f(r(g))^{-1}
 \tag{OCH.12}
\]
is a continuous global Bloch--Kato cocycle for
\(\operatorname{ad}(r_f)\), with the source's normalization.
Moreover, if this cocycle is a coboundary
\(X-\operatorname{ad}(r_f)(g)X\), then
\(X\in\overline{\mathbf Q}_\ell\operatorname{diag}(1,0,-1)\).

\emph{Proof.}
Differentiate the multiplication in \(\mathcal G\), including
its inverse-transpose component.  The product rule gives the
adjoint cocycle equation; (OCH.8) gives its normalization.  Its
continuity follows from OCH15's formal Taylor division for an
arbitrary supplied dual-number map; it does not presuppose continuity
of its second component.  Away from \(\ell\) it vanishes
on the killed inertia kernel.  Average on the remaining finite
inertia quotient to obtain a coboundary there.  At \(\ell\),
OCH15 makes the dual-number representation crystalline on
\(D.HL\); the extension criterion says exactly that (OCH.12)
has a crystalline period primitive.  Average over the finite
quotient and transport the original decomposition-group graph by
KTR1.  This is the full original global local condition.

For the last assertion put \(H_0=\operatorname{diag}(1,0,-1)\).
The derivative at \(c\) gives \(X=-JX^tJ\).
Differentiating the two marked zero entries in (OCH.11) gives
\([X,\zeta H_0]_{01}=[X,\zeta H_0]_{10}=0\).
Since \(2\) and \(\zeta\) are invertible, the first relation
kills \(X_{11},X_{02},X_{20}\) and pairs the remaining opposite
entries with a minus sign; the marked relations kill
\(X_{01},X_{10}\), hence also \(X_{12},X_{21}\).
Finally \(X_{22}=-X_{00}\).  Thus \(X=X_{00}H_0\).
This is an entrywise calculation over the actual coefficient ring
whenever \(2\) and \(\zeta\) are units.
% END OCH20

% BEGIN OCH21
\paragraph{OCH21. The exact point-rigidity hypotheses follow from InpZ and the marked automorphy input.}
\label{claim:OCH21}
Fix one \(D,g_1\) for which the original \texttt{Aut\allowbreak B} and
\texttt{Lem\allowbreak W} hold.  Then the original \texttt{Pt\allowbreak Hyp} holds for
every associated \(U,S\): the fixed-point assertion, every
nonpositive-weight base rigidity assertion, and full rigidity at
every two-sided point.

\emph{Proof.}
The first two assertions are OCH17 and OCH19.  At a two-sided
point use OCH16.  In type I, \texttt{Aut\allowbreak B} supplies exactly the
automorphic conjugate-self-dual representation
\(\rho_K\epsilon^{-1}\), for the given \(\iota\).
OCH18 supplies its irreducibility on the full cyclotomic kernel;
twisting does not change the restriction there.  Its pairing is
\(J\), since \(r(c)=1\) gives
\(\rho_K(cgc^{-1})=J\rho_K(g)^{-t}J\).
After the twist the pairing multiplier is \(\epsilon^{-2}\),
which is \(\epsilon^{1-3}\), and its extension to \(c\)
on the adjoint is precisely \(-JX^tJ\).
Thus all hypotheses, signs and multipliers of ADK6's literal
\texttt{Ad\allowbreak Vanish\allowbreak K} match.  Scalar twisting cancels in the
adjoint.  Every cocycle supplied by OCH20 is a coboundary.

In type II, KTR4--KTR8 prove the actual middle character is one
and construct the continuous irreducible outer extension with
its determinant and full cyclotomic condition.  The same marked
\texttt{Lem\allowbreak W} therefore applies; its universal quantifier supplies
modularity for \emph{every} original \texttt{Is\allowbreak Ext} extension.
KTR13 is exactly \texttt{Rig\allowbreak II} on the full polarized
nine-dimensional adjoint, so again each OCH20 cocycle is a
coboundary.  These are the actual predicates in \texttt{Inp\allowbreak Z},
constructed in KIZ2.  The sign convention for coboundaries causes
no change: the usual expression \(gX-X\) becomes the source's
\(X-gX\) by replacing \(X\) by \(-X\).
% END OCH21

% BEGIN OCH22
\paragraph{OCH22. Every two-sided point gives an original admissible representation.}
\label{claim:OCH22}
Under the same \texttt{Aut\allowbreak B} and \texttt{Lem\allowbreak W}, the actual
\(\rho_f:G_K\to\mathrm{GL}_3(\overline{\mathbf Q}_\ell)\)
of each two-sided point satisfies
\[
       \operatorname{Adm3}(\iota,e,N,\rho_f)
       \quad\text{or}\quad
       \operatorname{Adm2}(\iota,e,N|\operatorname{disc}K|,\rho_f).
 \tag{OCH.13}
\]

\emph{Proof.}
In type I, irreducibility is its defining condition, continuity
is OCH15, and the automorphic field is exactly \texttt{Aut\allowbreak B}.
At an away place OCH15 diagonalizes inertia with characters
\(\psi^{w_i}\), each killed by fixing the \(N\)-th roots
of unity.  At a place above \(\ell\), its ordered flag has
characters with inertia
\(\epsilon^{e},1,\epsilon^{-e}\) times those same finite
characters, and OCH15 supplies the potentially crystalline
alternative in the original local condition.  This checks every
field of \texttt{Adm3}, at every original prolongation.

In type II, KTR5--KTR7 give the actual outer representation \(W\)
of \(G\), with determinant \(\eta\), and
\(\rho_f\simeq W|_H\oplus1\) by the explicit coordinate
permutation.  \texttt{Lem\allowbreak W} gives the original newform of weight
\(2e+1\) attached to \(W\epsilon^{-e}\), with fixed \(\iota\).
Away from \(\ell\), the outer inertia characters are
\(\psi,\psi^{-1}\) when \(\eta\) is unramified.  If
\(\eta\) is ramified, coprimality of the quadratic discriminant
and \(N\) makes \(\psi\) trivial there, and the residual
quotient is of order two; its characters are \(1,\eta\).
Thus all finite characters are killed by fixing the
\(N|\operatorname{disc}K|\)-th roots, exactly the conductor
allowance in \texttt{Adm2}.
At \(\ell\), the eigenline of character \(\Psi\) in
OCH15 has zero middle coordinate, since some inertia element has
\(\Psi^2-1\) a unit.  Its two outer coordinates give a
nonzero eigenline for \(W\).  Normality of inertia makes it
stable under the decomposition group, and the determinant supplies
the second character, of inertia \(\Psi^{-1}\).
Their cyclotomic weights are \(e,-e\).  No ordered ratio is
\(\epsilon\): raise any proposed equality to the finite order
of the finite characters and evaluate an inertia element of
infinite cyclotomic order; it would force \(2e=1\) or
\(-2e=1\).  This supplies the source's nonexceptional-ratio
alternative for \(W\).  All fields of \texttt{Adm2}, including
the actual direct-sum isomorphism, now hold.
% END OCH22

\subsection{The integral Selmer-to-cotangent comparison}

% BEGIN OCH23
\paragraph{OCH23. A canonical normalized representative of each torsion Selmer class.}
\label{claim:OCH23}
Let \(\alpha=\Psi^{-1}\) act on \(E/\mathcal O\), with
its discrete topology, and let \(\mathcal S\) be the original
strict Selmer group \(\texttt{Sha\allowbreak U}\) with family
\(\texttt{sha\allowbreak Fam}(\ell N)\).  For every
\(b\in\mathcal O\setminus\{0\}\), its subgroup
\(\mathcal S[b]\) injects additively into the original
\(\texttt{sha\allowbreak Coc}(b)\).

\emph{Proof.}
Since \(\alpha(g_1)-1\) is a unit, replace a cocycle \(z\)
by
\[
 \nu[z](g)=z(g)-(\alpha(g)-1)
                     (\alpha(g_1)-1)^{-1}z(g_1).
 \tag{OCH.14}
\]
Changing \(z\) by a coboundary leaves this expression unchanged.
It is additive in the class, is zero at \(g_1\), and is zero
as a cocycle exactly when the class is zero.  For a class killed
by \(b\), all its values lie in \((E/\mathcal O)[b]\).
Use the exact isomorphism
\(\mathcal O/(b)\simeq(E/\mathcal O)[b]\),
\(\bar a\mapsto a/b+\mathcal O\), and call the resulting
scalar cocycle \(u\).  Its values are locally constant.

The decomposition group at \(\ell\) belongs to the original
\(\texttt{sha\allowbreak Fam}(\ell N)\), and every killed inertia element
belongs to one of its inertia groups.  Thus, on each such local
group, \(\nu[z](g)=(\alpha(g)-1)m\) for a local primitive.
On \(\mathrm{Kill}\), \(\Psi=1\), so the cocycle is zero.
At the distinguished local group, choose its original element
\(s\) with \(\Psi(s)^2-1\) a unit.  Then
\((\alpha(s)-1)bm=0\) forces \(bm=0\).  The local primitive
therefore already lies in \((E/\mathcal O)[b]\); it can be
represented in \(\mathcal O/(b)\).  No enlargement of the
annihilator is necessary.

Send \(u(g)\) to
\[
              \gamma(g)=u(g)(E_{10}-E_{21}).
 \tag{OCH.15}
\]
The full adjoint action of the base on this line is \(\alpha\):
\(-JX^tJ\) fixes \(E_{10}-E_{21}\), and diagonal
conjugation multiplies it by \(\Psi^{-1}\).
Thus (OCH.15) is a weight-\(-1\) matrix cocycle.
At \(c\), \(2u(c)=u(c^2)=0\), so \(u(c)=0\).
Together with \(u(g_1)=0\), the cocycle identity also gives
\(u(cg_1^{-1}c)=0\).  Hence all normalization equations hold.
A local scalar primitive \(m\) gives the source-sign matrix
primitive \(X=-m(E_{10}-E_{21})\), since
\(\operatorname{cob}X=X-\operatorname{ad}(r_0)X\).
This verifies every field of \texttt{sha\allowbreak Coc}; its \((1,0)\)
entry recovers \(u\), so the map is injective.  This proves the
literal \texttt{sha\_\allowbreak inj}, with the allowed choice \(a=b\).
% END OCH23

% BEGIN OCH24
\paragraph{OCH24. Matrix cocycles produce actual normalized square-zero deformations.}
\label{claim:OCH24}
For any \(\gamma\in\texttt{sha\allowbreak Coc}(a)\), there is a point
derivation \(D_\gamma:S\to\mathcal O/(a)\) at the actual
base point satisfying
\[
 D_\gamma(r(g)_{ij})=\Psi(g)^{w_j}\gamma(g)_{ij}.
 \tag{OCH.16}
\]
The map \(\gamma\mapsto D_\gamma\circ q\) is additive and
injective, annihilates every equation in \(Q\), and each of
its values lifts to an integral weight-\(-1\) point derivation
of \(A\).

\emph{Proof.}
In the actual complete ring
\(\mathcal O\oplus\mathcal O/(a)\), with square-zero second
summand, define the matrix of the lift by
\((1+\varepsilon\gamma(g))r_0(g)\), and keep component
\(\eta(g)\).  The cocycle identity is exactly the product
identity in \(\mathcal G\); the inverse is obtained by the
same square-zero formula.  Local constancy of \(\gamma\) and
adic continuity of \(\Psi\) give a continuous lift.  The
\texttt{sha\allowbreak Coc} normalization makes (OCH.8) zero, and its killed
inertia condition gives the required killed subgroup.

For the local condition choose the source's weight-\(-1\)
primitive \(X\) over \(\mathcal O/(a)\), and lift its two
entries to a weight-\(-1\) matrix \(Y\) over
\(\mathcal O\).  Conjugate the base representation over
\(\mathcal O\oplus\mathcal O\) by \(1+\varepsilon Y\).
OCH8 guarantees the exact \(I_t\)-condition on this integral
local deformation.  Its reduction has cocycle \(\gamma\)
on the local group.  Hence the global square-zero representation
satisfies every defining condition of \(S\), with an actual
local comparison map.  Universality gives a map from \(S\).
Its first component is the original base point, by uniqueness;
its second component is the desired derivation.  Multiplication
by the base diagonal gives exactly (OCH.16).

The integral local conjugation supplies the integral local derivative
required by OCH14.  Apply OCH14 to obtain the stated integral
lift on \(A\).  Because \(q(Q)=0\), the reduced derivation
annihilates every equation.  Uniqueness on actual entries in
(OCH.16), followed by continuity on the complete presentation,
proves additivity.  If \(D_\gamma\circ q=0\), the homogeneous
preimage of each entry supplied by OCH11 and the unit
\(\Psi(g)^{w_j}\) force every \(\gamma(g)_{ij}=0\).
This proves injectivity, including its original additive rather
than artificially strengthened linear scope.
% END OCH24

% BEGIN OCH25
\paragraph{OCH25. The full original lower bound, before any finiteness conclusion.}
\label{claim:OCH25}
For every actual model constructed above,
\[
             V\!\operatorname{len}(\Omega,c,e,\psi)
          \le \operatorname{cotLen}(x,A_{-1},Q.F)=M.\operatorname{len}
                   \quad\hbox{in }\mathbf N\cup\{\infty\}.
 \tag{OCH.17}
\]

\emph{Proof.}
For each nonzero \(b\), combine OCH23 and OCH24.  Their
additive injection from \(\mathcal S[b]\) to point derivations
annihilates the weight-\(-1\) equations \(Q.F\), and every
image lifts integrally with that weight.  Nonzero weights vanish
at the base.  These are precisely all hypotheses of SP1, the
proved generic \texttt{Lemma\allowbreak M'}; it gives
\(\operatorname{length}_{\mathcal O}\mathcal S[b]\le M.\operatorname{len}\).

Every class of the original discrete first cohomology group is
torsion.  A continuous cocycle on the compact Galois group has
finite image in the discrete module \(E/\mathcal O\); the
product of finitely many nonzero denominators annihilates it.
Consequently \(\mathcal S\) is the increasing union of its
\(\varpi^n\)-torsion submodules.  The length of an increasing
union is the supremum of these lengths: any finite strict chain
in the union is witnessed by finitely many elements, all in one
member of the union.  Thus the same bound holds for
\(\operatorname{length}_{\mathcal O}\mathcal S\), allowing
infinity on either side.  By RV1's exact quotient and local-family
identification, this is the original
\(V\!\operatorname{len}\), whose level is \(\ell N\).
This establishes \texttt{vlen\_\allowbreak le} without presupposing its
finiteness or the trace upper bound.
% END OCH25

\subsection{The upper bound and the original coefficient datum}

% BEGIN OCH26
\paragraph{OCH26. All five geometric hypotheses on the same model.}
\label{claim:OCH26}
Assume the original \texttt{Fin\allowbreak A} and \texttt{Pt\allowbreak Hyp} for the
same marked datum.  Let \(\sigma\in H\), and suppose that
the trace difference from the base is nonzero at every two-sided
point.  Then the five hypotheses \(M1\)--\(M5\) of SP6 hold
for the model \(A,Q\) and the original trace-difference element.

\emph{Proof.}
OCH12 gives \(M2\).  For \(j=-2,-1,0\), a linear
functional on the rational saturated cotangent quotient annihilating
the weight-\(j\) equations defines a weight-\(j\) point
derivation.  OCH13 lifts it to an actual first-order map on
\(S\).  Its tangent cocycle has that weight because each entry
has a preimage of its prescribed weight.  Base rigidity makes it
a coboundary; OCH20 makes its primitive a multiple of \(H_0\).
But \(H_0\) commutes with the base representation, including
the disconnected action, so this coboundary is zero.  The
derivation is zero on every entry and hence on all of \(A\),
by OCH11's density and the fixed-weight continuity established in
OCH13.  The finite-dimensional cotangent quotient therefore has
zero dual and is zero.  This gives exactly the three saturated
spanning assertions in \(M1\); its smoothness and \(r\ge1\)
were proved in OCH11.

A graded point annihilating both positive and negative weights can
be torus-rescaled to a local point by OCH12.  Its representation
is both upper and lower, hence the base by \texttt{Pt\allowbreak Hyp}.
Entry density then identifies its graded point with the base;
rescaling does not change it.  This is stronger than the required
\(M3\), so in particular proves that implication with its
original two one-sided alternatives.

For a two-sided graded point, rescale it to a local point and
rescale any two given derivations by the same weight factors.
OCH13 lifts both.  Full point rigidity and OCH20 put both tangent
cocycles in the one-dimensional space of coboundaries of \(H_0\).
The map from these derivations to cocycles is injective by entry
density, so the two derivations are linearly dependent.  Undoing
the common rescaling proves exactly \(M4\).

Finally choose homogeneous diagonal-entry preimages and let
\(a\in A_0\) be their sum at \(\sigma\), minus the base
trace.  Its base value is zero.  Trace is invariant under torus
conjugation; OCH12 therefore transfers the assumed nonzero trace
difference to every two-sided graded point.  This is \(M5\).
All five assertions concern the same \(A,Q,q,a\).
% END OCH26

% BEGIN OCH27
\paragraph{OCH27. A finite irredundant set of actual point orbits and its valuation bound.}
\label{claim:OCH27}
Under OCH26 there is a finite set \(T\) of strictly local,
two-sided maps \(S\to\overline{\mathbf Q}_\ell\), covering
all such point representations up to \(T(t)\)-conjugation,
with no two distinct members so conjugate.  Its model length is
finite and
\[
 \frac{M.\operatorname{len}}{e_{\mathcal O}}
       \le2\sum_{f\in T}v_\ell
        \bigl(\operatorname{tr}\rho_f(\sigma)
                    -\operatorname{tr}\rho_0(\sigma)\bigr).
 \tag{OCH.18}
\]
Here \(e_{\mathcal O}=v_{\mathcal O}(\ell)\), and
\(v_\ell(\ell)=1\).

\emph{Proof.}
Apply SP6, the full \texttt{Lemma\allowbreak M}, with normalized discrete
valuation \(e_{\mathcal O}v_\ell\).  It gives finite model
length and a finite covering set of two-sided graded points.
OCH12 moves each of these to an actual local point.  Both signs
are preserved: vanishing of all negative (respectively positive)
homogeneous coordinates is equivalent to upper (respectively lower)
triangularity of the actual entries.  The forward implication uses
the entry preimages, and the reverse uses their fixed-weight density.
Thus the points obtained are two-sided.  They cover the actual
points as well, because each actual point restricts along \(q\)
to a graded point and equality there determines all actual entries.

On this finite set choose one representative of each actual torus
orbit.  To justify its bound, rather than deleting terms from a
previous sum, restrict this new set again along \(q\).
It is still a full graded cover: lift any graded point by OCH12
and use the actual cover.  Conversely actual torus conjugacy
implies the same graded torus action on every homogeneous element,
by entry density and continuity.  The map \(f\mapsto f\circ q\)
is injective on our representatives, since equality gives equal
entry representations and hence the same orbit.  Apply SP6's
\emph{every-cover} inequality to this restricted set.  Its trace
values are the original ones and its indexing has no multiplicity.
Divide by the positive integer \(e_{\mathcal O}\) to obtain
(OCH.18).  This proves the literal irredundant \texttt{model\_\allowbreak upper}
consumer with all its point and valuation quantifiers.
% END OCH27

% BEGIN OCH28
\paragraph{OCH28. The unramified-coefficient OneChar conclusion.}
\label{claim:OCH28}
Assume \(e_{\mathcal O}=1\), the exact \(GD\) datum, the
given \(\iota\), \(\sigma\), and base trace \(t\), and
assume that every original \(\operatorname{Adm3}\) or
\(\operatorname{Adm2}\) representation has trace at
\(\sigma\) different from \(t\).  Under \texttt{Inp\allowbreak Z} and
\texttt{Automorphy\allowbreak Unramified\allowbreak Coefficients}, \(V\!\operatorname{len}\)
is finite and there is a finite set \(X\) of actual admissible
representations such that distinct members are not isomorphic and
\[
             V_n(\Omega,c,e,\psi)
                \le2\sum_{\rho\in X}
                       v_\ell(\operatorname{tr}\rho(\sigma)-t).
 \tag{OCH.19}
\]

\emph{Proof.}
Choose \(D\) by OCH9.  Invoke the original automorphy input
\emph{once}, for this \(D\) and this \(\iota\), obtaining
one marked \(g_1\) with \texttt{Fin\allowbreak A}, \texttt{Aut\allowbreak B}, and
\texttt{Lem\allowbreak W}.  Construct \(U,S\) by OCH9 and their model
by OCH10--OCH14, using the actual \texttt{Lemma\allowbreak P} in
\texttt{Inp\allowbreak Z}.  OCH21 supplies \texttt{Pt\allowbreak Hyp}; OCH22 makes
every two-sided point admissible.  The hypothesis on traces
therefore supplies OCH26's nonvanishing condition.
OCH27 gives the finite irredundant point set \(T\), and OCH25
gives the lower bound.  Its finite upper bound makes the original
Selmer length finite before taking its natural-number value.

Set \(X=\{\rho_f:f\in T\}\), where
\(\rho_f(\gamma)=r_f(\operatorname{galEmb}\gamma).\mathrm{left}\).
These are the actual source homomorphisms.  If two are isomorphic,
OCH16 gives a torus conjugacy of the two full point representations,
and the irredundancy of \(T\) gives equality of the points.
Thus the map \(T\to X\) is injective and \(X\) has the
required isomorphism-irredundancy.  Since \(\eta(\sigma)=1\),
its base trace is precisely the supplied
\(t=\Psi(\sigma)+1+\Psi(\sigma)^{-1}\).
Combine (OCH.17) and (OCH.18).  The source normalization of
\(V_n\) is length divided by \(e_{\mathcal O}\); the sum
is unchanged under the injection \(T\to X\).  This is
(OCH.19), with all three conclusions of the original predicate.
% END OCH28

% BEGIN OCH29
\paragraph{OCH29. An unramified coefficient datum with exactly the same character values.}
\label{claim:OCH29}
For every original character with \(\ell\nmid\operatorname{ord}\psi\),
there is an original unramified coefficient datum \(c'\) and a
Dirichlet character \(\psi'\) of the same modulus having
exactly the same values in the fixed
\(\overline{\mathbf Q}_\ell\).  The full \(GD\) predicate
is preserved.

\emph{Proof.}
Put \(m=\operatorname{ord}\psi\).  Choose in the fixed
algebraic closure a primitive \(m\)-th root of unity.  It lies
in the unramified extension whose residue field contains
\(\mu_m\): lift the simple roots of \(X^m-1\), using
\(\ell\nmid m\).  Its valuation ring is finite unramified
over \(\mathbf Z_\ell\), and hence is an allowed \(c'\)
with ramification index one.  All values of \(\psi\) on units
are powers of that root, and its nonunit values are zero.
Regard these same values in \(c'\); injectivity into the
algebraic closure verifies multiplicativity and every
Dirichlet-character condition.

Order, parity, and all primary components are unchanged.  For the
last statement, primary decomposition is the unique decomposition
of a finite-order element into its prime-primary and prime-to-that-
prime parts; injective coefficient maps preserve both powers and
orders and hence preserve those components.  The Galois characters
are formed by composition with the same cyclotomic restriction
map, so their kernels, including the \(GD.eta\) kernel
condition, are identical after inclusion in the fixed closure.
Consequently \(\operatorname{comp}_\ell(\psi')^2=1\),
\(\Psi'(-1)=\Psi(-1)\), and every value-dependent field of
\(GD\) is unchanged.  All other fields concern only
\(K,\ell,N,e\).  This proves the exact coefficient-transport
assertion, without asserting an embedding of the entire possibly
ramified original \(\mathcal O\) into \(\mathcal O'\).
% END OCH29

% BEGIN OCH30
\paragraph{OCH30. The full original OneChar, on every supplied coefficient datum.}
\label{claim:OCH30}
The literal original implication is proved:
\[
       \texttt{Inp\allowbreak Z}\ \wedge\
       \texttt{Automorphy\allowbreak Unramified\allowbreak Coefficients}
                    \quad\Longrightarrow\quad\texttt{One\allowbreak Char}.
 \tag{OCH.20}
\]
It retains every original closure, coefficient datum, fixed
\(\iota\), \(N\ne0\), \(e\), \(GD\),
\(\sigma\) and \(t\) quantifier.  In particular it adds
no lower bound on \(\ell\) or on \(e\) beyond the original
hypotheses.

\emph{Proof.}
Use OCH29 to pass from \((c,\psi)\) to
\((c',\psi')\) with exactly equal algebraic-closure values.
The two \(\Psi\)'s and their inverses agree at every element;
thus the specified \(t\) is unchanged.  The classes of
\(\operatorname{Adm3}(\iota,e,N)\) and
\(\operatorname{Adm2}(\iota,e,N|\operatorname{disc}K|)\)
representations do not refer to this auxiliary coefficient ring.
The trace-nonvanishing premise therefore stays exactly the same,
with the same \(\iota\) and \(\sigma\).
Apply OCH28 to \(c',\psi'\).  RV7's full coefficient
independence gives simultaneous finiteness of the two original
\(V\!\operatorname{len}\)'s and equality of their normalized
\(V_n\)'s.  Carry the same finite set of actual representations
\(X\) back without changing a matrix or an isomorphism condition.
Its inequality is (OCH.19) for the original coefficient datum.
These are exactly the conclusion fields in
\texttt{Kato\allowbreak Use.\allowbreak lean, One\allowbreak Char}; the proof never replaces its
irredundant finite set by a formal multiset or assumes finiteness
in order to obtain it.

KIZ2 supplies the \texttt{Inp\allowbreak Z} argument.  The separate assertion
\texttt{AUT\_\allowbreak UNRAM\_\allowbreak COEF} remains a premise of (OCH.20).
Accordingly this chapter proposes closure of the selected
\texttt{TASK\_\allowbreak ONE\_\allowbreak CHAR\_\allowbreak PROOF}, whose recorded contract is
this implication.  It does not assert that the separate automorphy
producer, the complete arithmetic installation, or RH has thereby
been completed.
% END OCH30


\section{The exact adelic bridge and the original measured setting block}
\label{sec:adelic-measures-v48}

\status{This chapter supplies the entire original \texttt{MExists},
including its literal \texttt{Meas\allowbreak Tie}, from its original antecedents.
It also supplies the generic \texttt{MGExists}, \texttt{MWCore},
\texttt{Inf\allowbreak Transfer}, prescribed-measure extension, and unipotent
covolume-one producers.  All infinity measures supplied by the source
are retained exactly.  The arguments below construct the measures;
they do not invoke the existence of a completed global setting.
The normalized nonzero Bessel model and the Siegel--Weil orbital and
fiber comparisons remain separate original producers.}

We use the source's conjugate-linear-first Hermitian convention.
Let \(E/F\) be a CM quadratic extension, let \(\sigma\) be its
nontrivial automorphism, and let \(J\) be a nondegenerate Hermitian
matrix on \(V=E^n\).  For an \(F\)-algebra \(R\), put
\[
 E_R=R\otimes_F E,\qquad
 G(R)=\{g\in\operatorname{GL}_n(E_R):g^*Jg=J\}.
 \tag{AMD.1}
\]
Here \(g^*=\sigma(g)^t\), and the source denotes \(G(R)\) by
\texttt{HA}\((R,\sigma,J)\).  Write \(\mathbb A_F\),
\(F_\infty\), and \(\mathbb A_{F,f}\) for its full, infinite,
and finite adelic rings.  A measure on a closed subgroup is viewed
as a measure on the ambient group by the pushforward along its
inclusion; this is exactly the source operation \(\operatorname{amb}\).

\subsection{One coordinate bridge, with the rational and compact sets}

% BEGIN AMD1
\paragraph{AMD1. Tensor splitting with the actual involution and topology.}
\label{claim:AMD1}
There are canonical topological ring isomorphisms, compatible with
\(E\) and its involution,
\[
 \mathbb A_F\otimes_F E\simeq
 (F_\infty\otimes_F E)\times(\mathbb A_{F,f}\otimes_F E),
 \qquad
 F_\infty\otimes_F E\simeq\prod_{v\mid\infty}\mathbb C.
 \tag{AMD.2}
\]
For an original \(I:\texttt{Aut\allowbreak Inst}\), the second map is the
source's \(\theta_v\) at each place, and the finite map followed
by tensor interchange is its \(\tau_F\).  The induced group map
\[
 \operatorname{split}:G(\mathbb A_F)
       \longrightarrow G(F_\infty)\times G(\mathbb A_{F,f})
 \tag{AMD.3}
\]
is a topological group isomorphism and is the literal
\texttt{Measures.\allowbreak split\allowbreak A}.

\begin{proof}
Choose \(\delta\ne0\) with \(\sigma\delta=-\delta\).
Then \(E=F\oplus F\delta\), and \(\delta^2\in F^\times\).
For every real place \(v\), \(\phi_v(\delta)\) is nonzero
and purely imaginary.  An element of the tensor algebra has the
unique expression \(a+b\delta\), with \(a,b\) in the scalar
ring.  The first map in (AMD.2) splits these two coefficients into
their infinite and finite parts.  At an infinite place it sends
\(a_v+b_v\delta\) to \(a_v+b_v\phi_v(\delta)\); the inverse
recovers \(a_v\) from the real part and \(b_v\) by division of
the imaginary part by the nonzero real number
\(\operatorname{Im}\phi_v(\delta)\).  These formulas prove
bijectivity, continuity in both directions, multiplicativity, and
compatibility with conjugation.  On a pure tensor they give
\(a\otimes x\mapsto a_v\phi_v(x)\) and
\(a\otimes x\mapsto x\otimes a_f\), precisely the two source
formulas.  The choices of \(\delta\) merely verify these canonical
maps and do not enter their definition.

Apply these maps to every matrix entry.  They preserve the equation
\(g^*Jg=J\).  Conversely, a matrix satisfying that equation has
the two-sided inverse \(J^{-1}g^*J\); thus the inverse is a
polynomial expression in its entries.  The matrix isomorphism
therefore restricts to the topological group isomorphism (AMD.3),
including for the topology that records both a matrix and its
inverse.  Its two components are exactly extension of scalars by
the two projections of \(\mathbb A_F\).  This is the definition
of \texttt{split\allowbreak A}.  This argument also covers \(n=0\).
\end{proof}
% END AMD1

% BEGIN AMD2
\paragraph{AMD2. The original bridge exists and its coordinate equations determine it.}
\label{claim:AMD2}
For every original \(I:\texttt{Aut\allowbreak Inst}\) there is a unique
continuous group isomorphism
\[
 b:I.G_\infty\times I.G_f\longrightarrow
                          G(\mathbb A_F)
 \tag{AMD.4}
\]
satisfying the original \texttt{Is\allowbreak Bridge}.  In particular the
whole original \texttt{Ad\allowbreak Bridge\allowbreak Iso} holds.  Its infinite and
finite factors \(b_\infty,b_f\) satisfy
\[
 \operatorname{split}\circ b=b_\infty\times b_f,
 \quad
 b_\infty(s)=\operatorname{split}(b(s,1))_\infty,
 \quad
 b_f(x)=\operatorname{split}(b(1,x))_f.
 \tag{AMD.5}
\]
They are exactly \texttt{to\allowbreak Inf\allowbreak I} and \texttt{to\allowbreak Fin\allowbreak I}, for every
supplied witness of \texttt{Ad\allowbreak Bridge\allowbreak Iso}.

\begin{proof}
The original real change-of-basis matrix satisfies
\(P_v^*\phi_v(J)P_v=J_{\epsilon_v}\).  Send a supplied real
factor \(s_v\) to \(P_vs_vP_v^{-1}\), and send its finite
factor by interchange of the two tensor factors.  AMD1 combines
these matrices into one matrix over \(\mathbb A_F\otimes_F E\).
The displayed identity for \(P_v\) shows that the real equations
\(s_v^*J_{\epsilon_v}s_v=J_{\epsilon_v}\) become exactly the
equations for \(G\).  The finite equation is unchanged by tensor
interchange.  This gives a group isomorphism.  By construction
\[
 P_v^{-1}\theta_v(b(s,x))P_v=s_v,
 \qquad \tau_F(b(s,x))=x,
 \tag{AMD.6}
\]
which are all the \texttt{Is\allowbreak Bridge} equations.

We check the possibly unfamiliar finite topology in the original
record.  Its identity neighborhoods have the cofinal subgroups
\[
 K(m)=\{g:g\widehat L=\widehat L,
                    (g-1)\widehat L\subseteq m\widehat L\},
                    \qquad m\ge1.
 \tag{AMD.7}
\]
A finite full lattice is sandwiched, after choosing an \(E\)-basis,
between two nonzero integral scalar multiples of \(\mathcal O_E^n\).
Its completed lattice is consequently a compact open module in
the finite adelic vector space.  The sets (AMD.7) are open for the
usual restricted-product topology: at almost all places they are
the full integral lattice stabilizer, and at the finitely many
remaining places they are ordinary matrix congruence subgroups.
Conversely an arbitrary restricted-product identity neighborhood
contains (AMD.7) for an integer with sufficiently large valuations
at its finitely many exceptional rational primes.  The source
axioms \texttt{topf\_\allowbreak open} and \texttt{topf\_\allowbreak basis} say exactly
that these same subgroups form a cofinal open identity basis for
\(I.G_f\).  The two group topologies are therefore equal.  This
is the finite-topology argument of FV4, with the original lattice
and congruence subgroups kept visible.  Thus the isomorphism above
and its inverse are continuous.

Finally, all \(\theta_v\) and \(\tau_F\) together separate
the scalar algebra by AMD1.  Equations (AMD.6) therefore determine
every entry of \(b(s,x)\).  Any other source bridge is the same
map.  Splitting its coordinates proves (AMD.5), as well as the
assertion for an arbitrarily supplied witness, rather than only
the witness chosen in this proof.
\end{proof}
% END AMD2

% BEGIN AMD3
\paragraph{AMD3. Rational points and the literal maximal-compact conditions.}
\label{claim:AMD3}
The bridge satisfies the entire original \texttt{Ad\allowbreak Bridge\allowbreak GF} and
\texttt{Ad\allowbreak Bridge\allowbreak K}; hence the full original
\texttt{Adelic\allowbreak Isomorphism} holds for every \(I\).

\begin{proof}
For \(\gamma\in G(F)\), scalar extension has real coordinates
\(P_v^{-1}\phi_v(\gamma)P_v\) and finite coordinate
\(\gamma\otimes1\).  These are exactly the two conditions in
\(I.GFset\).  Conversely those conditions specify the extension
of this same \(\gamma\), by AMD1.  Thus
\(b(I.GFset)=\texttt{rat\allowbreak HA}\), including both inclusions and
the same rational matrix.

For the infinite compact condition, let \(s^*J_\epsilon s=J_\epsilon\).
The source requires its entries between opposite signs to vanish.
This is equivalent to \(sJ_\epsilon=J_\epsilon s\).  Together
with the unitary equation it is equivalent to \(s^*s=1\): in
the reverse direction the two preserved forms imply that \(s\)
commutes with \(J_\epsilon\).  Under \(g=PsP^{-1}\), the
identity \(s^*s=1\) is precisely
\[
 g^*(P^{-1})^*P^{-1}g=(P^{-1})^*P^{-1},
 \tag{AMD.8}
\]
which uses the source's exact majorant \(\texttt{maj\allowbreak V}\).

For the finite condition, the completed lattice generated by
\(L\) is the product of its local completed lattices.  Here is
the finite-generation point needed for equality.  Choose finitely
many generators of \(L\).  Locally every vector in a completed
lattice is a combination of these same generators with integral
coefficients; choosing these coefficients at every place gives
integral adeles.  The resulting finite sum is exactly the original
adelic span.  The reverse inclusion is immediate by evaluation.
It follows that an adelic matrix and its inverse preserve this
span if and only if they preserve every local completed lattice.
Thus the source's \(I.Kset\) is exactly \(K_f^{HA}\) under
the tensor flip.  Combined with (AMD.8) this proves
\(b(KGA)=KHA\).  AMD2's uniqueness proves both assertions for
every source bridge, as quantified in \texttt{Ad\allowbreak Bridge\allowbreak GF} and
\texttt{Ad\allowbreak Bridge\allowbreak K}.  No choice of a different lattice or a
conjugate compact subgroup has entered the conclusion.
\end{proof}
% END AMD3

\subsection{The standard flag groups and their Haar products}

An original standard frame \(P=(e_1,\ldots,e_r;
f_1,\ldots,f_r;z)\) satisfies
\[
 h(e_i,e_j)=h(f_i,f_j)=0,\quad h(e_i,f_j)=\delta_{ij},
 \quad h(z,e_i)=h(z,f_i)=0,\quad h(z,z)=-T,\quad T\ne0.
 \tag{AMD.9}
\]
Put \(U=\langle e_i\rangle_E\), \(U'=\langle f_i\rangle_E\),
\(V_0=(U\oplus U')^\perp=Ez\perp W\).  These are direct
decompositions, because their indicated Gram matrices are
nondegenerate.  In particular \(2r+1\le n\).

% BEGIN AMD4
\paragraph{AMD4. The original subgroups are closed and split into their two adelic factors.}
\label{claim:AMD4}
For every frame (AMD.9), the source's groups \(N(R),H(R),S(R)\)
are closed locally compact second-countable groups for
\(R=F_v,F_\infty,\mathbb A_{F,f},\mathbb A_F\).  They are
exactly the following groups:
\[
\begin{split}
 S(R)&=\{g:ge_i-e_i\in\langle e_j:j<i\rangle_{E_R},\\
       &\hspace{37mm}gz-z\in U_R\},\\
 H(R)&=\{g:ge_i=e_i, gf_i=f_i, gz=z\},\\
 N(R)&=\{g:ge_i-e_i\in\langle e_j:j<i\rangle_{E_R},\\
       &\hspace{20mm}gv-v\in U_R\quad\hbox{for every }v\in U_R^\perp\}.
\end{split}
 \tag{AMD.10}
\]
The map (AMD.3) restricts to topological group isomorphisms
\(N(\mathbb A_F)\simeq N(F_\infty)\times N(\mathbb A_{F,f})\)
and likewise for \(H\) and \(S\).

\begin{proof}
The first two displayed conditions are the source's
\texttt{uni\allowbreak Flag(F0)} and \texttt{stab\allowbreak T(tup\allowbreak H)} with the ordered
tuple \(F0=(e_1,\ldots,e_r,z)\).  The third is its
\texttt{Nsub}; since \(z\in U_R^\perp\), it already implies
the extra \(z\)-condition in \(S\).  The spans and orthogonal
spaces in (AMD.10) are scalar extensions of fixed direct
summands of the finite-dimensional \(E\)-space.  Choose bases
once.  Membership then means that finitely many specified matrix
blocks vanish or are the identity.  In particular the universal
condition on \(v\) is equivalent to testing a finite basis of
\(U^\perp\).  These are closed coordinate conditions, and
they commute with the two scalar projections in (AMD.2).
The converse also holds: vanishing in both components is
vanishing in the product algebra.  This proves closedness and
all the asserted splitting isomorphisms.

The local scalar algebras are finite products of local fields;
the adelic ones are restricted products with compact open
integer modules.  They are locally compact and second countable:
there are countably many finite places, countable local bases,
and countably many choices of a finite exceptional set and local
basis elements.  The equation \(g^*Jg=J\), together with its
polynomial inverse from AMD1, makes \(G(R)\) a closed subgroup
of the corresponding matrix space.  The same properties therefore
hold for the closed groups in (AMD.10).  The argument also proves
these properties for the pointwise stabilizer of any finite
tuple, including a linearly dependent or degenerate tuple.
\end{proof}
% END AMD4

% BEGIN AMD5
\paragraph{AMD5. The exact semidirect product, without a definiteness restriction.}
\label{claim:AMD5}
For every standard frame, \(H\simeq U(W)\) and multiplication
induces an algebraic isomorphism of varieties and a homeomorphism
\[
                    N(R)\times H(R)\longrightarrow S(R),
                    \qquad(n,h)\longmapsto nh.
 \tag{AMD.11}
\]
As groups, \(S=N\rtimes H\).  This includes standard frames
for which the real group \(H\) is noncompact.

\begin{proof}
A unitary map fixing the nondegenerate space
\(U\oplus U'\oplus Ez\) pointwise preserves its orthogonal
complement \(W\), and its restriction to \(W\) is an arbitrary
unitary map.  This proves the assertion about \(H\), with the
section \(j:U(W)\to H\) acting identically on the fixed space.

An element \(s\in S\) preserves \(U\), its complete flag,
and \(U^\perp\).  It induces an isometry on
\(U^\perp/U\simeq V_0\), fixes the image of \(z\), and
therefore preserves \(W=z^\perp\) in that quotient.  Call its
restriction \(q(s)\in U(W)\).  This is a homomorphism, and
\(q\circ j\) is the identity.  By (AMD.10), its kernel is
exactly \(N\): fixing both \(z\) and \(W\) in the quotient
is the assertion that every \(v\in U^\perp\) is fixed
modulo \(U\).  Thus the inverse to (AMD.11) is
\[
       s\longmapsto\bigl(sj(q(s))^{-1},\,j(q(s))\bigr).
 \tag{AMD.12}
\]
In an adapted basis \(q(s)\) is a fixed matrix block.  Matrix
inversion on the unitary group is polynomial.  Both (AMD.11)
and (AMD.12) are consequently algebraic and continuous over
each of the scalar rings in AMD4.  All these statements follow
from \(T\ne0\); no sign condition on \(T\) or on \(W\)
has been used.
\end{proof}
% END AMD5

% BEGIN AMD6
\paragraph{AMD6. Unipotence, unimodularity, and the exact modulus of the stabilizer action.}
\label{claim:AMD6}
The groups \(N(R)\) and \(H(R)\) are unimodular.  For every
\(h\in H(R)\), conjugation by \(h\) preserves every Haar
measure on \(N(R)\).  These statements hold at finite places
and on both adelic factors, as well as at infinity.

\begin{proof}
Use the ordered blocks \(U,V_0,U'\), and reverse the order
of the dual basis when reading the last block as a complete
flag.  An element of \(N\) is upper block triangular, has
unipotent upper triangular action on \(U\), identity on
\(V_0\), and the inverse dual action on \(U'\).  It is
therefore unipotent.  Its Lie algebra consists precisely of
\[
 X=\begin{pmatrix}
 A&B&C\\0&0&-J_0^{-1}B^*\\0&0&-A^*
 \end{pmatrix},
 \qquad A\hbox{ strictly upper triangular},\qquad C^*=-C,
 \tag{AMD.13}
\]
where \(J_0\) is the Gram matrix on \(V_0\).  Indeed these
are exactly the flag equations and \(X^*J+JX=0\).
The finite polynomials \(\exp X\) and \(\log(1+Y)\) are
inverse.  They preserve the flag equations, and
\(X^*J+JX=0\) is equivalent to \((\exp X)^*J\exp X=J\).
For the reverse implication use
\(J^{-1}(\log g)^*J=\log(g^{-1})=-\log g\).
Thus these polynomials identify \(N\), as a variety, with
the nilpotent \(F\)-Lie algebra (AMD.13).

An inner automorphism of a unipotent group is unipotent on its
Lie algebra and has determinant one.  For \(H\simeq U(W)\),
the adjoint determinant is also one: after a splitting-field
extension this is conjugation by \(g\in\operatorname{GL}_m\)
on \(M_m\), whose determinant is
\((\det g)^m(\det g^{-1})^m=1\).  This polynomial identity
descends to \(F\).  Haar measure in a local analytic chart
changes under an automorphism by the absolute value of its
derivative determinant at the identity.  Applying this to inner
automorphisms proves local unimodularity of \(N\) and \(H\).

For the additional action of \(H\) on \(N\), write
\(V_0=Ez\perp W\) in (AMD.13).  Conjugation fixes the
\(A\)- and \(C\)-coordinates and the column of \(B\)
corresponding to \(z\).  On the remaining \(r\) rows of
\(B\) it acts by the dual of its action on \(W\).  Its
\(F\)-determinant is therefore
\[
       N_{E/F}(\det_E h)^{-r}=1.
 \tag{AMD.14}
\]
The last equality follows by taking determinants of
\(h^*J_Wh=J_W\).  This proves measure preservation at every
local place, including split places and residue characteristic
two.  Local Haar measures can be normalized on compact open
subgroups at almost every finite place; their restricted product
is a Haar measure on the finite adelic group.  An adelic element
belongs to these integral subgroups at almost every place, and
the local moduli just computed are one at every place.  The
restricted-product measures are therefore unimodular and the
action preserves them.  Every other Haar measure is a positive
scalar multiple, so the assertions hold for every normalization.
AMD4 gives the same conclusion on the full product.  In the
case \(r=0\), \(N=1\), and the calculation reads as the
empty determinant one.
\end{proof}
% END AMD6

% BEGIN AMD7
\paragraph{AMD7. The actual subgroup product is Haar, with no inserted scalar.}
\label{claim:AMD7}
If \(\mu_N\) and \(\mu_H\) are any given Haar measures on
\(N(R)\) and \(H(R)\), respectively, then
\[
       \mu_S=\operatorname{mult}_*(\mu_N\times\mu_H)
 \tag{AMD.15}
\]
is a Haar measure on the exact \(S(R)\).  As ambient measures,
\[
 \operatorname{amb}\mu_S
   =\operatorname{mult}_*
       (\operatorname{amb}\mu_N\times\operatorname{amb}\mu_H).
 \tag{AMD.16}
\]

\begin{proof}
Under (AMD.11), left multiplication by \(n_0h_0\) is the map
\((n,h)\mapsto(n_0h_0nh_0^{-1},h_0h)\).  AMD6 and left
invariance of the two given measures show, by product integration,
that it preserves \(\mu_N\times\mu_H\).  This proves left
invariance of (AMD.15).  It is nonzero and finite on compact
sets because (AMD.11) is a homeomorphism and a product of Haar
measures is a nonzero Radon measure finite on compact sets.
It is therefore Haar.  Composing the multiplication map with
the subgroup inclusions gives exactly multiplication in the
ambient group; functoriality of pushforwards proves (AMD.16).
There is no comparison with a separately chosen normalization
of \(\mu_S\), and hence no unrecorded multiplicative constant.
\end{proof}
% END AMD7

\subsection{Weak source measures and exact infinity transfer}

% BEGIN AMD8
\paragraph{AMD8. The source's weak Haar predicate is sufficient.}
\label{claim:AMD8}
Let \(K\) be a closed unimodular subgroup of a locally compact
second-countable group.  A nonzero, sigma-finite Borel measure
\(\nu\) on the ambient group which is supported on \(K\)
and invariant under left multiplication by every element of
\(K\) is the ambient pushforward of a genuine Haar measure
on \(K\).  These are exactly the four assumptions in the
original \texttt{Is\allowbreak Haar\allowbreak On\allowbreak Set}; local finiteness need not be
added to them.

\begin{proof}
A closed embedding identifies Borel subsets of \(K\) with
ambient Borel subsets contained in \(K\).  Pull back \(\nu\)
by this identification.  Support on \(K\) proves that its
pushforward is \(\nu\); sigma-finiteness, nonzero mass, and
left invariance pass to this measure \(\nu_K\).

Choose any ordinary Haar measure \(\lambda\) on \(K\).
Second countability implies sigma-compactness, so \(\lambda\)
is sigma-finite.  Unimodularity makes it right invariant and
inversion invariant.  For any Borel sets \(A,B\subseteq K\),
Tonelli's theorem applied to the nonnegative indicator function
gives
\[
\begin{split}
 \int_K\nu_K(gA\cap B)\,d\lambda(g)
        &=\lambda(A)\nu_K(B),\\
 \int_K\nu_K(A\cap g^{-1}B)\,d\lambda(g)
        &=\lambda(B)\nu_K(A).
\end{split}
 \tag{AMD.17}
\]
For the first equality, for fixed \(x\in B\) the admissible
\(g\)'s form \(xA^{-1}\).  For the second, for fixed
\(x\in A\) they form \(Bx^{-1}\).  The two integrands on
the left are equal by left invariance of \(\nu_K\).  Thus
\[
                \lambda(A)\nu_K(B)
                     =\lambda(B)\nu_K(A).
 \tag{AMD.18}
\]
All statements so far are valid as equalities of extended
nonnegative integrals; no local finiteness of \(\nu_K\)
has been used.

Sigma-finiteness and nonzero mass supply a Borel \(A\) with
\(0<\nu_K(A)<\infty\).  Choose a compact \(B_0\) with
\(0<\lambda(B_0)<\infty\).  The right side of (AMD.18) for
these sets is finite and positive.  Consequently both
\(\lambda(A)\) and \(\nu_K(B_0)\) are finite and positive:
a zero factor would give zero, while an infinite factor paired
with a positive factor would give infinity.  Applying (AMD.18)
again, now to arbitrary \(B\), proves
\[
             \nu_K(B)=
                \frac{\nu_K(A)}{\lambda(A)}\lambda(B).
 \tag{AMD.19}
\]
This is a finite positive scalar multiple of Haar measure, and
is therefore Haar.  The proof establishes precisely the missing
regularity instead of assuming it.
\end{proof}
% END AMD8

% BEGIN AMD9
\paragraph{AMD9. The full original infinity-transfer theorem.}
\label{claim:AMD9}
The literal \texttt{Inf\allowbreak Transfer} holds for every original
\(I\), every supplied bridge witness, every standard frame,
and every pair of measures satisfying its two original
\texttt{Is\allowbreak Haar\allowbreak On\allowbreak Set} predicates.  Its conclusion retains the
given ambient measures exactly.

\begin{proof}
Under \(b_\infty\) from AMD2, a rational vector is sent from
the real coordinates \(P_v^{-1}\phi_v(x)\) to its scalar
extension in \(F_\infty\otimes_F E\).  Matrix multiplication,
the Hermitian form, and the span of each initial part of the
frame commute with this map.  The original real \(\texttt{In\allowbreak S}\)
and \(\texttt{In\allowbreak N}\) conditions therefore become the first
and third conditions of (AMD.10).  In the third condition the
extra \(z\)-clause is automatic because \(z\in U^\perp\).
The original \(\texttt{Hset}\), which fixes the union of the
two frame ranges and \(\{z\}\), becomes exactly \(H(F_\infty)\).
Thus \(b_\infty\) restricts to homeomorphic group isomorphisms
on the two specified closed real subgroups.

By AMD6 these subgroups are unimodular, even if \(H(F_\infty)\)
is noncompact.  Apply AMD8 to the two supplied ambient measures
\(\nu_N,\nu_H\), and push their subgroup Haar measures forward
under the two restricted isomorphisms.  Call the results
\(\mu_{N,\infty},\mu_{H,\infty}\).  A topological group
isomorphism carries Haar measure to Haar measure.  By composition
of the two inclusions and the original bridge,
\[
 \operatorname{amb}\mu_{N,\infty}=(b_\infty)_*\nu_N,
 \qquad
 \operatorname{amb}\mu_{H,\infty}=(b_\infty)_*\nu_H.
 \tag{AMD.20}
\]
AMD2 identifies \(b_\infty\) with the source's \texttt{to\allowbreak Inf\allowbreak L}
for the supplied witness, not merely a convenient replacement.
Equations (AMD.20), together with the two Haar assertions, are
every field of the original \texttt{Inf\allowbreak Transfer}.
\end{proof}
% END AMD9

\subsection{Every field of the original measured block}

% BEGIN AMD10
\paragraph{AMD10. The full global-group measure constructor.}
\label{claim:AMD10}
The original \texttt{MGExists} holds for every \(I\), every
given sigma-finite Haar measure \(\mu_f\) on \(I.G_f\), and
every \(r\in\mathbb N\).  The three measures specified by
its conclusion are the exact pushforwards of the original
global, infinite, and finite measures.

\begin{proof}
Use the one bridge \(b\) from AMD2.  Let \(\mu_\infty\)
be the source's product of the fixed real Haar measures
\(\texttt{REAL.\allowbreak mu\allowbreak G}\), so its original full measure is
\(\mu_\infty\times\mu_f\).  Define
\[
 \mu_G=b_*(\mu_\infty\times\mu_f),\qquad
 \mu_{G,\infty}=(b_\infty)_*\mu_\infty,\qquad
 \mu_{G,f}=(b_f)_*\mu_f.
 \tag{AMD.21}
\]
These are Haar measures.  For every finite place choose any
Haar measure \(\mu_{G,v}\) on \(G(F_v)\).  For every tuple
\(X_0\) of length \(2r+2\), choose any Haar measure \(\mu_B(X_0)\)
on its pointwise stabilizer in \(G(\mathbb A_F)\).
AMD4 proves the needed closedness, local compactness, and second
countability even for a degenerate tuple; no nonsingularity
condition on \(X_0\) is imposed here.

These are the five entries of the original \texttt{Meas\allowbreak G}.
All five Haar clauses hold.  The remaining clause follows
from (AMD.5) and product pushforward:
\[
        \operatorname{split}_*\mu_G
            =\mu_{G,\infty}\times\mu_{G,f}.
 \tag{AMD.22}
\]
Thus the entire \texttt{Meas\allowbreak G.\allowbreak Is\allowbreak Text\allowbreak Norm} holds.  Equations
(AMD.21), with \(b_\infty,b_f\) identified in AMD2, are
exactly the three equalities in \texttt{MGExists}.  The freely
chosen local and tuple-stabilizer measures have no further
normalization requirement in this original contract.
\end{proof}
% END AMD10

% BEGIN AMD11
\paragraph{AMD11. The whole generic frame-measure core.}
\label{claim:AMD11}
The full original \texttt{MWCore} holds for every CM extension,
every nondegenerate Hermitian matrix, every original standard
frame, and every pair of given infinity Haar measures
\(\mu_{N,\infty},\mu_{H,\infty}\).  Both infinity measures
are retained as equalities, not merely up to scale.

\begin{proof}
Choose arbitrary Haar measures \(\mu_{N,f}\) and \(\mu_{H,f}\)
on the finite groups.  By the restricted splitting isomorphisms
in AMD4 define the global measures by
\[
\begin{split}
 \mu_N&=(\operatorname{split}_N^{-1})_*
                  (\mu_{N,\infty}\times\mu_{N,f}),\\
 \mu_H&=(\operatorname{split}_H^{-1})_*
                  (\mu_{H,\infty}\times\mu_{H,f}).
\end{split}
 \tag{AMD.23}
\]
They are Haar.  Define \(\mu_{S,f}\) by (AMD.15) from these
two finite measures, and choose arbitrary Haar measures
\(\mu_{S,v}\) on \(S(F_v)\) for every finite place.
These are exactly the eight entries of \texttt{Meas\allowbreak W}.
AMD7 proves the Haar property for \(\mu_{S,f}\); all other
Haar properties follow from the construction.

It remains to check the three actual normalization equalities.
Pushing (AMD.23) into the ambient groups and then splitting gives
\[
\begin{split}
 \operatorname{split}_*(\operatorname{amb}\mu_H)
   &=\operatorname{amb}\mu_{H,\infty}
                       \times\operatorname{amb}\mu_{H,f},\\
 \operatorname{split}_*(\operatorname{amb}\mu_N)
   &=\operatorname{amb}\mu_{N,\infty}
                       \times\operatorname{amb}\mu_{N,f}.
\end{split}
 \tag{AMD.24}
\]
The inclusion squares commute by AMD4, so these are literal
pushforward identities.  The third equality is (AMD.16) over
\(\mathbb A_{F,f}\).  Thus every one of the eleven clauses of
\texttt{Meas\allowbreak W.\allowbreak Is\allowbreak Text\allowbreak Norm} holds.  Its infinity entries were
never changed.  This proves \texttt{MWCore} under precisely
the original standard-frame hypothesis, without asserting the
unguarded auxiliary statement for arbitrary nonstandard frames.
\end{proof}
% END AMD11

% BEGIN AMD12
\paragraph{AMD12. The full original \texttt{MExists} and all same-bridge ties.}
\label{claim:AMD12}
For every original \(D\), the entire \texttt{MExists\allowbreak Core}\((D)\)
and \texttt{MExists}\((D)\) hold.  For the literal original
antecedent blocks \(A,L,P,R_f,X_u,B_u,\Pi,M_u,W_b\), one
obtains \(M_b\) and a single bridge witness for which the
constructed setting has the complete original \texttt{Meas\allowbreak Tie}.

\begin{proof}
First fix \(A,L\), a given finite Haar measure \(\mu_f\),
and the word-indexed weak infinity measures
\(\nu_N(u),\nu_H(u)\) in \texttt{MExists\allowbreak Core}.  AMD10 gives
one \(\texttt{Meas\allowbreak G}\) and one bridge \(b\) for \(A.I\).
For each word use the original coordinate frame
\[
       P_u=(e_C(A,L),f_C(A,L),z_C(A,L,u)).
 \tag{AMD.25}
\]
This is standard.  Indeed the coordinate map sends the original
Hermitian form to \(A.I.J\); hence the six pairing identities
are exactly the fields \(L.ee,L.ff,L.ef,L.zu\_\allowbreak orth,L.zu\_\allowbreak norm\).
Its scalar is \(T=(D.\mathrm{mem}(u)).A.T\), viewed in \(A.I.F\),
and it is nonzero because the original integer \(T\) is prime.
These verify every field of the source's \texttt{frame\_\allowbreak std}
for this actual frame.

Apply AMD9 with the already chosen bridge and the given weak
measures at this word.  Apply AMD11 to the resulting two
infinity Haar measures, keeping them exactly.  The resulting
\(W_u\) has the complete normalization record and satisfies
\[
 \operatorname{amb}(W_u.\mu_{N,\infty})
                   =(b_\infty)_*\nu_N(u),\qquad
 \operatorname{amb}(W_u.\mu_{H,\infty})
                   =(b_\infty)_*\nu_H(u).
 \tag{AMD.26}
\]
This is also the full original \texttt{MWExists\allowbreak Std} construction.
Choose such \(W_u\) for every word.  There is no new bridge
choice at a word, since the original bridge was fixed first.
The \(\texttt{Meas\allowbreak G}\) and this family are precisely the
four fields of \(M_b:\texttt{MBlk}\).  Equations (AMD.21)
and (AMD.26) give all three group-measure ties and both
word-indexed subgroup-measure ties in \texttt{Meas\allowbreak Tie\allowbreak Core}.
This proves \texttt{MExists\allowbreak Core} with its full quantifiers.

For \texttt{MExists}, take \(\mu_f=B_u.\mu_f\), with its
original Haar and sigma-finiteness fields, and take the original
\(M_u.\mu_N,M_u.\mu_H\).  Its two weak-Haar fields are
exactly the antecedents just used.  In the source's dependent
constructor \(\texttt{mk}\), the relevant fields are these
same measures and these same coordinate maps.  Consequently
\texttt{Meas\allowbreak Tie\allowbreak Core} becomes the literal \texttt{Meas\allowbreak Tie}
displayed in \texttt{Hida\allowbreak Setting}, without a change of variables
or a scalar adjustment.  The other supplied blocks, including
\(W_b\), are retained as given.  No gap hypothesis, no Bessel
existence assertion, and no Siegel--Weil formula is needed to
prove this measured-block constructor.
\end{proof}
% END AMD12

\subsection{Prescribing global measures while fixing the infinite factors}

% BEGIN AMD13
\paragraph{AMD13. Absorbing the sole global Haar scalar into the finite factor.}
\label{claim:AMD13}
Let \(e:A\times B\simeq C\) be a topological group
isomorphism of locally compact second-countable groups.  For
every given Haar measure \(\alpha\) on \(A\) and every
given Haar measure \(\gamma\) on \(C\), there is a Haar
measure \(\beta\) on \(B\) such that
\[
                       e_*(\alpha\times\beta)=\gamma.
 \tag{AMD.27}
\]
In particular the original \texttt{exists\_\allowbreak haar\_\allowbreak map\_\allowbreak prod\_\allowbreak eq}
is supplied with the prescribed \(A\)-measure unchanged.

\begin{proof}
Choose any Haar measure \(\beta_0\) on \(B\).  The
pushforward \(\gamma_0=e_*(\alpha\times\beta_0)\) is Haar.
Haar uniqueness gives \(\gamma=c\gamma_0\) for a finite
positive real number \(c\).  Put \(\beta=c\beta_0\).
It is Haar, and product integration and pushforward commute
with this finite scalar multiplication, giving (AMD.27).
The first factor \(\alpha\) is not rescaled.  This argument
also works when one or more of the groups is trivial; a Haar
measure there is a finite positive mass at its sole point.
\end{proof}
% END AMD13

% BEGIN AMD14
\paragraph{AMD14. The entire original prescribed-measure extension.}
\label{claim:AMD14}
For every original standard frame, prescribe Haar measures
\(\mu_{N,\infty},\mu_{H,\infty},\mu_N,\mu_H\) and a
family of Haar measures \(\mu_{S,v}\) on the exact local
groups.  There is an original \(W:\texttt{Meas\allowbreak W}\) with
\(W.\texttt{Is\allowbreak Text\allowbreak Norm}\) which has each of these five
specified entries exactly.  This is the whole original
\texttt{exists\_\allowbreak meas\allowbreak W\_\allowbreak prescribed}.

\begin{proof}
Apply AMD13 to \(\operatorname{split}_N^{-1}\) with the
given infinite measure and the given global \(N\)-measure.
This supplies a finite Haar measure \(\mu_{N,f}\) for which
the first identity in (AMD.23) holds with the prescribed
\(\mu_N\).  Apply it separately to
\(\operatorname{split}_H^{-1}\), obtaining \(\mu_{H,f}\)
and the second identity in (AMD.23).  Neither infinite measure
nor either global measure has changed.

Define \(\mu_{S,f}\) by AMD7, and use the supplied family
\(\mu_{S,v}\) as the local entry.  The same proof of the
eleven normalization clauses in AMD11 applies.  In particular
the finite \(S\)-measure is the exact product pushforward of
the two newly chosen finite factors.  The source imposes no
additional equality relating this finite measure to the
independently supplied local \(S\)-measures in this contract.
All five requested equalities therefore hold simultaneously.
\end{proof}
% END AMD14

\subsection{The actual unipotent rational quotient and its covolume}

% BEGIN AMD15
\paragraph{AMD15. A compact set of representatives for the original unipotent quotient.}
\label{claim:AMD15}
For every original standard frame, the rational subgroup
\(N(F)\subset N(\mathbb A_F)\) is exactly the source's
\(\texttt{rat\allowbreak Of}\).  It is countable, discrete and closed,
and there is a compact set \(K\subset N(\mathbb A_F)\)
such that
\[
                    N(\mathbb A_F)=N(F)K.
 \tag{AMD.28}
\]
This supplies the full original \texttt{nsub\_\allowbreak cocompact}.

\begin{proof}
The defining flag equations have coefficients in \(F\), after
writing \(E=F\oplus F\delta\).  A rational matrix whose
adelic extension satisfies these equations already satisfies
them over \(F\), since scalar extension is injective.  The
same applies to its unitary equation.  Thus intersection with
the source's rational ambient unitary subgroup is exactly
\(N(F)\), as asserted.

The polynomial logarithm of AMD6 identifies \(N(F)\) with
a finite-dimensional nilpotent \(F\)-Lie algebra
\(\mathfrak n(F)\), and similarly over \(\mathbb A_F\).
Choose an \(F\)-basis of that Lie algebra.  The diagonal
embedding of \(F\) in \(\mathbb A_F\) is discrete and
closed, so the same holds for \(F^{\dim\mathfrak n}\) and
hence for \(N(F)\).  Countability follows from countability
of the number field.

We give the needed compactness argument explicitly.  In the
additive case choose a closed fundamental parallelepiped for
\(\mathcal O_F\subset F_\infty\) and multiply it by
\(\prod_{v<\infty}\mathcal O_{F,v}\).  This is compact and
its translates by \(F\) cover \(\mathbb A_F\): the Chinese
remainder theorem, after clearing the finitely many denominators,
first makes every finite component integral, and translation by
\(\mathcal O_F\) then moves the infinite component into the
parallelepiped.  This also verifies the asserted discreteness
of \(F\): an infinite neighborhood avoiding the nonzero
integral lattice points, times the integral finite adeles,
meets the diagonal \(F\) only at zero.

Since \(\mathfrak n\) is nilpotent, it has a flag of
\(F\)-ideals
\(0=\mathfrak n_0\subset\mathfrak n_1\subset\cdots
\subset\mathfrak n_d=\mathfrak n\), of successive dimension
one, with \([\mathfrak n,\mathfrak n_i]\subseteq\mathfrak n_{i-1}\).
To construct it, choose a nonzero central \(F\)-vector and
induct in the quotient; its center is nonzero whenever the
nilpotent quotient is nonzero.  Exponentiation gives normal
subgroups \(N_i\), with each quotient the additive group.
Choose \(X_i\in\mathfrak n_i\) mapping to its fixed basis
in \(\mathfrak n_i/\mathfrak n_{i-1}\).  The finite
Baker--Campbell--Hausdorff formula, or induction through these
central quotients, shows that
\[
 (t_1,\ldots,t_i)\longmapsto
             \exp(t_1X_1)\cdots\exp(t_iX_i)
 \tag{AMD.29}
\]
is a polynomial bijection onto \(N_i\), with polynomial
inverse.  All its coefficients are in \(F\), so it applies
both to rational and to adelic points.

Inductively suppose \(N_{i-1}(\mathbb A_F)=N_{i-1}(F)K_{i-1}\)
with \(K_{i-1}\) compact.  Write an element of \(N_i\) as
\(u\exp(tX_i)\), and choose \(a\in F\) with \(t-a\)
in the fixed compact additive representative set \(C\).
Then
\[
 u\exp(tX_i)=\exp(aX_i)
   \bigl(\exp(-aX_i)u\exp(aX_i)\bigr)\exp((t-a)X_i).
 \tag{AMD.30}
\]
The middle factor is in \(N_{i-1}(\mathbb A_F)\), because
that subgroup is normal.  Reduce it by the induction hypothesis.
The resulting left factor is rational and the right factor lies
in the compact set \(K_{i-1}\exp(CX_i)\).  This proves
(AMD.28) at the next stage and hence at \(i=d\).
For the zero-dimensional group the compact set is \(\{1\}\).
\end{proof}
% END AMD15

% BEGIN AMD16
\paragraph{AMD16. The literal covolume-one Haar measure.}
\label{claim:AMD16}
For every original standard frame there is a Haar measure
\(\mu_N\) on \(N(\mathbb A_F)\) with
\(\texttt{covol}(\mu_N)=1\), for exactly the source's
choice-based covolume definition.  This is the whole original
\texttt{exists\_\allowbreak haar\allowbreak Measure\_\allowbreak }\allowbreak
\texttt{unipotent\allowbreak Radical\_\allowbreak }\allowbreak\texttt{covolume\_\allowbreak one}.

\begin{proof}
Put \(G=N(\mathbb A_F)\), \(\Gamma=N(F)\), and choose
any Haar measure \(\lambda\) on \(G\).  AMD15 shows that
\(\Gamma\backslash G\) is compact.  Choose a relatively
compact open neighborhood \(V\) of the identity with
\(VV^{-1}\cap\Gamma=\{1\}\).  Such a neighborhood exists
by discreteness.  The quotient map is open, so compactness
gives finitely many \(x_1,\ldots,x_m\) such that the images
of \(Vx_j\) cover the quotient.  Set
\[
 B_j=Vx_j\setminus\bigcup_{i<j}\bigcup_{\gamma\in\Gamma}
                                     \gamma Vx_i,
 \qquad B=\bigcup_{j=1}^m B_j.
 \tag{AMD.31}
\]
These are Borel sets, because \(\Gamma\) is countable.
No two points of one \(Vx_j\) have the same orbit, by the
condition on \(VV^{-1}\), and (AMD.31) removes every orbit
already represented at an earlier stage.  Thus \(B\) meets
each orbit exactly once.  It is a Borel fundamental domain
contained in the finite union of the compact sets
\(\overline Vx_j\).  Hence \(\lambda(B)<\infty\).
Also \(\lambda(B)>0\), since the countably many translates
\(\gamma B\) cover \(G\) and a Haar measure is nonzero.

Every measurable fundamental domain has this same measure.
Indeed partition one of them into its intersections with the
countably many translates of the other, translate each piece
back, and use left invariance and countable additivity.  The
same argument works for the source's almost-everywhere
disjointness and covering definition after discarding the
corresponding null sets.  In particular its chosen fundamental
domain has the finite positive measure \(\lambda(B)\).
Multiplying \(\lambda\) by a positive finite constant
preserves the null sets, the fundamental-domain property, and
the preceding equality of measures.

Define \(\mu_N=\lambda(B)^{-1}\lambda\).  It is Haar,
its fundamental domains have measure one, and the real value
of that measure is one.  The source's definition selects the
measure of a fundamental domain and takes its real part when
such a domain exists.  Existence and choice independence have
just been proved, so its value is exactly one for this
\(\mu_N\).  This proves the stated source contract, rather
than only compactness of the rational quotient.
\end{proof}
% END AMD16

% BEGIN AMD17
\paragraph{AMD17. The covolume normalization can coexist with every prescribed infinity and local measure.}
\label{claim:AMD17}
For an original standard frame, fix arbitrary infinity Haar
measures \(\mu_{N,\infty},\mu_{H,\infty}\), an arbitrary
global Haar measure \(\mu_H\), and arbitrary local Haar
measures \(\mu_{S,v}\).  One original normalized
\(W:\texttt{Meas\allowbreak W}\) retains all these entries and has
\(\texttt{covol}(W.\mu_N)=1\).  It can be chosen for each
original word using the same already fixed infinity bridge.

\begin{proof}
Choose the global \(\mu_N\) of AMD16 and apply AMD14 with
this measure and all the specified entries.  Its finite factor
absorbs whatever positive scalar is required; neither infinity
factor, the global \(H\)-measure, nor any local \(S\)-measure
is altered.  The resulting \(W\) has the eleven original
normalization clauses and the required covolume.  If the
infinity measures came from AMD9, their exact ambient ties
(AMD.20) therefore survive this replacement.  Applying the
construction independently at every word changes no bridge
and gives the literal preservation step used in
\texttt{Rescale\allowbreak Step3}.

The other inputs to that original step are distinct: the
compatible global orbital measures and stabilizer measures,
the normalized local fiber measures, and their identities from
the original Siegel--Weil premise.  This claim says that once
those prescribed \(H\)- and local \(S\)-Haar measures have
been produced, the unipotent covolume and all required infinity
ties can be imposed simultaneously.  It does not assume that
the prescribed measures already satisfy any additional orbital
or fiber identity in order to construct \(W\).
\end{proof}
% END AMD17

\subsection{Exact effect on the remaining global construction}

AMD12 removes the entire measured-block producer from the
remaining work, and AMD3 supplies the full adelic isomorphism,
including the rational and compact subsets.  Together with the
eight dependent blocks in ASG18, this gives every original
setting block except the normalized nonzero Bessel block
\(W_b\).  For any such original \(W_b\), the dependent
constructor now gives an actual \(X\) with its complete
\texttt{Meas\allowbreak Tie}.  The decay assertion associated with a
properly constructed \(W_b\) is not inferred from measure
existence.

The later \texttt{Rescale\allowbreak Step3} still needs its original
Siegel--Weil and orbital/fiber comparison suppliers.  AMD14,
AMD16, and AMD17 reconstruct its prescribed-measure and
unipotent-normalization portions while preserving the original
infinity measures.  The Witt/genus/level/representation
transport into the arithmetic columns also remains a separate
same-object construction.  None of these distinctions closes
\texttt{GLOBAL\_\allowbreak SPECTRAL}, the arithmetic counting obligation,
or the full RH/GRH reconstruction by itself.


\section{Arithmetic models for the original principal uniformization}
\label{sec:principal-uniformization-v48}

\status{This packet constructs arithmetic suppliers and individual
moduli points for the original principal-uniformization theorem.
The original field, host,
perfect alternating lattices, level ideals and Hermitian space are
retained.  The arithmetic classical input is Hasse--Minkowski,
already identified in MF4; the final subsection separately names
the classical complex-torus and algebraic-geometric foundations
used for individual moduli points.  The proofs below do not assume a Hermitian
local-to-global theorem or the desired principal uniformization.
The precise remaining holomorphic-family and normalized-frame inputs
are recorded at the end.  In particular, an arithmetic lattice model
alone does not establish the common principal-chart obligation.}

Put \(F=E^{\bar{\ } }\), \(R=\mathcal O_E\), and
\(R_0=\mathcal O_F\).  The original \texttt{Host} supplies a CM field
\(E\), an everywhere unramified quadratic extension \(E/F\), and
\(\mathrm{rt}^2=-d\), where \(d\) is a positive prime and
\(d\equiv7\pmod {16}\).  Its chosen complex square root is
\(\mathrm{rt}_0\).  Write
\[
 \Psi=\{\tau:E\longrightarrow\mathbb C:
                       \tau(\mathrm{rt})=\mathrm{rt}_0\}.
 \tag{PUN.1}
\]
Complex conjugation takes \(\mathrm{rt}\) to \(-\mathrm{rt}\), so
\(\Psi\) contains exactly one member of each conjugate pair.
Throughout this packet Hermitian forms are \emph{linear in the first
variable}, as in the original \texttt{herm}.  An original datum
\(D\) consists of a finite projective \(R\)-module \(L\), its perfect
alternating \(\mathbb Z\)-pairing \(B\), the adjoint identity
\(B(ax,y)=B(x,\bar a y)\), and the original compatible positive
complex structure with multiplicities \((p,q)\).

An original level is denoted \(l=(M,\mathfrak b,A)\).  In particular
\(M\ge3\) is odd, \(A>0\), and
\[
 \mathfrak b\bar{\mathfrak b}=MR,\qquad
 \mathfrak b+\bar{\mathfrak b}=R,\qquad
                  N=l.M'=\operatorname{lcm}(M,A).
 \tag{PUN.2}
\]
No freeness of \(L\) over \(R\) is imposed.  It is finite free as a
\(\mathbb Z\)-module.  The empty matrices and zero lattice have their
usual meanings in rank zero.

\subsection{The norm equation and the same Hermitian space}

% BEGIN PUN1
\paragraph{PUN1. The literal quadratic norm theorem.}
\label{claim:PUN1}
The classical Hasse--Minkowski theorem in MF4 implies the whole
original \texttt{Published.\allowbreak Hasse\allowbreak Norm\allowbreak St}: for every number field
\(K\), every \(\delta,x\in K^\times\), if
\(x=a^2-\delta b^2\) is solvable at every real embedding and every
finite completion, it is solvable in \(K\).

\begin{proof}
If \(\delta=s^2\), take
\(a=(x+1)/2\), \(b=(x-1)/(2s)\).  Otherwise apply
Hasse--Minkowski to the nondegenerate ternary quadratic form
\[
                  Q(U,V,W)=U^2-\delta V^2-xW^2.
 \tag{PUN.3}
\]
At the specified completions there is an isotropic vector with
\(W=1\).  At a complex completion a square root of \(x\) gives
one as well.  A nonzero global isotropic vector cannot have \(W=0\):
then \(V\ne0\) would make \(\delta=(U/V)^2\), and \(V=0\)
would force all three coordinates to vanish.  Division by \(W\)
gives the required solution.  This preserves every field and every
place quantified in the original statement, without an auxiliary
extension or an omitted dyadic case.
\end{proof}
% END PUN1

% BEGIN PUN2
\paragraph{PUN2. Hermitian local-to-global isometry by a quadratic induction.}
\label{claim:PUN2}
Let \(E/F\) be a CM quadratic extension and let \(h,h'\) be
nondegenerate Hermitian spaces of the same finite dimension.
If they are isometric at every finite completion of \(F\) and at
every real embedding, they are isometric over \(E\).

\begin{proof}
Diagonalize \(h'\) and let \(a\in F^\times\) be its first
coefficient.  The quadratic form \(q(v)=h(v,v)\) on the underlying
\(F\)-space of \(h\) is nondegenerate: its polar form is
\(\operatorname{Tr}_{E/F}h(v,w)\), whose radical is zero by the
nondegeneracy of \(h\) and of the field trace.  The form
\(q(v)-a t^2\) is locally isotropic everywhere.  Hasse--Minkowski
supplies a nonzero global zero \((v,t)\).

If \(t\ne0\), then \(h(v/t,v/t)=a\).  If \(t=0\), choose
\(w\) with \(h(v,w)=1\) and replace \(w\) by
\(w-h(w,w)v/2\).  The resulting \(w_0\) has norm zero and
\(h(v,w_0)=1\).  Thus \(v+(a/2)w_0\) has norm \(a\).
In either case split off a global line of norm \(a\).

At each completion, two vectors of the same nonzero norm are
carried to one another by an isometry.  For a field extension this
is the explicit reflection calculation of MF3, with the arguments
interchanged to convert its conjugate-linear-first convention to
the present convention.  If its first denominator vanishes, the
intermediate vector \(-w\) makes the alternative denominator
\(2a\ne0\).  For a split quadratic algebra a Hermitian space is
a perfect pairing between two vector spaces; extend the specified
nonzero paired vectors to dual bases.  This gives the same assertion
without dividing by a possible zero divisor.  Consequently the two
orthogonal complements are locally isometric at every place.
Induction gives an isometry of the complements, and adjoining the
chosen lines gives the desired isometry.  The dimension-zero case
starts the induction.  This argument uses only the stated quadratic
theorem; no separate Landherr theorem is being assumed.
\end{proof}
% END PUN2

% BEGIN PUN3
\paragraph{PUN3. The signature forced by the original positive complex structure.}
\label{claim:PUN3}
Suppose \(\phi:L\hookrightarrow E^r\) has full span, \(h=h_G\)
is nondegenerate, \(\bar\xi=-\xi\ne0\), and
\[
              B(x,y)=\operatorname{Tr}_{E/\mathbb Q}
                              (\xi h(\phi x,\phi y)).
 \tag{PUN.4}
\]
Then \(r=p+q\).  At \(\tau\in\Psi\), if
\(\operatorname{Im}\tau\xi>0\), \(h_\tau\) has \(q\)
positive and \(p\) negative directions; if the imaginary part is
negative, these multiplicities are reversed.  In particular two
original data of the same \((p,q)\), represented using the same
\(\xi\), have identical archimedean Hermitian signatures.

\begin{proof}
Tensoring the injection with \(E\) gives an isomorphism, so the
rank is the original rank \(p+q\).  Under
\(E\otimes_{\mathbb Q}\mathbb R\simeq\prod_{\tau\in\Psi}\mathbb C\),
the pairing (PUN.4) becomes
\(\sum_\tau2\operatorname{Re}(\tau\xi\,h_\tau)\).
The original complex structure \(J\) commutes with \(E\), hence
with the idempotents and complex scalars in this product.  It
therefore acts by a complex-linear operator \(J_\tau\) on each
\(\mathbb C^r\), with \(J_\tau^2=-1\).

The two real trace identities in \texttt{Datum.\allowbreak dom}, first on
\(E\) and then by real-linear extension on
\(E\otimes\mathbb R\), give
\[
 \operatorname{tr}_{\mathbb C}J_\tau=i(p-q).
 \tag{PUN.5}
\]
Together with \(J_\tau^2=-1\), this means that its \(+i\) and
\(-i\) eigenspaces have dimensions \(p\) and \(q\).
The invariance of \(B\), tested also after multiplying one argument
by arbitrary complex scalars in the same product factor, gives
\(h_\tau(J_\tau u,J_\tau v)=h_\tau(u,v)\).  Thus these two
eigenspaces are orthogonal.  Write \(\tau\xi=is\) with
\(s\in\mathbb R^\times\).  On them the original positivity
\(B(Ju,u)>0\) says respectively
\(-2s h_\tau(u,u)>0\) and \(2s h_\tau(u,u)>0\).
The asserted signature follows.  This also proves the required
diagonal normal form with first \(q\) signs positive when \(s>0\):
choose orthogonal bases in the two definite eigenspaces and rescale
each vector by the positive square root of its absolute norm.
Conjugate embeddings have the same signatures by conjugating all
coordinates.
\end{proof}
% END PUN3

% BEGIN PUN4
\paragraph{PUN4. Local isometry from the perfect integral trace pairing.}
\label{claim:PUN4}
Under the original \texttt{Host.\allowbreak unram}, let two full lattices in
rank \(r\) Hermitian spaces satisfy (PUN.4), with the same nonzero
\(\xi\), and let both resulting alternating integral pairings be
perfect.  At every finite place of \(F\) their Hermitian spaces
are isometric, including the places above \(2\).

\begin{proof}
Let \(C=\xi^{-1}\mathfrak D_{E/\mathbb Q}^{-1}\).  The exact
trace-duality calculation of PCL2--PCL3 gives, for either lattice,
\[
          L=\{v:h(v,L)\subseteq C\}.
 \tag{PUN.6}
\]
This identity commutes with localization and completion: a finite
projective lattice and its dual are finite projective, and the
pairing matrix may be computed in a local basis.  Fix a finite
place \(v\) of \(F\), put \(A=\widehat{\mathcal O}_{F,v}\)
and \(B=\mathcal O_E\otimes_{\mathcal O_F}A\).  The original
unramified hypothesis means that \(B\) is either \(A\times A\)
or the integers of an unramified quadratic extension.  In either
case every conjugation-stable fractional ideal of \(B\) is
\(cB\) for some \(c\in F_v^\times\): in the field case choose
a power of a uniformizer of \(F_v\), and in the split case
conjugation stability makes the two valuations equal.

The ideal \(C\) is conjugation stable, since
\(\bar\xi=-\xi\) and the different is stable.  Choose the same
\(c\) for both lattices.  Equation (PUN.6) says exactly that
\(c^{-1}h\) is a unimodular Hermitian form on a free \(B\)-module
of rank \(r\).  MF2 diagonalizes this form over \(B\) and removes
each diagonal unit by a unit norm.  Its proof uses the surjective
residue trace and successive corrections, so it includes residue
characteristic \(2\).  In the split case dual bases for the two
factors give the same conclusion.  Both scaled forms are therefore
isometric to the same standard form.  Multiplication by \(c\)
gives the asserted local isometry of the original forms.
\end{proof}
% END PUN4

% BEGIN PUN5
\paragraph{PUN5. Every component lattice embeds in the prescribed Hermitian space.}
\label{claim:PUN5}
Let \(D,D'\) be original data for the same host, with
\(D'.p=D.p\) and \(D'.q=D.q\).  Suppose the prescribed
\((G,\phi,\xi)\) satisfies the full original Hermitian-coordinate
contract (PUN.4), including injectivity, full span, Hermitian
symmetry and \(\det G\ne0\).  There is an \(R\)-linear
injection \(\phi':D'.L\hookrightarrow E^r\) with full span and
\[
        B'(x,y)=\operatorname{Tr}_{E/\mathbb Q}
                      (\xi h_G(\phi'x,\phi'y)).
 \tag{PUN.7}
\]
The matrix \(G\) and scalar \(\xi\) are unchanged.

\begin{proof}
For \(\xi\ne0\), PCL1 supplies Hermitian coordinates
\((G',\psi)\) for \(D'\) using this same \(\xi\).  The ranks
agree by PUN3.  The two forms have the same infinite signatures
by PUN3 and are isometric at every finite place by PUN4.  PUN2
provides an \(E\)-linear isometry
\(g:(E^r,h_{G'})\longrightarrow(E^r,h_G)\).
Take \(\phi'=g\psi\).  It is injective, has full span, and
satisfies (PUN.7) by the isometry identity.

For the literal source's additional case \(\xi=0\), (PUN.4)
and perfection imply \(D.L=0\); full span then gives \(r=0\).
The original rank equality and the equal signatures give
\(D'.L=0\), since a finite projective module over a domain of
rank zero is zero.  The unique map into \(E^0\) proves the
assertion.  Thus this is the full same-space producer, with a
quadratic local-to-global proof replacing its longer norm-ideal
implementation.
\end{proof}
% END PUN5

\subsection{The actual principal group and all lattice-model fields}

% BEGIN PUN6
\paragraph{PUN6. The original half-level condition is the full congruence condition.}
\label{claim:PUN6}
Let \(g\) be an \(R\)-linear automorphism of the original \(L\)
preserving \(B\).  If
\((g-1)\bar{\mathfrak b}L\subseteq ML\), then
\((g-1)L\subseteq ML\).  Consequently the original group
\(\operatorname{princ}(l,G,D,\phi)\) is exactly
\[
 \{g\in U(G):g\Lambda=\Lambda,
                        (g-1)\Lambda\subseteq N\Lambda\},
 \qquad \Lambda=\phi L,\quad N=\operatorname{lcm}(M,A).
 \tag{PUN.8}
\]

\begin{proof}
Choose \(i\in\mathfrak b\) and \(j\in\bar{\mathfrak b}\)
with \(i+j=1\).  Put \(y=gx-x\).  The assumption gives
\(jy=Me_1\).  For any \(z\in L\), write
\(g(\bar i z)=\bar i z+Me_2\); this is permitted because
\(\bar i\in\bar{\mathfrak b}\).  The adjoint and isometry
identities give
\[
 B(iy,z)=B(gx-x,\bar i z)=-M B(gx,e_2).
 \tag{PUN.9}
\]
Thus \(z\mapsto B(iy,z)/M\) is an integral linear functional.
Perfection gives \(w\in L\) representing it, and injectivity of
the pairing gives \(iy=Mw\).  Adding the equations for \(iy\)
and \(jy\) proves \(y=M(w+e_1)\).

An element stabilizing \(\Lambda\) induces an actual automorphism
of \(L\), because \(\phi\) is an isomorphism onto its range.
Its Hermitian isometry preserves \(B\), so the preceding argument
applies to the literal half-level condition.  The additional
\(A\)-condition gives
\((g-1)\Lambda\subseteq M\Lambda\cap A\Lambda=N\Lambda\),
where the last equality follows in a \(\mathbb Z\)-basis of
\(\Lambda\).  The converse follows from \(M\mid N\),
\(A\mid N\).  This proves equality with the original group,
not just a finite-index inclusion.
\end{proof}
% END PUN6

% BEGIN PUN7
\paragraph{PUN7. Finite index with the source's coset order.}
\label{claim:PUN7}
The actual principal group has finite index in
\(U(G,\Lambda)=\{g\in U(G):g\Lambda=\Lambda\}\).
There is a finite set \(S\subseteq\operatorname{GL}_r(E)\)
such that every \(g\in U(G,\Lambda)\) has an \(s\in S\) with
\(s^{-1}g\in\operatorname{princ}(l,G,D,\phi)\).

\begin{proof}
The set \(\Lambda/N\Lambda\) is finite since \(N>0\) and
\(\Lambda\) is finite free over \(\mathbb Z\).  Reduction gives
a homomorphism from \(U(G,\Lambda)\) to the automorphism group
of this finite module.  Its kernel is exactly (PUN.8).  Choose
one representative of each element of the finite image.  Equal
reductions of \(g\) and its representative \(s\) give
\(s^{-1}g\) in that kernel, which is the precise order required
by the original \texttt{Gamma\_\allowbreak index} field.
\end{proof}
% END PUN7

% BEGIN PUN8
\paragraph{PUN8. Freeness on the original full flag domain.}
\label{claim:PUN8}
For every \(g\) in the original principal group, every
\(x\in\operatorname{GG}(\Psi,G,p_0)\), and every original
upper-triangular \(b\in\operatorname{Bor}(\Psi,r)\),
\(g\cdot x=x b\) implies \(g=1\).

\begin{proof}
At each \(\tau\in\Psi\), the successive column spaces of
\(x_\tau\) form a complete flag whose restrictions of
\(h_\tau\) are nondegenerate: their determinants are precisely
the nonzero leading minors in the definition of \(\operatorname{GG}\).
An isometry preserving this flag preserves the one-dimensional
orthogonal difference between each two consecutive terms.  On
each such nondegenerate line it acts by a scalar of absolute value
one.  It is therefore diagonalizable in a fixed basis with all
powers bounded.  The equality \(g x=x b\) ensures flag
preservation.  Conjugate embeddings give the same bounds; hence
the real-linear operators induced by all powers of \(g\) are
bounded on \(\Lambda\otimes\mathbb R\).

Choose a \(\mathbb Z\)-basis of \(\Lambda\).  Its integral
endomorphisms form a discrete lattice in the real endomorphism
space, so a bounded set contains only finitely many of them.
Two nonnegative powers of \(g\) coincide.  Since \(g\) is
invertible, it has finite order on \(\Lambda\), and hence on
\(E^r\).  By PUN6 it acts as the identity modulo \(M\).

For completeness the required torsion argument uses only that
\(M\) is odd and at least three.  Choose an odd prime \(p\mid M\).
If a nonidentity integral matrix of finite order is the identity
modulo \(p\), a power has prime order \(q\) and is still of
the form \(1+p^a C\), where \(a\ge1\) is maximal and an entry
of \(C\) is not divisible by \(p\).  Expand its \(q\)-th power.
For \(q\ne p\), division by \(p^a\) and reduction modulo \(p\)
gives \(qC=0\).  For \(q=p\), division by \(p^{a+1}\) gives
\(C=0\pmod p\): every other term has an additional factor
\(p\), including the last because \(a(p-1)-1\ge1\).
Both are contradictions.  Thus \(g=1\), including when some
positive or negative signature multiplicity is zero.  In rank
zero the assertion is immediate.
\end{proof}
% END PUN8

% BEGIN PUN9
\paragraph{PUN9. One transvection modulus for all original integral isotropic vectors.}
\label{claim:PUN9}
For the same \(\Lambda,G,\xi_0\), with
\(\bar\xi_0=-\xi_0\), there is one integer \(m>0\) such that
for both \(\kappa_0=2\xi_0\) and \(\kappa_0=-2\xi_0\),
every \(v\in\Lambda\) with \(h_G(v,v)=0\), and every
\(s\in E\) with \(\bar s=s\) and \(s\) integral over
\(\mathbb Z\), the actual matrix
\(\operatorname{transvMat}(G,v,ms\kappa_0)\) belongs to the
original principal group.

\begin{proof}
Choose finite \(R\)-generators of \(\Lambda\).  Clear the
denominators of all their pairings multiplied by \(2\xi_0\),
using one positive rational integer \(d_0\).  Sesquilinearity
then gives
\(d_0(2\xi_0)h_G(\Lambda,\Lambda)\subseteq R\).
Set \(m=N d_0\).  For \(\alpha=ms\kappa_0\) define
\[
        T_\alpha(x)=x+\alpha h_G(x,v)v.
 \tag{PUN.10}
\]
This is exactly the original matrix: its \((a,b)\)-entry is
\(\delta_{ab}+\alpha v_a\sum_kG_{bk}\bar v_k\).
As \(\bar\alpha=-\alpha\) and \(h_G(v,v)=0\), expansion
gives \(h_G(T_\alpha x,T_\alpha y)=h_G(x,y)\), and
\(T_\alpha^{-1}=T_{-\alpha}\).  Because \(s\in R\), both
maps preserve \(\Lambda\), and
\((T_\alpha-1)\Lambda\subseteq N\Lambda\).
PUN6 puts them in the actual principal group.  The same \(m\)
works for both signs, every isotropic vector and every such \(s\);
it was chosen before these variables.
\end{proof}
% END PUN9

% BEGIN PUN10
\paragraph{PUN10. The whole original lattice-model supplier.}
\label{claim:PUN10}
Under the source's hypotheses for
\texttt{lattice\allowbreak Model\_\allowbreak of\_\allowbreak same\_\allowbreak signature\_\allowbreak hasse\allowbreak Norm}, every
datum \(D'\) with the same \((p,q)\) has an actual
\texttt{Lat\allowbreak Model}\((l,D',\xi_0,G)\).  All its fields hold in
the prescribed space and on the original component lattice.

\begin{proof}
Apply PUN5 with \(\xi=-\xi_0\) to obtain the actual map
\(\phi':D'.L\hookrightarrow E^r\).  It has the original
perfect trace identity and full span.  Finite generation of
\(\Lambda'=\phi'D'.L\) permits a single nonzero integral
denominator for \(h_G(\Lambda',\Lambda')\), by clearing the
finitely many pairings of generators.  PUN7 supplies the exact
finite-index field; PUN8 supplies freeness for every original
flag frame and triangular multiplier; PUN9 supplies the one
positive transvection modulus with both original signs of
\(\kappa_0\).  These are precisely the fields of
\texttt{Lat\allowbreak Model}; none is inferred merely from commensurability
of an auxiliary lattice.
\end{proof}
% END PUN10

% BEGIN PUN11
\paragraph{PUN11. Actual flag-domain models for any supplied family uniformization.}
\label{claim:PUN11}
Fix the original \(X_o,c,\sigma\), a \emph{supplied}
\(F_u:\texttt{Family\allowbreak Unif}\) on these objects, and the original
Hermitian coordinates \((G,\phi,\xi_0)\), with
\(\operatorname{Im}\tau\xi_0<0\) for every \(\tau\in\Psi\).
There is a \(\texttt{Models\allowbreak Of}\ F_u\ r\) satisfying the
whole original \texttt{Model\allowbreak Base} and \texttt{Neg\allowbreak Xi}, and every
one of its matrices and imaginary scalars is exactly \(G\) and
\(\xi_0\).

\begin{proof}
The family record has finitely many component data \(D_i\),
each with \((D_i.p,D_i.q)=(D.p,D.q)\).  Apply PUN10 to each
component, choosing its injection \(\phi_i\) and put
\(\Lambda_i=\phi_iD_i.L\).  The inverse map onto its range
gives the required actual integral equivalence
\(D_i.L\simeq_R\Lambda_i\).

Define \(X_i\) using \(\Psi\) in (PUN.1), \(G\),
\(p_0=q\), \(\Lambda_i\), the principal group of PUN6, and
the constant character \(\chi(\lambda,g)=1\).  PUN3 applied
to \(-\xi_0\) gives the required signature.  PUN10 gives
finite generation, full span, the denominator, finite index and
freeness.  Take \(\kappa_0=2\xi_0\), whose imaginary part is
negative at every \(\tau\in\Psi\), and the modulus of PUN9.
All character laws hold identically, with finite-order exponent
one.  This constructs every field of the original \texttt{Data}.

The \texttt{Model\allowbreak Base} trace and group fields are their defining
identities.  Its domain-stability field also holds literally.
Left multiplication by a unitary element preserves each Gram
matrix.  For right multiplication by an invertible upper
triangular matrix, the leading \(j\)-column Gram determinant
is multiplied by the positive factor
\(\prod_{a<j}|b_{aa}|^2\).  Thus every defining strict sign
inequality of \(\operatorname{GG}\) is preserved.  Rank,
CM type, and \texttt{Neg\allowbreak Xi} are the prescribed choices above.
This proves the exact arithmetic model producer; the existence
of \(F_u\) itself has not been assumed as a conclusion.
\end{proof}
% END PUN11

\subsection{The exact cusp-center lattice, in every cusp rank}

The next statements apply to \emph{every} original
\(X:\texttt{Data}\ E\ r\), not only to PUN11's models.
Fix an original cusp \(\mathrm{cu}:\texttt{Cusp\allowbreak FD}\ X\ \rho\),
with linearly independent totally isotropic vectors
\(e_1,\ldots,e_\rho\), and write
\(Z(\beta)=\texttt{z\allowbreak Mat\allowbreak E}\ X\ e\ \beta\).

% BEGIN PUN12
\paragraph{PUN12. The cusp-center parameter space and its faithful additive action.}
\label{claim:PUN12}
The original skew-Hermitian parameter space
\[
 \mathfrak z=\{\beta\in M_\rho(E):\bar\beta=-\beta^t\}
 \tag{PUN.11}
\]
has \(\mathbb Q\)-dimension \(\rho^2[E:\mathbb Q]/2\).
The original matrices satisfy
\[
 Z(0)=1,\quad Z(\beta+\beta')=Z(\beta)Z(\beta'),\quad
 Z(\beta)^{-1}=Z(-\beta),
 \tag{PUN.12}
\]
and \(\beta\mapsto Z(\beta)\) is injective.
For \(\beta\in\mathfrak z\), \(Z(\beta)\) is an isometry
of the original Hermitian form.

\begin{proof}
Choose a nonzero anti-invariant \(w\in E\).  Then
\(E=F\oplus Fw\).  Each diagonal entry of \(\beta\) lies in
\(Fw\), and each pair above and below the diagonal is determined
by one arbitrary element of \(E\).  This gives dimension
\(\rho[F:\mathbb Q]+\rho(\rho-1)[E:\mathbb Q]/2
=\rho^2[E:\mathbb Q]/2\).

In the source's convention the matrix acts by
\[
 Z(\beta)x=x+\sum_a\Big(\sum_b\beta_{ab}h_G(x,e_b)\Big)e_a.
 \tag{PUN.13}
\]
Every displacement is in the span of the \(e_a\), and all
functionals \(h_G(-,e_b)\) vanish on that span.  Consequently
any product of two displacement operators is zero, proving
(PUN.12).  Nondegeneracy of \(h_G\) and independence of the
\(e_b\) make \(x\mapsto(h_G(x,e_b))_b\) surjective.
If the displacement is zero, independence of the \(e_a\)
therefore gives \(\beta=0\).  Finally, expand the pairing of
the two sides of (PUN.13).  The term involving two
displacements vanishes by total isotropy.  The two remaining
cross terms cancel precisely because
\(\beta_{ab}+\bar\beta_{ba}=0\).  This proves the isometry
claim with the actual matrix and its actual order.
\end{proof}
% END PUN12

% BEGIN PUN13
\paragraph{PUN13. Finite generation and full rational span of the actual center subgroup.}
\label{claim:PUN13}
The original subgroup
\[
 Z_\Gamma=\texttt{z\allowbreak Gam}\ X\ \mathrm{cu}
       =\{\beta\in\mathfrak z:Z(\beta)\in\Gamma\}
 \tag{PUN.14}
\]
is a finitely generated abelian group, and every
\(\beta\in\mathfrak z\) has a positive integral multiple in
\(Z_\Gamma\).  In particular its rational span is exactly
\(\mathfrak z\).

\begin{proof}
PUN12 shows that (PUN.14) is an additive subgroup.  Choose finite
\(R\)-generators \(v_1,\ldots,v_t\) of the original
\(\Lambda\), and map
\[
 \beta\longmapsto
        ((Z(\beta)-1)v_1,\ldots,(Z(\beta)-1)v_t).
 \tag{PUN.15}
\]
This is additive and injective: the \(v_i\) span \(E^r\),
and PUN12 is faithful.  For \(\beta\in Z_\Gamma\) its
image lies in \(\Lambda^t\).  The latter is finite free over
\(\mathbb Z\); all its subgroups are finitely generated.
This proves finite generation without first assuming a lattice
or a bounded fundamental domain for the center.

For any \(\beta\in\mathfrak z\), clear the finitely many
denominators of \((Z(\beta)-1)v_i\) relative to \(\Lambda\).
There is a positive integer \(d_0\) with
\(d_0(Z(\beta)-1)\Lambda\subseteq\Lambda\).
By (PUN.12), \(Z(jd_0\beta)\) and its inverse preserve
\(\Lambda\) for every integer \(j\), and PUN12 says they
are unitary.  The original \texttt{Gamma\_\allowbreak index} field gives
finitely many cosets of \(\Gamma\) in the unitary lattice
stabilizer.  Two of the nonnegative powers
\(Z(jd_0\beta)\) have the same coset.  Their quotient is
\(Z((j-i)d_0\beta)\in\Gamma\), for some \(j>i\).
Thus \((j-i)d_0\beta\in Z_\Gamma\), proving the required
positive multiple and the full rational span.  Neither normality
of \(\Gamma\) nor a uniform congruence level was used.
\end{proof}
% END PUN13

% BEGIN PUN14
\paragraph{PUN14. The literal all-rank \texttt{Lat\allowbreak Z} producer.}
\label{claim:PUN14}
For every CM number field \(E\), every \(r,\rho\in\mathbb N\),
every component type \(I\), and every actual family
\(X:I\to\texttt{Data}\ E\ r\), the original
\(\texttt{FJ.\allowbreak Lat\allowbreak Z}\ X\ \rho\) holds.

\begin{proof}
For each original cusp of each component, PUN13 identifies the
actual \(Z_\Gamma\) as a finitely generated subgroup of a
\(\mathbb Q\)-vector space.  It is torsion free, so the
structure theorem for finitely generated abelian groups gives
a \(\mathbb Z\)-basis.  Its length is the dimension of its
rational span, hence by PUN12--PUN13 it is exactly
\(\rho^2[E:\mathbb Q]/2\).  The original CM-type condition
pairs the embeddings into two complementary equally large sets,
so \(|X.\Psi|=[E:\mathbb Q]/2\).  The basis length is therefore
the literal \(\texttt{n\allowbreak Z}\ X\ \rho=|X.\Psi|\rho^2\).
Regard the basis vectors as the
original skew-Hermitian matrices.  For every matrix \(\beta\),
membership in \(\texttt{z\allowbreak Gam}\) is equivalent to membership
in their \(\mathbb Z\)-span, and they are \(\mathbb Z\)-linearly
independent.  These are precisely the two fields of the original
\texttt{Is\allowbreak ZBasis}.  The construction can be chosen for each cusp,
which is exactly \texttt{Lat\allowbreak Z}'s quantifier order.  If \(\rho=0\)
the zero group has its empty basis; if no such cusp exists the
corresponding assertion is vacuous.  No rank or component
nonemptiness restriction has been introduced.
\end{proof}
% END PUN14

\subsection{Periods and the level-one exponent dictionary}

% BEGIN PUN15
\paragraph{PUN15. Periods for every original ordered type function.}
\label{claim:PUN15}
For every original datum \(D\), and every original
\(\mathrm{ty}:\{1,\ldots,g\}\to\operatorname{Hom}(E,\mathbb C)\)
with \(\texttt{Is\allowbreak Type\allowbreak Fn}(H,p,q,\mathrm{ty})\),
\[
 g=[E:\mathbb Q](p+q)/2,
 \qquad \exists P:L\longrightarrow\mathbb C^g,
                           \texttt{Is\allowbreak Period}(D,\mathrm{ty},P).
 \tag{PUN.16}
\]
Such type functions exist.  This proves the two unconditional
period-existence clauses of the original
\texttt{type\allowbreak Fn\_\allowbreak period\_\allowbreak exist}, without assuming \texttt{N1c}.

\begin{proof}
Let \(V=L\otimes_{\mathbb Z}\mathbb R\), and give it the
complex structure supplied in \texttt{Datum.\allowbreak dom}, declaring
\(i v=Jv\).  The \(E\)-action is complex linear.  As \(E\)
is separable, its complexification splits into its embedding
eigenspaces.  For \(a\in E\) the original trace identities
give
\[
 \operatorname{tr}_{\mathbb C}(a\mid V,J)
 =\tfrac12\bigl(\operatorname{tr}_{\mathbb R}a
                  -i\operatorname{tr}_{\mathbb R}(Ja)\bigr)
 =\sum_{\tau\in\Psi}
                       \bigl(p\tau(a)+q\bar\tau(a)\bigr).
 \tag{PUN.17}
\]
The distinct embeddings are linearly independent, so the
\(\tau\)-eigenspace has exactly the prescribed dimension.
Choosing a basis in each eigenspace, and ordering it according
to the supplied \(\mathrm{ty}\), gives an \(E\)-equivariant
complex-linear isomorphism \(T:(V,J)\simeq\mathbb C^g\).
It also gives the dimension formula; choosing any ordering
proves the existence of a type function.

Define \(P(x)=T(1\otimes x)\).  This is additive, injective,
and has the exact coordinate action
\(P(ax)_j=\mathrm{ty}_j(a)P(x)_j\).  A \(\mathbb Z\)-basis
of \(L\) is a real basis of \(V\), so \(P(L)\) is discrete
and spans \(\mathbb C^g\) over \(\mathbb R\).
On \(V\) define
\[
                 H(u,v)=B(Ju,v)+iB(u,v).
 \tag{PUN.18}
\]
The relations \(J^2=-1\), \(B(Ju,Jv)=B(u,v)\) and
alternation show that \(H\) is Hermitian and complex linear
in the first variable.  Its real diagonal value is
\(B(Ju,u)>0\) for \(u\ne0\).  Transporting it through \(T\)
gives a positive Hermitian matrix \(H_m\) satisfying
\(\operatorname{Im}H_m(Px,Py)=B(x,y)\).  These are all the
literal fields of \texttt{Is\allowbreak Period}.  The construction also
works for the empty type and zero lattice.
\end{proof}
% END PUN15

% BEGIN PUN16
\paragraph{PUN16. An adapted cusp pseudobasis and surjectivity of the exact pairing map.}
\label{claim:PUN16}
Assume the original \(\Lambda\subset E^r\) has the perfect
pairing \(B(x,y)=\operatorname{Tr}_{E/\mathbb Q}(-\xi_0h_G(x,y))\).
For any rational totally isotropic subspace \(W\), put
\(C=\Lambda\cap W\).  It has a pseudobasis
\[
                    C=\bigoplus_{j=1}^{\rho}\mathfrak a_j e_j,
                    \qquad \mathfrak a_j\ne0,
 \tag{PUN.19}
\]
with \(e_j\) an \(E\)-basis of \(W\).  For this basis the
map
\[
 \Lambda\longrightarrow\prod_j
        \mathfrak D^{-1}\bar{\mathfrak a}_j^{-1},
 \quad x\longmapsto(-\xi_0h_G(x,e_j))_j,
 \qquad\mathfrak D=\mathfrak D_{E/\mathbb Q},
 \tag{PUN.20}
\]
is onto.

\begin{proof}
The submodule \(C\) is finite and torsion free over the Dedekind
domain \(R\), hence projective.  It spans \(W\): an arbitrary
vector of \(W\) has a nonzero integral multiple in \(\Lambda\)
by fullness of \(\Lambda\).  To obtain (PUN.19), choose a
nonzero coordinate functional on \(W\).  Its image on \(C\)
is a nonzero fractional ideal, hence projective, so the exact
sequence onto that ideal splits.  The kernel has rank one less.
Inducting gives the direct sum of fractional ideals and its
corresponding rational basis.  This argument permits nonprincipal
ideals and does not make \(C\) free over \(R\).

Moreover \(C\) is a primitive \(\mathbb Z\)-sublattice of
\(\Lambda\): if \(n x\in C\) with \(n\ne0\) and
\(x\in\Lambda\), then \(x\in W\).  Thus
\(\Lambda/C\) is finite free over \(\mathbb Z\), and every
integral functional on \(C\) extends to \(\Lambda\).
For any tuple \(t_j\in\mathfrak D^{-1}\bar{\mathfrak a}_j^{-1}\),
the formula
\[
          \sum_j a_j e_j\longmapsto
                  \operatorname{Tr}_{E/\mathbb Q}
                                      \bigl(\sum_j t_j\bar a_j\bigr)
 \tag{PUN.21}
\]
is integral on \(C\).  Extend it to \(\Lambda\) and use
perfection to represent the extension by \(B(x,-)\).
Testing (PUN.21) with one coefficient in \(\mathfrak a_j\)
at a time and using nondegeneracy of the trace pairing shows
\(-\xi_0h_G(x,e_j)=t_j\) for all \(j\).
Conversely these coordinates necessarily lie in the displayed
ideals by integrality on \(C\).  This proves the exact image,
not only its rational span.
\end{proof}
% END PUN16

% BEGIN PUN17
\paragraph{PUN17. The actual center congruence is an entrywise fractional-ideal condition.}
\label{claim:PUN17}
Suppose \(\Gamma\) is the original principal group and \(A=1\).
In an adapted cusp basis as in (PUN.19), set
\[
            I_{ij}=M\mathfrak a_i\mathfrak D\bar{\mathfrak a}_j.
 \tag{PUN.22}
\]
For a skew-Hermitian matrix \(\beta\),
\[
 \beta\in Z_\Gamma
       \quad\Longleftrightarrow\quad
           \beta_{ij}/\xi_0\in I_{ij}\quad\hbox{for every }i,j.
 \tag{PUN.23}
\]

\begin{proof}
PUN6 identifies \(\Gamma\) with the full congruence kernel
modulo \(M\).  A displacement \((Z(\beta)-1)x\) lies in
\(W\).  Since \(M\Lambda\cap W=MC\), the congruence
condition is exactly
\[
  \sum_j\beta_{ij}h_G(x,e_j)\in M\mathfrak a_i
                      \quad(x\in\Lambda,\ 1\le i\le\rho).
 \tag{PUN.24}
\]
PUN16 says that \(-\xi_0h_G(x,e_j)\) ranges independently
over \(\mathfrak D^{-1}\bar{\mathfrak a}_j^{-1}\).
Consequently (PUN.24) is equivalent, entry by entry, to
\[
 (\beta_{ij}/\xi_0)
        \mathfrak D^{-1}\bar{\mathfrak a}_j^{-1}
                                      \subseteq M\mathfrak a_i,
\]
which is exactly (PUN.23), by invertibility of nonzero fractional
ideals.  Conversely the same calculation proves the congruence
for every \(x\); the inverse has displacement with the opposite
sign and hence preserves \(\Lambda\) too.  PUN12 supplies
unitarity.  Thus no stabilizer condition was lost in reducing
to entries.  Notice \(I_{ji}=\bar I_{ij}\), because the
different is conjugation stable.
\end{proof}
% END PUN17

% BEGIN PUN18
\paragraph{PUN18. The factor of two in the exact Fourier phase.}
\label{claim:PUN18}
Let \(\chi\) be Hermitian and let \(I_{ij}\) be as in
(PUN.22).  The phase in the original \texttt{pr\allowbreak E} satisfies
\[
 \bigl(\operatorname{prE}(\chi,\beta)\in\mathbb Z
                 \text{ for every }\beta\in Z_\Gamma\bigr)
 \quad\Longleftrightarrow\quad
 \operatorname{Tr}_{E/\mathbb Q}(\chi_{ij}s)\in\mathbb Z
                 \text{ for every }s\in I_{ij},\ i,j.
 \tag{PUN.25}
\]
This includes the original choices \(\kappa_0=\pm2\xi_0\).

\begin{proof}
Set \(S=\beta/\xi_0\).  Then \(S\) is Hermitian, and PUN17
is the condition \(S_{ij}\in I_{ij}\).  The definition of
\texttt{pr\allowbreak E}, summing the real parts over one CM type, gives
\[
 \operatorname{prE}(\chi,\beta)
       =\pm\frac12\operatorname{Tr}_{E/\mathbb Q}
                              \Bigl(\sum_{i,j}\chi_{ij}S_{ij}\Bigr).
 \tag{PUN.26}
\]
The sign is immaterial for integral membership.  If \(i\ne j\),
test (PUN.26) on the matrix with entries \(s,\bar s\) in
positions \((i,j),(j,i)\).  Its half-trace is
\(\operatorname{Tr}_{E/\mathbb Q}(\chi_{ij}s)\).
For a diagonal position test the real entry \(s+\bar s\).
Because \(\chi_{ii}\in F\), its half-trace is the same
quantity.  This proves the forward implication even at primes
above \(2\).

Conversely the right side implies
\(\chi_{ij}I_{ij}\subseteq\mathfrak D^{-1}\), since
each ideal is closed under multiplication by \(R\).
For an admissible \(S\) the sum
\(u=\sum_{i,j}\chi_{ij}S_{ij}\) lies in
\(\mathfrak D^{-1}\), and conjugating and interchanging the
indices gives \(\bar u=u\).  The original host element
\(w_0\in R\) has \(w_0+\bar w_0=1\), by
\(2w_0=1+\mathrm{rt}\).  Therefore
\[
 \tfrac12\operatorname{Tr}_{E/\mathbb Q}(u)
       =\operatorname{Tr}_{E/\mathbb Q}(w_0u)\in\mathbb Z.
 \tag{PUN.27}
\]
This proves integrality of (PUN.26).  The trace-one element,
not an unjustified division by \(2\) in the integer ring,
is what supplies the required even trace.
\end{proof}
% END PUN18

% BEGIN PUN19
\paragraph{PUN19. The whole original exponent match when \(l.A=1\).}
\label{claim:PUN19}
For every supplied family \(F_u\), every \texttt{Models\allowbreak Of}
satisfying the original \texttt{Model\allowbreak Base}, and every
\(\rho\in\mathbb N\),
\[
                  l.A=1\quad\Longrightarrow\quad
                         \texttt{FJ.\allowbreak Expo\allowbreak Match}\ M.X\ 1\ \rho\ l.M.
 \tag{PUN.28}
\]
Thus the asserted Fourier exponents use the original coefficient
lattice \(\texttt{Xc\allowbreak FD}\ 1\ \mathrm{cu}\), every original
cusp basis, and the exact modulus \(l.M\).

\begin{proof}
In the adapted basis of PUN16, the exact definitions of
\(\texttt{crd\allowbreak I}\), \(\texttt{dual\allowbreak ZE}\) and
\(\texttt{Xc\allowbreak FD}\) give
\[
 \begin{split}
 \texttt{crd\allowbreak I}(1,\mathrm{cu})&=\prod_i\mathfrak a_i,\\
 \texttt{Xc\allowbreak FD}(1,\mathrm{cu})&=
                      \prod_i\mathfrak D^{-1}\mathfrak a_i^{-1},\\
 \texttt{dual\allowbreak OE}(\texttt{Xc\allowbreak FD}(1,\mathrm{cu}))&=
                      \prod_i\mathfrak D\mathfrak a_i.
 \end{split}                                                   \tag{PUN.29}
\]
The first equality follows by expanding a vector in the
pseudobasis; the second is the trace-dual calculation of PCL2;
the third is ordinary fractional-ideal duality.  The trace dual
is already an \(R\)-module, so the additional span in
\texttt{Xc\allowbreak FD} does not enlarge it.

For a nonzero fractional ideal \(J\),
\(\operatorname{Tr}(zJ)\subseteq\mathbb Z\) is equivalent
to \(zJ\subseteq\mathfrak D^{-1}\).  With \(J=I_{ij}\),
the right side of (PUN.25) is therefore equivalent to
\[
     M\chi_{ij}(\mathfrak D\mathfrak a_i)
                     (\mathfrak D\bar{\mathfrak a}_j)\subseteq R
                            \qquad\text{for all }i,j.
 \tag{PUN.30}
\]
Testing vectors supported in one coordinate and then adding
shows that this is precisely
\[
    M\sum_{i,j}u_i\chi_{ij}\bar v_j\in R
       \quad\text{for all }u,v\in
                     \texttt{dual\allowbreak OE}(\texttt{Xc\allowbreak FD}(1,\mathrm{cu})).
 \tag{PUN.31}
\]
PUN18 proves equivalence with integral phase on the actual
\(\texttt{z\allowbreak Gam}\), which is the required statement in this
basis.

To verify the original arbitrary-basis quantifier, write the
given cusp basis as \(e_a^0=\sum_b P_{ab}e_b\), with
\(P\in\operatorname{GL}_\rho(E)\).  Direct substitution in
(PUN.13) gives
\[
 Z_{e^0}(\beta)=Z_e(P^t\beta\bar P).
 \tag{PUN.32}
\]
The corresponding changes of the Hermitian argument and dual
coordinates are \(\chi^0=P\chi\bar P^t\) and
\(u=u^0P\), \(v=v^0P\).  The identities
\[
 \operatorname{tr}((\chi^0)^t\beta)
       =\operatorname{tr}(\chi^tP^t\beta\bar P),
 \qquad
 u^0\chi^0\bar v^{0t}=(u^0P)\chi\overline{(v^0P)}^{\,t}
 \tag{PUN.33}
\]
preserve the phase and the integral expression.  The coordinate
pairing shows that \(\texttt{dual\allowbreak OE}(\texttt{Xc\allowbreak FD})\) transforms
by exactly this row-vector map; using \(P^{-1}\) gives both
directions.  Hence (PUN.31) and the phase condition transport
back to the given cusp, not to a substitute basis.  The empty
rank case has zero sums and is included.  Applying this to
each component proves the literal all-\(\rho\) endpoint.
\end{proof}
% END PUN19

\subsection{Actual moduli points from integral periods}

\paragraph{Classical complex-geometric foundations used in this subsection.}
We use the Riemann algebraicity criterion for a complex torus,
Appell--Humbert classification and its positivity criterion,
GAGA for proper complex schemes, and the representing dual abelian variety with its
Poincare bundle.  The universal property of the dual is used for
\emph{all} test schemes.  These are named external foundations,
not numbered reconstructions of their classical proofs.
The analytic references are Conrad, \emph{Abelian Varieties:
Geometry, Parameter Spaces, and Arithmetic}, Theorems 2.2.3,
3.1.3 and 4.3.2, and Example 4.1.4; the all-scheme universal
property is recorded in Milne, \emph{Abelian Varieties},
\S8, especially Remark 8.7 and the paragraph preceding
Remark 8.15.\footnote{Author texts:
\href{https://math.stanford.edu/~conrad/papers/cmchina.pdf}{Conrad's notes}
and \href{https://www.jmilne.org/math/CourseNotes/AV.pdf}{Milne's notes}.
The metadata lists each foundation separately.  None of these
references is being cited for the original \texttt{N1c} record.}
The explicit CM line, commutator sign, integral markings and
literal moduli installation are proved next.

% BEGIN PUN20
\paragraph{PUN20. The canonical primitive CM line on a period torus.}
\label{claim:PUN20}
Let \(P:L\to\mathbb C^g\) be an original \texttt{Is\allowbreak Period}
for \(D\), and let \(H_m\) be its positive Hermitian form.
The quotient \(T=\mathbb C^g/P(L)\) has an algebraic abelian
variety model with the given \(R\)-action and a rigidified,
symmetric, cubical, ample line bundle \(\mathcal L\) such that
the exact \texttt{SAB.\allowbreak adj}, \texttt{SAB.\allowbreak prim} and
\texttt{Out.\allowbreak KTriv} fields hold.  The exponential map is the
original quotient map, with its actual kernel \(P(L)\), and
its tangent map is an isomorphism in both directions in the
source's dual-number sense.

\begin{proof}
For \(\lambda=P(a)\) put
\[
       \alpha(\lambda)=\exp\bigl(\pi i B(a,w_0a)\bigr)\in\{1,-1\}.
 \tag{PUN.34}
\]
The action of \(w_0\) is well defined on the lattice.  Using
\(w_0+\bar w_0=1\) and the adjoint identity shows
\[
 B(a+b,w_0(a+b))-B(a,w_0a)-B(b,w_0b)
                                  \equiv B(a,b)\pmod2.
\]
Thus \(\alpha\) is the semicharacter for the alternating
form \(\operatorname{Im}H_m\).  The factors
\[
 j(\lambda,z)=\alpha(\lambda)
       \exp\bigl(\pi H_m(z,\lambda)+\pi H_m(\lambda,\lambda)/2\bigr)
 \tag{PUN.35}
\]
satisfy the lattice cocycle law, as direct multiplication
shows.  They define an analytic line bundle on \(T\).
The named Riemann and Appell--Humbert theorems make \(T\)
projective and this line ample; GAGA supplies the algebraic
variety, bundle, group law and all the \(R\)-endomorphisms.
The quotient's canonical trivial frame at \(0\) fixes its
rigidification.  Since \(\alpha(-\lambda)=\alpha(\lambda)\),
the bundle is symmetric.

Here are the original additional line identities.  Write
\(\mathcal B(\mathcal M)=m^*\mathcal M\otimes
p_1^*\mathcal M^{-1}\otimes p_2^*\mathcal M^{-1}\otimes
e^*\mathcal M\) for the normalized difference bundle.
For an Appell--Humbert bundle with form \(H_m\), the pullback
\((1,w_0)^*\mathcal B(\mathcal M)\) has form
\[
 H_m(u,w_0v)+H_m(w_0u,v)=H_m(u,v)
\]
and semicharacter exactly (PUN.34), independently of the
semicharacter of \(\mathcal M\).  Indeed the semicharacter
ratio is
\(\alpha_{\mathcal M}(a+w_0a)\alpha_{\mathcal M}(a)^{-1}
\alpha_{\mathcal M}(w_0a)^{-1}=\exp(\pi iB(a,w_0a))\).
Taking \(\mathcal M=\mathcal L\) proves the literal primitive
identity.  The adjoint relation for \(B\), and hence for
\(H_m\), proves
\((a,1)^*\mathcal B(\mathcal L)
\simeq(1,\bar a)^*\mathcal B(\mathcal L)\) for every
\(a\in R\): their mixed Hermitian forms and their mixed
semicharacters agree.  This is precisely \texttt{SAB.\allowbreak adj}.
The third finite difference of these quadratic factors cancels.
It gives the normalized cubical trivialization; its coherence
identities hold because the ratios are global holomorphic units
on powers of the compact connected torus and are one along the
normalizing axes.  GAGA transfers these identities, not only
their Chern classes, to algebraic line-bundle isomorphisms.

Perfection of \(B\) identifies its integral lattice with its
dual.  The polarization map \(T\to T^\vee\) associated to
\(\mathcal L\) is therefore an analytic isomorphism; GAGA
makes it an algebraic isomorphism.  For an arbitrary commutative
\(\mathbb C\)-algebra \(A_0\), a point \(x\) whose
translation-difference bundle is locally trivial over
\(\operatorname{Spec}A_0\) maps to zero in the representing
dual abelian variety.  The isomorphism just proved forces
\(x=0\).  This is the full \texttt{KTriv} test, including
nonreduced algebras, rather than only a statement about complex
points.  Algebraic ampleness supplies \texttt{Rel\allowbreak Ample} in its
printed form: choose a very ample power, use its projective
coordinate sections, and on each nonzero locus their ratios
generate the affine coordinate ring.

Finally the quotient map is holomorphic and locally biholomorphic.
Algebraic regular functions are analytic on their open domains,
giving the exact \texttt{Hol\allowbreak Unif.\allowbreak hol} condition.  Analytic
coordinates at the origin identify its derivative with an
isomorphism of complex tangent spaces.  A tangent vector of a
smooth complex scheme is the derivation, or equivalently the
map from the dual numbers, evaluating every local function at
its directional derivative.  This is exactly
\texttt{Is\allowbreak Diff}/\texttt{Is\allowbreak Deriv\allowbreak At0}; both required uniqueness
clauses follow.  Reading the pullback frame in any algebraic
bundle atlas gives the exact \texttt{Hol\allowbreak Unif.\allowbreak fr} and
\texttt{Is\allowbreak AH} identities (PUN.35), normalized at the origin.
\end{proof}
% END PUN20

% BEGIN PUN21
\paragraph{PUN21. The original negative Weil-pairing sign.}
\label{claim:PUN21}
For the algebraic object of PUN20, every \(n>0\) and
\(a,b\in L\), the original \texttt{Out.\allowbreak Weil\allowbreak Is} relation on
\(e(Pa/n),e(Pb/n)\) has scalar
\[
                   \exp\bigl(-2\pi i B(a,b)/n\bigr).
 \tag{PUN.36}
\]
It remains this scalar after every field embedding
\(\mathbb C\to k\).

\begin{proof}
For \(x=Pa/n\), a generator of the source's
\(\mathcal L_z^{\otimes n}\otimes
\mathcal L_{z+x}^{\otimes(-n)}\) is represented on the
covering space by
\[
                         a_x(z)=\exp(-\pi n H_m(z,x)).
 \tag{PUN.37}
\]
Its lattice transformation is the correct quotient of the
factors (PUN.35): the apparent difference between
\(H_m(\lambda,x)\) and \(H_m(x,\lambda)\) has exponential
\(\exp(2\pi iB(\lambda,Pa))=1\).
It descends to a nowhere-zero analytic section, which is
algebraic by GAGA.  In the exact multiplication order of
\texttt{Comm\allowbreak Is\allowbreak Of}, the commutator is
\[
 \frac{a_x(z)a_y(z+x)}{a_y(z)a_x(z+y)}
      =\exp\bigl(-\pi n(H_m(x,y)-H_m(y,x))\bigr)
      =\exp\bigl(-2\pi i B(a,b)/n\bigr).
 \tag{PUN.38}
\]
Every other pair of generators differs by nonzero constants,
because the abelian variety is proper and connected; these
constants cancel.  Thus the definition's universal condition on
all generators holds, not merely a computation for a chosen
pair.  The exhibited isomorphisms and their scalar equality
are algebraic.  Pulling them back along \(\mathbb C\to k\)
preserves the identities and the scalar is its image in \(k\).
Proper geometrically connected base change gives the same
constant-ratio assertion over \(k\).  This verifies the literal
local-trivialization definition of \texttt{Weil\allowbreak Is}.
\end{proof}
% END PUN21

% BEGIN PUN22
\paragraph{PUN22. Every original marked period gives an actual moduli point.}
\label{claim:PUN22}
The whole original \texttt{exists\_\allowbreak point\_\allowbreak of\_\allowbreak is\allowbreak Period} is
supplied by PUN20--PUN21 and the original \texttt{Moduli.\allowbreak rep}.
In particular retain all its data: \(D',t,P\), the compatible
maps \(f_n:D.L\to D'.L\), multipliers \(\nu_n\), the fixed
\(\sigma:K_A(l)\to\mathbb C\), the hypotheses
\(D'.p=D.p\), \(D'.q=D.q\), \(p+q\ge1\),
\(\sigma(\rho)=\mathrm{rt}_0\), the given ordered
\(\texttt{Is\allowbreak Type\allowbreak Fn}\), the given \(\texttt{Is\allowbreak Period}\), and
\[
 \begin{gathered}
 f_{nm}\equiv f_n\pmod n,\qquad
 f_n:D.L/nD.L\simeq_R D'.L/nD'.L,\\
 B'(f_na,f_nb)\equiv\nu_nB(a,b)\pmod n,\qquad
 \sigma(\zeta_A)=\exp(-2\pi i\nu_A/A),\qquad t=f_N.
 \end{gathered}                                                   \tag{PUN.39}
\]
There is a point \(y\) of the supplied \(X_o\) over
\(\sigma\), with the source's actual exponential map,
kernel, tangent bijections, \(R\)-action, all Weil pairings,
and both printed level identities.

\begin{proof}
Construct the abelian variety and its primitive line from
\(D',P\) by PUN20.  For \(n>0\) define
\[
 v_n(a)=e(P(f_na)/n),\qquad
                         \zeta_n=\exp(-2\pi i\nu_n/n).
 \tag{PUN.40}
\]
The induced map on \(L/nL\) is an isomorphism, and the
quotient description identifies all \(n\)-torsion points with
\(\tfrac1nP(L')/P(L')\).  Hence \(v_n\) has the exact
kernel \(nL\) and the exact torsion image.  Addition and
\(R\)-linearity are immediate.  The first congruence in
(PUN.39) gives the tower relation.

The multiplier \(\nu_n\) is a unit modulo \(n\): the two
perfect alternating forms modulo \(n\), transported by the
isomorphism \(f_n\), differ by this multiplier.  Since the
rank is positive, an integral symplectic pair with pairing one
tests its invertibility.  The same pair and the congruences
in (PUN.39) give \(\nu_{nm}\equiv\nu_n\pmod n\).
Thus \(\zeta_n\) is primitive and
\(\zeta_{nm}^m=\zeta_n\).  PUN21 gives
\(e_n(v_n(a),v_n(b))=\zeta_n^{B(a,b)}\).
For \(n=0\), set \(v_0=1\), \(\zeta_0=1\); no positive-level
clause is imposed there.  These formulas construct every
field of the original \texttt{Tate\allowbreak Iso}.

The same construction is valid at every algebraically closed
field \(k\) and every \(\beta:\mathbb C\to k\), with no
continuity condition on \(\beta\).  Indeed \([n]\) is an
algebraic finite etale isogeny in characteristic zero: on the
complex torus it is a finite covering of degree \(n^{2g}\),
and GAGA gives finiteness and the invertible differential.
Its kernel is therefore the constant finite etale scheme on
the complex torsion points and remains so after base change.
All torsion points over \(k\) are exactly their scalar
extensions.  The eigenbasis of the complex Lie algebra also
base changes; every embedding of the number field into \(k\)
is an extension through \(\beta\) of one of its complex
embeddings.  It therefore gives the precise \texttt{Det\allowbreak Cond}
for every \(k\)-valued point in \texttt{Obj.\allowbreak geom}.

Choose \(i\in\mathfrak b\), \(j\in\bar{\mathfrak b}\)
with \(i+j=1\).  Define the half-level marking by
\(\eta(a)=v_M(ja)\), and define \(\eta_A(a)=v_A(a)\).
The first kills \(\mathfrak bL\), since
\(j\mathfrak b\subseteq MR\), and it equals \(v_M(a)\)
on \(\bar{\mathfrak b}L\), since
\(i\bar{\mathfrak b}\subseteq MR\).  The second kills
\(AL\).  Along with PUN20, these give every field of the
original \texttt{Obj}, including its full geometric test
over every \(\beta\).  Its base values are
\(\sigma(\rho)\) and \(\sigma(\zeta_A)\).

Apply the given representing property to this object over
the reduced scheme \(\operatorname{Spec}\mathbb C\).
It gives the actual \(y:\operatorname{Spec}\mathbb C\to X_o\)
and an isomorphism with the pullback of the universal object.
Transport the quotient exponential map through this
isomorphism.  Its group, holomorphic, tangent, kernel and
Weil fields are unchanged.  Finally, for \(n=M,A\), the
congruence \(f_N\equiv f_n\pmod n\) gives
\(e(P(t a)/n)=v_n(a)\).  On the prescribed half-lattice
this is exactly \(\eta(a)\), and for all \(a\) it is
\(\eta_A(a)\) at level \(A\).  This establishes both
literal level identities at the specified point.
\end{proof}
% END PUN22

% BEGIN PUN23
\paragraph{PUN23. Equality of actual period points is exactly marked integral isometry.}
\label{claim:PUN23}
Take two sets of original period-point data
\((D',t,P,y,e)\) and \((D'',t',P',y',e')\) satisfying the
full conclusions used in \texttt{point\_\allowbreak eq\_\allowbreak iff\_\allowbreak isometry\_\allowbreak closed}.
Then \(y=y'\) if and only if there are an integral
\(R\)-linear equivalence \(\gamma:D'.L\simeq D''.L\)
and a complex-linear equivalence \(h:\mathbb C^g\simeq\mathbb C^g\)
such that
\[
 \begin{gathered}
 B''(\gamma a,\gamma b)=B'(a,b),\qquad P'\gamma=hP,\\
 \gamma t(a)-t'(a)\in MD''.L\quad(a\in\bar{\mathfrak b}D.L),\\
 \gamma t(a)-t'(a)\in AD''.L\quad(a\in D.L).
 \end{gathered}                                                   \tag{PUN.41}
\]

\begin{proof}
If \(y=y'\), the two exponential maps are holomorphic covering
homomorphisms of the same complex torus, with invertible
derivatives.  Lifting from the simply connected space
\(\mathbb C^g\), or using the universal cover of a complex
Lie group, gives a unique additive holomorphic equivalence
\(h\) with \(e'h=e\).  An additive holomorphic map is
complex linear.  The kernel identities give
\(hP(L')=P'(L'')\), hence the integral equivalence
\(\gamma\).  Compatibility with the \(R\)-action implies
its \(R\)-linearity.  Comparing the Weil scalars for every
positive \(n\) shows
\(B''(\gamma a,\gamma b)-B'(a,b)\in n\mathbb Z\) for every
\(n\); it is therefore zero.  Comparing the two specified
level markings gives exactly the last two congruences of
(PUN.41).

Conversely (PUN.41) gives an analytic isomorphism of complex
tori, which is algebraic by GAGA and respects the \(R\)-action.
The given Weil-pairing identities for all \(n\) identify each
fiber line's integral Riemann form with \(B'\) or \(B''\),
by the same divisibility argument.  The primitive condition
on an original \texttt{SAB} then determines its semicharacter:
the calculation in PUN20 says that
\((1,w_0)^*\mathcal B(\mathcal M)\) has precisely (PUN.34).
Thus the isomorphism preserves the actual primitive line,
not just its polarization class.  Normalize the line
isomorphism at zero; the cubical isomorphism follows by
uniqueness of the normalized cube trivialization.  The two
congruences in (PUN.41) preserve \(\eta\) and \(\eta_A\).
We have an isomorphism of the original moduli objects.
Uniqueness in the supplied \texttt{Moduli.\allowbreak rep}, applied to
\(\operatorname{Spec}\mathbb C\), forces \(y=y'\).
This proves both directions of the literal endpoint without
restricting to a chosen representative family.
\end{proof}
% END PUN23

% BEGIN PUN24
\paragraph{PUN24. A marked genus exists for every prescribed base embedding.}
\label{claim:PUN24}
For every original \(\sigma:K_A(l)\to\mathbb C\) satisfying
\(\sigma(\rho)=\mathrm{rt}_0\), the original
\texttt{In\allowbreak Genus}\((D,l,\sigma,D',t)\) is inhabited.
One can take \(D'=D\).  Together with PUN15 and PUN22, this
constructs a point of the original moduli space over that
\(\sigma\) whenever \(p+q\ge1\).

\begin{proof}
Write \(\sigma(\zeta_A)=\exp(-2\pi i k/A)\), with
\(k\) a unit modulo \(A\).  Choose
\(u\in\widehat{\mathbb Z}^{\times}\) reducing to \(k\)
modulo \(A\): at primes dividing \(A\) use the indicated
unit residue, and choose one at the other primes.
At every finite place of \(F\), the original unramified
quadratic integer algebra has surjective norm on units by
MF2, including at \(2\).  Therefore choose
\(s\in(\mathcal O_E\otimes\widehat{\mathbb Z})^\times\)
with \(s\bar s=u\).  The choices are independent at the
finite places and combine in this profinite product.

For each \(n>0\), choose \(s_n\in R\) lifting the residue
of \(s\) in \(R/nR\), and choose an integer \(\nu_n\)
lifting \(u\bmod n\).  Put \(f_n(a)=s_na\).
Reduction is invertible, \(f_{nm}\equiv f_n\pmod n\), and
the adjoint identity gives
\[
 B(f_na,f_nb)=B(a,\bar s_ns_nb)
                              \equiv\nu_nB(a,b)\pmod n.
\]
Take \(t=f_N\).  The required root-of-unity identity is the
chosen one at \(A\).  These are all fields of the original
\texttt{In\allowbreak Genus}; the omitted \(n=0\) values may be chosen
arbitrarily.  If \(A=1\) the congruence imposes no restriction
and choose \(u=1\).  The construction neither changes the
coefficient embedding nor assumes a global integral norm
solution.  PUN15 supplies a period of \(D\) for each desired
type, and PUN22 gives its actual moduli point.
\end{proof}
% END PUN24

\subsection{The exact remaining principal-chart boundary}

The arithmetic outputs above are complete at their printed
scope.  PUN1 supplies the literal quadratic norm theorem;
PUN10 supplies the whole original lattice model on the same
Hermitian space.  For any supplied original
\(\texttt{Family\allowbreak Unif}\), PUN11 constructs its flag-domain
models with the prescribed \(G_m\), \(\xi_0\), lattices,
principal groups, and trivial character.  PUN14 supplies
\(\texttt{Lat\allowbreak Z}\) for every family of original data and every
cusp rank.  PUN19 supplies the original exponent match,
on the original coefficient lattice and phase, for every
\(\rho\) when \(l.A=1\).  These results do not require a
uniformization to be silently substituted for the one in the
source.

PUN20--PUN24 additionally construct actual primitive moduli
objects from every marked period, their original Weil signs
and level markings, equality precisely by marked integral
isometry, and nonemptiness at every prescribed base embedding.
They use the supplied original \(\texttt{Moduli}\) record;
existence of that representing scheme is not asserted here.
They do not yet assemble the finite family required by
\(\texttt{N1c}\).  The still-needed original producers are
\texttt{exists\_\allowbreak uniformisation\_\allowbreak at\_\allowbreak point\_\allowbreak closed},
\texttt{exists\_\allowbreak local\_\allowbreak period\_\allowbreak lift}, the holomorphy of the
moduli point and exponential for arbitrary holomorphic period
families, and the simultaneous theta-frame family with an
algebraic section on every fiber.  Their finite-component
assembly must preserve the same base embedding, type, level
markings, and chosen period at the base point.

After that family is constructed, the remaining original
\texttt{Frames\allowbreak OK}, \texttt{Cover\allowbreak Tri}, and \texttt{Fib\allowbreak Tri}
producers must still establish one simultaneous system of
normalized framed points: the precise Gaussian norm,
triangular weight character, arithmetic-group invariance,
holomorphic readings, cover, and fiber comparison.  The
resulting \(\texttt{Is\allowbreak Unf}\) must use those same points,
not merely an isomorphic unframed quotient.  Finally the
\texttt{From\allowbreak Test} descent must use that same family.  Only
then may the original
\texttt{Frame\allowbreak Unif.\allowbreak exists\_\allowbreak unf\_\allowbreak has\allowbreak Princ} be assembled from
these arithmetic suppliers and the analytic constructors.

Accordingly this packet supplies substantial original
subconstructors while leaving the complete principal-chart
existence statement open.  It does not close the common-field,
common-metric, common-measure, and common-phase obligations,
the full unconditional integral cusp spanning statement,
higher covers, or the all-tuple class-sum producers.


\section{One actual compatibility stage for all later choices}
\label{sec:common-compatibility-v48}

\status{GCF1--GCF18 construct the original common local comparison stage
on the supplied true-column data.  The output is one \texttt{GJStage},
chosen before every later \(W,V,S,N\), with the full original
\texttt{GIs2}: true volumes, the original numerical local scales, and
\texttt{Column\allowbreak Data\allowbreak Split\allowbreak Primes}.  In particular all five fields of
\texttt{Group\allowbreak Compatibility\allowbreak All} are constructed together on these
stages.  The proof also preserves the actual comparison map through
arbitrary division-lattice changes.  It retains the original arithmetic
host and column input, and does not manufacture an eventual realization
of an arbitrary already supplied \texttt{Fam\allowbreak P} by these stages.
The selected design class already supplies its top estimate; its
remaining family-level issue is the simultaneous realization and its
equality with the given \texttt{Extra.\allowbreak cmp}.}

\paragraph{The precise objects being constructed.}
The accepted source is the pin recorded in this component's provenance.
Write \(E/F\) for the original CM extension and \(a\mapsto\bar a\)
for its fixed involution.  Its finite relative primes are unramified,
as required by the original host.  Let \(\mathbb I_E\) be its idele
group, \(U_E\) the subgroup of finite-component units (with arbitrary
infinite components), and
\[
 U^1(\mathbb A_{F,f})=
 \{z\in\mathbb A_{E,f}^{\times}:z\bar z=1\},\qquad
 U^1(\widehat{\mathcal O}_F)=U^1(\mathbb A_{F,f})
                                  \cap\widehat{\mathcal O}_E^{\times}.
 \tag{GCF.1}
\]
These are exactly the original rank-one unitary group and its integral
subgroup under the finite-adelic base-change identification; in rank
one a matrix is identified with its determinant.  For the true-column
datum \(I\), put \(S=S_G\cup S_Q\), the source's \texttt{u\allowbreak SG}.
Its \texttt{Lev\allowbreak Ok} says that \(S\cap\bar S=\varnothing\), and that
\(S\cup\bar S\) includes every auxiliary-level prime.  Conversely every prime of
\(S_G=(B.\mathrm{primes}(X)).\mathrm{vGs}\) lies above an
auxiliary level.  Thus \(S\cup\bar S\) is exactly the finite
set where the defining quotient imposes its congruences.
The actual finer determinant group is
\[
 D_{T,Q}=\mathbb I_E/K_{D,Q},\qquad
 U_{O,Q}=\operatorname{im}(U_E\longrightarrow D_{T,Q}).
 \tag{GCF.2}
\]
Here \(K_{D,Q}\) is generated by finite unit ideles congruent to one
at \(S\cup\bar S\), and by the original local \(F^{\times}\)
ideles.  This is the same quotient as HF14--HF15, including its
principal-unit congruences; no quotient by all units is taken until a
relative determinant is formed.

% BEGIN GCF1
\paragraph{GCF1. Hilbert--90 lifts in the actual finite adeles.}
\label{claim:GCF1}
Every \(z\in U^1(\mathbb A_{F,f})\) has an idele lift \(a\in
\mathbb I_E\) satisfying
\[
                         z=a_f/\bar a_f.             \tag{GCF.3}
\]
Two such lifts differ, in their finite components, by an idele of
\(F\).  Every such fixed finite idele, with arbitrary infinite
components, belongs to \(K_{D,Q}\).

At a split place, write \(z=(t,t^{-1})\) and take the lift
\((t,1)\).  At a nonsplit place the elementary field argument applies:
\(a=1+z\) works unless \(z=-1\), in which case choose a nonzero
trace-zero element.  These local choices are units almost everywhere.
For completeness, at any unramified nonsplit place where \(z\) is a
unit, reduce to the quadratic residue-field extension.  Some residue
element \(b\) has \(b+z\bar b\ne0\).  Otherwise putting \(b=1\)
would give \(z=-1\), and the asserted identity for all \(b\) would
make the nontrivial residue involution the identity.  Lift this \(b\)
integrally and set \(a=b+z\bar b\).  It is a unit and satisfies
\(a=z\bar a\).  The split construction is also a unit at every
place where \(z\) is integral.  Thus the local lifts and their inverses
are restricted products, and give a finite idele.  Assigning arbitrary
nonzero infinite components gives (GCF.3).  Dividing two solutions
shows directly that their quotient is fixed by conjugation.

To verify the last assertion with the source's \emph{rational} local
generators, let \(b\) be fixed and consider the finite set consisting
of its nonunit components and the congruence places.  At each underlying
place \(v\) of \(F\), density of \(F\) in \(F_v\) gives
\(k_v\in F^{\times}\) with the same valuation as \(b_v\), and with
\(b_v/k_v\equiv1\) whenever a congruence is required.  Dividing by
the finite product of \(\operatorname{loc}_v(k_v)\) leaves a unit
idele with the prescribed congruences.  The latter and the finite
product both belong to the defining subgroup in (GCF.2).
Infinite components cause no condition in that subgroup.  This also
proves independence of the global lift without changing the original
coefficient field or adding a choice of local representatives to the
quotient.
\hfill\(\square\)
% END GCF1

% BEGIN GCF2
\paragraph{GCF2. The exact congruence criterion for the finer quotient.}
\label{claim:GCF2}
For an idele \(b\), the following two conditions are equivalent to
\(b\in K_{D,Q}\):
\[
 \begin{aligned}
  \operatorname{ord}_P(b_P)&=
                 \operatorname{ord}_{\bar P}(b_{\bar P})
                        &&\text{for every }P,\\
  b_P/(\bar b)_P&\equiv1\pmod P &&\text{for }P\in S.
 \end{aligned}                                                   \tag{GCF.4}
\]
The first equality includes a prime fixed by conjugation.  The second
condition uses the actual conjugation between the two local fields.

Each generator of \(K_{D,Q}\) satisfies (GCF.4).  A local
\(F^{\times}\) idele is fixed; a congruence unit has order zero
and both paired components are one modulo their maximal ideals.
The conditions are multiplicative and closed under inverses, proving
necessity for the subgroup generated by these elements.

Conversely, suppose (GCF.4).  At a place below \(P\in S\), the
extension is split because \(P\ne\bar P\).  Choose
\(k_v\in F^{\times}\) sufficiently close to the \(\bar P\)
component of \(b\).  Its quotient with that component is a principal
unit, and the second equation in (GCF.4) makes the quotient at \(P\)
a principal unit as well.  At a place with no congruence requirement,
choose a local \(F\)-uniformizer of the common order in (GCF.4).
It has that same order upstairs because the original extension is
unramified.  A rational element of \(F\) with this local order can
be used, by density.  Only finitely many places require these choices:
the congruence set is finite, and an idele has finitely many nonunit
components.  Dividing \(b\) by their finite local product leaves an
element of the defining congruence-unit subgroup.  This proves
sufficiency.

In particular, forgetting the congruences gives the original identity
\[
 K_{D,Q}U_E=K_DU_E
 =\{b:\operatorname{ord}_P b_P
                  =\operatorname{ord}_{\bar P}b_{\bar P}
                    \text{ for every }P\}.             \tag{GCF.5}
\]
This agrees with HF15 and will preserve its positive-part relative
determinant representative.
\hfill\(\square\)
% END GCF2

% BEGIN GCF3
\paragraph{GCF3. The canonical norm-one determinant map, with all scalar laws.}
\label{claim:GCF3}
Define
\[
 d_Q:U^1(\mathbb A_{F,f})\longrightarrow D_{T,Q},\qquad
                d_Q(z)=[a]\quad(a_f/\bar a_f=z).        \tag{GCF.6}
\]
GCF1 makes this well defined, and products of lifts show it is a
homomorphism.  It is surjective: for any representative \(a\), use
\(z=a_f/\bar a_f\).  Under the rank-one matrix dictionary, (GCF.6)
is the original \texttt{u1Dt}; it is characterized by
\(d_Q(z)=[a]\) for \emph{every} \texttt{Is\allowbreak Det\allowbreak Lift}, not merely for a
selected lift.

For a lift \(a\),
\[
 \operatorname{ord}_P z_P=
       \operatorname{ord}_P a_P-
                  \operatorname{ord}_{\bar P}a_{\bar P}.
 \tag{GCF.7}
\]
Thus \(z\) is integral if and only if the orders of \(a\) are
balanced: integrality of a norm-one element includes integrality of
its inverse, since the inverse is its conjugate.  By GCF2,
\[
 \ker d_Q=
 \{z\in U^1(\widehat{\mathcal O}_F):z_P\equiv1\pmod P
                                                    \ (P\in S)\}.
 \tag{GCF.8}
\]
This is open in the actual finite-adelic topology.  Indeed the integral
unit subgroup is open in the restricted product of local unit groups,
and finitely many principal-unit conditions define an open subgroup.
The identification \(E\otimes_F\mathbb A_{F,f}\simeq
\mathbb A_{E,f}\) has this topology: decompose a local tensor product
by the primes above \(v\); its integral subring is the product of the
integral closures, and then take the restricted product.  Consequently
(GCF.8) supplies the source's \texttt{dt\_\allowbreak ker} and \texttt{dt\_\allowbreak open}
on its specified tensor-product model.

The image of the integral subgroup is exactly \(U_{O,Q}\).
One inclusion follows by taking a unit idele lift.  For the reverse
inclusion, any integral \(z\) has a lift with balanced orders by
(GCF.7); (GCF.5) writes its class as the class of a unit idele.
Finally, for rational \(z\in E^{\times}\) of norm one, choose
\(k\in E^{\times}\) with \(z=k/\bar k\).  The explicit solutions
in HF14 give the original rational determinant map
\(e_{1,Q}(z)=[\operatorname{diag}k]\).  The diagonal lift in
(GCF.6) therefore proves
\[
 d_Q(z_f)=e_{1,Q}(z),\qquad
 d_Q\bigl(U^1(\widehat{\mathcal O}_F)\bigr)=U_{O,Q}.    \tag{GCF.9}
\]
These are all six scalar determinant fields: the map, surjectivity,
kernel inclusion, open kernel, integral-unit image, and rational-point
law.  The last law also covers any rank-one unitary matrix whose
determinant is the specified rational element.
\hfill\(\square\)
% END GCF3

% BEGIN GCF4
\paragraph{GCF4. Reading the determinant of an actual decorated point.}
\label{claim:GCF4}
Let \(p\) be an original framed decorated self-dual lattice of rank
\(n\), let \(g\) realize its extension to the prescribed \(S_Q\)
frames, and let \(a\) be a Hilbert--90 lift of \(\det g\).  Then
\[
                         [a]=\operatorname{detPt}(p)
                                       \quad\text{in }D_{T,Q}.
 \tag{GCF.10}
\]
Here every frame is the one in the given decoration.  Choose its
rational local basis matrix \(F_P\), and put
\(d_P=\operatorname{ord}_P\det F_P\).  The local duality argument
in HF15 gives \(d_{\bar P}=-d_P\).  The source's point idele
\(a_p=\operatorname{ptIdl}(p)\) has components
\[
 (a_p)_P=
 \begin{cases}
    \det F_P,&P\in S,\\
    1,&\bar P\in S,\\
    \pi_P^{\max(d_P,0)},&P,\bar P\notin S.
 \end{cases}                                                   \tag{GCF.11}
\]
The set \(S\) contains no conjugate pair, so these cases are
unambiguous.  They give
\(\operatorname{ord}_P(a_p)_P-
\operatorname{ord}_{\bar P}(a_p)_{\bar P}=d_P\).
The actual carried-lattice equality gives
\(\operatorname{ord}_P\det g_P=d_P\), hence \(a/a_p\) has
balanced orders by (GCF.7).

At a selected prime put \(K_P=g_P^{-1}F_P\).  Both matrices carry
the standard local lattice onto the given local lattice; therefore
\(K_P\) and its inverse are integral.  The residue-frame condition
in \texttt{Is\allowbreak Pt\allowbreak Of} gives \(K_P\equiv1\pmod P\) at an auxiliary
prime.  At a prescribed rational-frame prime it gives \(K_P=1\).
These assertions can be checked column by column after multiplying
each rational local basis vector by a denominator outside \(P\).
Reduction of the determinant consequently gives
\(\det K_P\equiv1\pmod P\).  Since
\(\det g_P\det K_P=\det F_P\), (GCF.11) implies
\[
 \frac{(a/a_p)_P}{\overline{(a/a_p)}_P}
                 =(\det K_P)^{-1}\equiv1\pmod P.
\]
GCF2 proves \(a/a_p\in K_{D,Q}\), establishing (GCF.10).

Now suppose the original \texttt{Reads} datum realizes the class
by \(y k\), where \(k\) is unitary, preserves the standard lattice,
and is the identity at all the selected underlying places.  Its
determinant is integral and is one at those places, so GCF3 gives
\(d_Q(\det k)=1\).  Therefore (GCF.10) becomes
\[
                \operatorname{detPt}(p)=d_Q(\det y).
 \tag{GCF.12}
\]
For the reference tuple, the twist-prime set is empty by its literal
definition.  The \texttt{Reads.\allowbreak k\_\allowbreak off} clause therefore supplies this
identity condition at the auxiliary and \(S_Q\) places; the type
primes are disjoint from them.  This verifies the original
\texttt{point\allowbreak Determinant\_\allowbreak eq\_\allowbreak }\allowbreak\texttt{norm\allowbreak One\allowbreak Class\_\allowbreak of\_\allowbreak level} on the actual
class representatives, including rank zero, where all determinants
are one.
\hfill\(\square\)
% END GCF4

% BEGIN GCF5
\paragraph{GCF5. All column determinant readings and the adelic relative law.}
\label{claim:GCF5}
Take the actual \(C:\mathrm{ColsTiedSym}(I.\mathrm{rec7}(l,Z))\)
and its original \(\mathrm{StdC}\).  On its underlying Hermitian
columns, use \(d_Q\) from GCF3.  The six scalar fields in that claim
and the following two readings give the full
\texttt{Symmetric\allowbreak Hermitian\allowbreak Columns} record:
\[
 \begin{aligned}
 D.\mathrm{detCls}(u,0,c)&=d_Q(\det y_{u,0,c}),\\
 D.\mathrm{detCls}(u,i,c)&=d_Q(\det y_{u,i,c})
                                  &&(1\le i\le m).
 \end{aligned}                                                   \tag{GCF.13}
\]
For step zero use the endpoint \texttt{Reads} in \(\mathrm{StdC}\)
and the \texttt{Car0} representative equality.  For
\(1\le i<i+1\le m\), use its middle \texttt{Reads} and the
corresponding chart carrier.  For \(i=m\), use its terminal
\texttt{Reads}, the singleton terminal chart, and \texttt{Car\allowbreak M}.
Each is precisely (GCF.12).  The carriers only reindex the matrix
basis by the supplied rank equality; a simultaneous row and column
reindexing leaves the determinant unchanged.  Thus (GCF.13) holds for
the original representatives and for every rank-one unitary matrix
having their determinant, as demanded by \texttt{det0} and
\texttt{det\allowbreak I}.

There is also the full relative law for \emph{every} rank \(N\) and
every pair of self-dual lattices \(M,M'\).  Suppose an actual unitary
finite adelic matrix \(g\) carries \(M\) onto \(M'\), and a
rank-one unitary element \(z\) has determinant \(\det g\).  Choose
local basis matrices for the two lattices.  In those bases the
carried-lattice condition is an invertible integral matrix, so
\[
 \operatorname{ord}_P\det g_P=d_P(M')-d_P(M)=:\Delta_P.
 \tag{GCF.14}
\]
This uses local bases only.  It does not presume either global lattice
is free.  Completion of a finite lattice and localization give the
same integral span at \(P\); this verifies that the source's
\texttt{Carries} has exactly the interpretation used in (GCF.14).
For a lift \(a\) of \(z\), (GCF.7) gives the same differences
\(\Delta_P\).  The source relative idele is the positive-part
product in HF15,
\[
              a_{M,M'}=\prod_P
                  \operatorname{loc}_P(\pi_P)^{\max(\Delta_P,0)}.
\]
Self-duality gives \(\Delta_{\bar P}=-\Delta_P\), so this product
has the same conjugate-order differences as \(a\).
Their ratio belongs to \(K_{D,Q}U_E\) by (GCF.5).  Passing to the
quotient by \(U_{O,Q}\) proves exactly
\[
       D.\operatorname{relDet}(N,M,M')=[d_Q(z)]
                                  \quad\text{in }D_{T,Q}/U_{O,Q}.
 \tag{GCF.15}
\]
This is \texttt{Symmetric\allowbreak Relative\allowbreak Determinant}, with its original
positive-part convention and no rationality restriction on \(g\).
\hfill\(\square\)
% END GCF5

% BEGIN GCF6
\paragraph{GCF6. The full Hermitian extension preserves the given local ties.}
\label{claim:GCF6}
For the original \(I,\mathrm{Ok},\mathrm{Pins},l,Z,C,\mathrm{StdC}\),
suppose true volume compatibility is supplied and fix an actual
\(g:\mathrm{SymmetricGroupCompatibility}(D,C.\mathrm{herm.base})\).
There exists the source's \(F_u:\mathrm{ColsFullS}(C)\) such that
\[
 F_u.\mathrm{g}=g,\qquad
 I.\mathrm{StdDt}(l,Z,F_u.\mathrm{base.dt}).             \tag{GCF.16}
\]
The first equality is equality after the original dependent
identifications; it preserves the supplied local ties themselves.

Extend \(C.\mathrm{herm.base}\) by the eight fields proved in
GCF3 and GCF5.  Take its \texttt{at\allowbreak T}, level-zero compatibility,
adelic compatibility, cofactors, auxiliary-generator identity,
stabilizer maps, residue-field levels, standard lattice and base
compatibility directly from the supplied \(C.\mathrm{herm}\).
Insert the given \(g\), the supplied volume identity, (GCF.15), and
the original positive-volume field \(B.V_{\rm pos}\).
This lists every field of \texttt{Full\allowbreak Tied\allowbreak S}.  Forgetting the new
determinant, volume and local-group fields returns exactly
\(C.\mathrm{herm}\), which is the defining equality in
\texttt{Cols\allowbreak Full\allowbreak S}.  The chosen group field is still \(g\).
Finally the \texttt{Std\allowbreak Dt} assertion quantifies over every idele lift
of every rank-one element; it is exactly the universal characterization
of (GCF.6).  This proves the entire original
\texttt{exists\_\allowbreak full\allowbreak Hermitian\allowbreak Columns\_\allowbreak true\allowbreak Columns} endpoint.
\hfill\(\square\)
% END GCF6

\subsection{The given period root needs no new coefficient field}

% BEGIN GCF7
\paragraph{GCF7. The finite torsion action with its exact scaling factor.}
\label{claim:GCF7}
Let \(\Lambda\) be a complex period lattice and write
\(G=g_2(\Lambda)\), \(T=g_3(\Lambda)\).  Suppose an automorphism
\(\sigma\) of \(\mathbb C\) fixes \(G^3\) and \(T^2\).
There is \(\eta\in\mathbb C\), \(\eta^6=1\), with
\[
 \sigma G=\eta^2G,\qquad \sigma T=\eta^3T.             \tag{GCF.17}
\]
For every torsion class \(t\pmod\Lambda\) there is a class
\(t'\pmod\Lambda\) such that, for every scalar \(c\) preserving
\(\Lambda\) and fixed by \(\sigma\),
\[
 \begin{split}
 ct\in\Lambda&\ \Longleftrightarrow\ ct'\in\Lambda,\\
 \sigma\wp_\Lambda(ct)&=\eta\wp_\Lambda(ct')
                                      \quad(ct\notin\Lambda).
 \end{split}                                                   \tag{GCF.18}
\]

If \(G\ne0\), put \(u=\sigma G/G\), so \(u^3=1\);
if \(G=0\), take \(u=1\).  The equality for \(T^2\) gives
\(\sigma T=T\) or \(-T\).  Accordingly choose
\(\eta=u^2\) or \(-u^2\).  These choices prove (GCF.17), including
the zero-invariant cases.  Choose \(m\) with \(m^2=\eta^{-1}\).
Then \(m\ne0\), and the scaling laws PA1 give
\[
 g_2(m\Lambda)=\sigma G,\qquad
 g_3(m\Lambda)=\sigma T,\qquad m^{12}=1.
\]
Apply coordinatewise \(\sigma\) to the elliptic curve
\(y^2=x^3-Gx/4-T/4\), and uniformize the image by
\(\mathbb C/(m\Lambda)\).  TA1 identifies precisely this
conjugate lattice, and TA5 proves that addition, torsion and CM
multiplication are transported by their algebraic formulas.  A
CM endomorphism with differential \(c\) becomes one with differential
\(\sigma c=c\).  Dividing the image point's parameter by \(m\)
therefore gives \(t'\) with the first assertion in (GCF.18).
The identity
\(\wp_{m\Lambda}(mz)=m^{-2}\wp_\Lambda(z)=\eta\wp_\Lambda(z)\)
gives the second.  This argument applies a possibly discontinuous
\(\sigma\) to algebraic curve identities, not to analytic limits or
complex Taylor series.
\hfill\(\square\)
% END GCF7

% BEGIN GCF8
\paragraph{GCF8. The distinguished two-torsion point fixes the sign of \(g_3\).}
\label{claim:GCF8}
Suppose \(G^3-27T^2=1\), and a scalar \(w\) preserves
\(\Lambda\), satisfies \(w(1-w)=2k\) for \(k\in\mathbb Z\),
and admits a class \(t\) with
\[
       2t\in\Lambda,\qquad wt\notin\Lambda,\qquad
                                      (1-w)t\notin\Lambda.
 \tag{GCF.19}
\]
Every complex field automorphism fixing \(w,G^3,T^2\) fixes
\(T\) itself.

By GCF7 it transports \(t\) to a class \(t'\) with exactly
(GCF.19).  TA11 proves \(t'=t\pmod\Lambda\): on the
two-dimensional \(\mathbb F_2\)-space of two-torsion, \(w\) and
\(1-w\) are complementary idempotents, and the class with both
components nonzero is unique.  Put
\(x=\wp_\Lambda(wt)\) and \(y=\wp_\Lambda((1-w)t)\).
Both are abscissae of nonzero two-torsion points and satisfy
\[
                 4X^3-GX-T=0.                        \tag{GCF.20}
\]
They are distinct: equality of two \(\wp\)-values means equality
up to sign modulo the lattice; since both classes are two-torsion,
this would make their sum \(t\) zero and contradict (GCF.19).
Subtracting their two equations gives
\(G=4(x^2+xy+y^2)\).  Substitution into either equation gives
\[
 \{4(x-y)(2x+y)(x+2y)\}^2=G^3-27T^2=1.               \tag{GCF.21}
\]
The factor in braces is therefore \(1\) or \(-1\), so \(\sigma\)
fixes it.  On the other hand (GCF.18) and uniqueness of \(t\)
multiply each of \(x,y\) by \(\eta\); the same factor is multiplied
by \(\eta^3\).  Thus \(\eta^3=1\).  Equation (GCF.17) proves
\(\sigma T=T\).

The original CM ideal lattices meet the existence clause in (GCF.19).
Here \(w=(1+i\sqrt d)/2\), and the host gives
\(d\equiv7\pmod8\), so \(w(1-w)=(d+1)/4\) is even.
For the unit ideal PA2 gives \(t=\delta/2\) directly.
For an arbitrary nonzero ideal \(\mathfrak a\subset\mathcal O_K\),
\(\mathcal O_K/2\simeq\mathbb F_2\times\mathbb F_2\), with
\(w\) the two complementary idempotent coordinates.  Invertibility
of \(\mathfrak a\) makes \(\mathfrak a/2\mathfrak a\) a
rank-one module over each factor.  Its unique element nonzero in
both factors, divided by two and multiplied by \(\delta\), gives
the required class on \(\delta\iota(\mathfrak a)\).
This verifies the source's half-point input without extending the
coefficient number field.
\hfill\(\square\)
% END GCF8

% BEGIN GCF9
\paragraph{GCF9. Three-torsion controls the cubic ambiguity of \(g_2\).}
\label{claim:GCF9}
Suppose \(G^3-27T^2=1\), \(\beta\Lambda\subset\Lambda\),
\(\beta^2=-d\), where \(d>0\) and \(3\nmid d\).
There is \(v\in\mathbb C\), \(v^3=1\), such that every
automorphism fixing \(\beta,G^3,T^2\) fixes \(G/v\).

First choose a three-torsion class \(t\) not killed by either
\(\beta-1\) or \(\beta+1\).  Multiplication by \(\beta\)
has an integral matrix of trace zero on a basis of the lattice:
its complex eigenvalues are \(\beta,\bar\beta=-\beta\).
On \(\Lambda/3\Lambda\simeq\mathbb F_3^2\), neither
\(\beta-1\) nor \(\beta+1\) is the zero map, since the
trace of \(1\) or \(-1\) is nonzero in \(\mathbb F_3\).
Their two kernels have at most three elements each and cannot cover
the nine-element space.  Choose a class outside their union and
divide its lattice representative by three.  This gives
\[
 3t\in\Lambda,\qquad(\beta-1)t\notin\Lambda,\qquad
                              (\beta+1)t\notin\Lambda.          \tag{GCF.22}
\]
Multiplication by \(\beta\) is invertible on three-torsion, since
\(3\nmid d\).  Its square is \(1\) or \(-1\) there, so it
acts as an involution on the four nonzero abscissae, which identify
opposite classes.  By (GCF.22),
\(x=\wp(t)\) and \(y=\wp(\beta t)\) form a two-element orbit.

Every nonzero three-torsion abscissa \(X\) satisfies
\[
 3X^4-\tfrac32GX^2-3TX-G^2/16=0.                    \tag{GCF.23}
\]
For an explicit derivation, on \(y^2=x^3-Gx/4-T/4\) the tangent
slope is \((3x^2-G/4)/(2y)\).  A nonzero three-torsion point has
\(y\ne0\) and \(2P=-P\), so its doubled abscissa is \(x\).
The duplication formula gives the slope squared equal to \(3x\);
substitution yields (GCF.23).  The four opposite pairs have distinct
abscissae by the degree-two \(\wp\) map, so they are exactly the
four roots of (GCF.23).  Their sum is zero.

For two distinct roots \(x,y\), subtracting their equations yields
\(T=(x+y)(x^2+y^2)-G(x+y)/2\).  Subtracting their equations after
multiplication by the other root also yields
\[
 -3xy\{(x+y)^2-xy\}+\tfrac32Gxy-G^2/16=0.
\]
Eliminating \(T\) with these two identities gives
\[
                   \{G-3(x+y)^2\}^3=G^3-27T^2=1.     \tag{GCF.24}
\]
Set \(v=G-3(x+y)^2\), which is therefore nonzero.

Apply GCF7 to an automorphism fixing the stated data.  It sends the
chosen orbit to another two-element orbit \(\{x',y'\}\) of the
same involution.  The two orbits either agree or are disjoint.
In the latter case they are all four roots of (GCF.23), so
\(x'+y'=-(x+y)\).  In both cases
\((x'+y')^2=(x+y)^2\).  Consequently
\[
 \sigma v=\eta^2v,\qquad
                 \sigma(G/v)=G/v,                   \tag{GCF.25}
\]
as claimed.  Repeated roots or an unspecified pairing are not used:
the opposite three-torsion classes give the four distinct roots and
the original CM scalar gives their involution.
\hfill\(\square\)
% END GCF9

% BEGIN GCF10
\paragraph{GCF10. The prescribed twelfth root has invariants in the given field.}
\label{claim:GCF10}
Take the original \(D\), \(\mathrm{FieldOK}\), and a prescribed
\(\delta\) satisfying
\(\delta^{12}=\Delta_{\iota(\mathrm{ccK})}\).
Let \(M=\operatorname{im}(D.\mathrm{embLL})\subset\mathbb C\),
and \(\Lambda=\delta\iota(\mathrm{ccK})\).
PA1 gives \(\Delta_\Lambda=1\), and the first \texttt{Field\allowbreak OK}
clause and homothety invariance of \(j\) give
\[
  G^3=j(\iota(\mathrm{ccK}))/1728\in M,\qquad
                    T^2=(G^3-1)/27\in M.             \tag{GCF.26}
\]
The image of the original CM field is in \(M\), so it contains
\(w=(1+i\sqrt d)/2\) and \(\beta=2w-1\).
The host has \(3\nmid d\).  GCF8 and GCF9 show that all
automorphisms of \(\mathbb C\) fixing \(M\) fix \(T\) and
\(G/v\), for a fixed \(v^3=1\).  Hence
\[
                              T\in M,\qquad G/v\in M.
 \tag{GCF.27}
\]
The fixed-field assertion here is the one used in TA5: an algebraic
element outside a characteristic-zero field has a different conjugate,
and an element transcendental over it can be moved by extending a
transcendence basis.  Either embedding extends to an automorphism of
\(\mathbb C\).  Thus an element fixed by all these automorphisms is
in \(M\).

For arbitrary \(M\) this already proves the original
\texttt{exists\_\allowbreak delta\_\allowbreak data}: put \(u=v^{-1}\) and
\(\delta'=u^2\delta\).  Then \((\delta')^{12}=\delta^{12}\),
\[
             g_2(u^2\Lambda)=uG=G/v\in M,\qquad
                         g_3(u^2\Lambda)=T\in M.       \tag{GCF.28}
\]
For the actual \(\mathrm{Base}\), no period is changed: its
coefficient field contains all third roots of unity, since its
prescribed roots field contains \(\mathbb Q(\mu_{3^4})\).
Thus \(v\in M\), and (GCF.27) gives \(G,T\in M\) for the
\emph{given} \(\delta\).  Equivalently one may compare it with
\(\delta'\): their ratio has twelfth power one, so its fourth
power is a third root of unity and its sixth power is \(\pm1\).
The scaling factors for \(g_2,g_3\) already lie in \(M\).

All remaining \texttt{Data\allowbreak In} coordinates are
\(D.\mathrm{embLL}(D.\mathrm{embE}(\sigma_\tau a))\), and
therefore also lie in \(M\).  This supplies the complete original
\texttt{Data\allowbreak In} for \(\mathrm{tor0}(D,\mathrm{ccK},\delta)\)
without enlarging \(D.\mathrm{LL}\), changing its embedding, or
renormalizing the originally chosen period root.
\hfill\(\square\)
% END GCF10

% BEGIN GCF11
\paragraph{GCF11. The original \(\Omega\) formula holds for the true global stage.}
\label{claim:GCF11}
Use GCF10 on the original base with its third roots of unity, and
write \(e:D.\mathrm{LL}\simeq M\) for its image-field isomorphism.
For every original infinity label \(\tau\), let \(d_\tau\in M\)
be the image of the source's chosen generator
\(\mathrm{gen}(D.\mathrm{host},D.\mathrm{embE},(D.\mathrm{eTau}(\tau))_1)\).
The second \texttt{Field\allowbreak OK} clause gives
\[
                 \ell_T(D,M,\tau)=\mathcal O_Md_\tau.
 \tag{GCF.29}
\]
Indeed the source's \(\ell_T\) is generated by the images of
\(a b\), where \(a\in A\) and \(b\) belongs to the original
different ideal, with the stated embedding at \(\tau\).
This is exactly \(\mathrm{ellLat}\).  The generator identity in
the coefficient field carries to (GCF.29) under the isomorphism
\(e\); in each direction, express an element as a finite integral
linear combination and map every coefficient.  The ideal is nonzero,
so \(d_\tau\ne0\).

For \(s:D.\mathrm{LL}\hookrightarrow\mathbb C\), put
\(\rho=s\circ e^{-1}\).  The source's \(\mathrm{conjLat}\)
and the specified \(\mathrm{lamConj}(M,\rho,\Lambda)\) both
have discriminant one and the same \(j\)-invariant.  PA1 and TA1
then show they differ by a homothety \(\zeta\) with
\(\zeta^{12}=1\).  To check the last normalization, equality of
\(j\) and of the discriminants gives equality of the cubes of
\(g_2\) and squares of \(g_3\).  The scalar calculation in GCF7
provides a simultaneous homothety, also when one invariant is zero;
uniqueness of the lattice with its two invariants is TA1.
Its absolute value is one.  Their areas are therefore equal.

The original \texttt{true\allowbreak Gx} is defined using the area of
\(\mathrm{conjLat}\) and the \emph{same} chosen generator.
Its positive-value guard is inactive, since the area is positive
and \(\rho(d_\tau)\ne0\).  Thus it obeys the exact identity
\[
 \mathrm{trueGx}.\Omega(\tau,s)
   =\left\{
       \frac{\pi\,|\rho(d_\tau)|^2}
       {\operatorname{area}(\mathrm{lamConj}(M,\rho,\Lambda))}
     \right\}^{-1}.                                  \tag{GCF.30}
\]
The original \(\mathrm{OmegaCompatible}\) consists precisely of
the image field \(M\), its isomorphism \(e\), its given complex
embedding, \texttt{Data\allowbreak In}, and (GCF.29)--(GCF.30) for every
\(\tau,s\).  GCF10 supplies \texttt{Data\allowbreak In}; the preceding argument
supplies every other field.  Consequently \(\mathrm{retie}(D)\)
has this complete property.  If the original later stage already has
\(G.\mathrm{gx}=\mathrm{trueGx}\), the formula holds for that
stage itself.  In the selected application \(\mathrm{ccK}=\mathcal
O_K\), and the ambient input \(I.J.\delta\) is exactly the period
root just used.  This proves the literal
\texttt{om\allowbreak Tied\_\allowbreak with\allowbreak True}, for all later \(W,V,S,N\).
\hfill\(\square\)
% END GCF11

\subsection{Constructing and preserving the same local comparison}

% BEGIN GCF12
\paragraph{GCF12. All five compatibility fields on the actual recursive record.}
\label{claim:GCF12}
Fix a supplied scale \(\mathrm{sc}\), original level-zero data
\(L_0\), forms \(T_u\), \texttt{Std\allowbreak Tied}, \texttt{Joint0},
coefficient-field level-zero characters, and the original companion
identity
\[
              L_{b,u}+\mathfrak a_u L_{b,u}^{\#}=\Lambda_0.
 \tag{GCF.31}
\]
There is one family \(g_{u,w}\) of the actual local models, with
their exact rank identities, for which replacing the \texttt{gat}
field gives an original record \(D^g\), together with
\[
 G_t:\mathrm{GatTies}(D^g),\quad c:\mathrm{LatCmp}(D^g),\quad
 \mathcal G:\mathrm{GroupCompatibilityAll}(D^g,G_t,c).
 \tag{GCF.32}
\]
Its level-zero datum is the given \(L_0\); its forms are the given
\(T_u\), transported to that record.  It also satisfies the full
\texttt{Is\allowbreak Rec\allowbreak G}, including its joint equality of
\((D^g,G_t,c,\mathcal G.\mathrm{herm},\mathcal G.\mathrm{forms})\).

Here are the actual choices and all the fields.  LI1--LI8 construct
the local completion identifications \(I_{u,w}\), inverse away
characters \(\psi_{u,w}\), and the embeddings into the same
\(\mathbb C_p\).  SF1--SF7 use (GCF.31) to choose one
\(B\in\mathrm{GL}_k(E)\) with
\[
 B^*JB=1,\qquad
 x\in\Lambda_0\ \Longleftrightarrow\ B^{-1}x\in\mathcal O_E^k,
 \tag{GCF.33}
\]
and make every local frame the specified base change of that \(B\).
Take the original \(\mathrm{packC}\) and \(\mathrm{gatOf}\)
with these choices, with \(C_e=1\), and with the original
\(r_{\mathrm{sc}}\) and \(\mathrm{ep}_{\mathrm{sc}}\).
CR1--CR6 give the exact equality of its local coefficient rank
with \(\dim_{\mathbb C}(D.\mathrm{mem}(u).A.Mu0)\).
Consequently replacing the \texttt{gat} field is legitimate and is
exactly \(\mathrm{recC}(D,\mathrm{sc},1,L_0,I,\psi,h_d)\).

Take \(\mathcal G.\mathrm{herm}\) to be \(L_0\) with the
canonical local ties of this construction, and take the transported
\(T_u\) as \(\mathcal G.\mathrm{forms}\).  VE8--VE9 provide all
six original \texttt{Cmp\allowbreak Hyp0P} fields on these same objects:
polynomial value existence, evaluation injectivity, point existence,
all presentations, inverse-character identity, and integral lattice
span.  Define \(c\) by those unique value polynomials and the
original coefficient embedding.  The exact equality is
\[
                \mathcal O_w\text{-span}\,c(\mathcal L_u)=M_{u,w}.
 \tag{GCF.34}
\]
Its \texttt{Cmp\allowbreak Tied} is now (GCF.33), the common-frame identities,
and these original value/evaluation/character identities, with
\(G_t\)'s original positions.  Thus the third compatibility field
is supplied by a proved map, including its integral span.

For the fourth field the four rings in
\texttt{Ring\allowbreak Compatible\allowbreak At} are already the canonical
\(WQ(p,n),QQ(p,n),ORec(D,w,n),KfRec(D,w,n)\).
Use the four identity isomorphisms.  They commute with the two
fraction-field embeddings, the maps \(\iota\) and \(j\), and the
specified image of \(w_U\), because these are the maps used in
\(\mathrm{packC}\).  This proves all five commuting equations
for every original \(u,w,v\); no isomorphism with a different
coefficient ring is substituted.  For the fifth field use the actual
\texttt{Hida\allowbreak Setting}, original orbit number, dominant weights and
rank \(k\), the chosen \(I_{u,w},\psi_{u,w}\), and its literal
\(\mathrm{gatOfLevelP}\) equality.  The \texttt{Hyp\allowbreak At\allowbreak P} condition
holds because these split local models carry the original self-dual
base lattice by the specified frames, as LI5 and SF6 show.
This is precisely \texttt{Gat\allowbreak Is}.

The original recursive witness lists exactly these \(D,L_0,I,\psi,
h_d,B,i_K,\lambda\) and the six-field comparison just constructed.
All five entries in its dependent equality are the definitions in
(GCF.32), so it supplies \texttt{Is\allowbreak Rec\allowbreak G}, rather than just equality
of their ranks.  LI6 also supplies, for every original position
\((v,s)\) and every \(x\in\mathcal O_F\),
\[
 x\in\mathfrak p_v
 \ \Longleftrightarrow\
  0<\nu_w\bigl(\tau_{G_t(v,s)}(x)\bigr).             \tag{GCF.35}
\]
It includes \(x=0\) and uses the original normalized valuation and
coefficient embedding.  Finally, the recursive numerical fields are
\[
 r=\begin{cases}r_G(\mathrm{sc})&p<p_0(\mathrm{sc},k),\\1&
                                  p\ge p_0(\mathrm{sc},k),
       \end{cases}\qquad
 e=\left\lceil f k^2\log_p a_{\mathrm{sc}}\right\rceil_+.
 \tag{GCF.36}
\]
These are the two original \texttt{Gat\allowbreak Num} equalities, with the
unchanged \(C_e=1\).  This proof constructs the local models jointly
with the comparison.  It does not assert that an arbitrary previously
chosen \(D.\mathrm{gat}\) or \(c\) equals these constructed data.
\hfill\(\square\)
% END GCF12

% BEGIN GCF13
\paragraph{GCF13. The local choices precede every later stage.}
\label{claim:GCF13}
Fix the original \(B,C\), and let \(Z,Z'\) be any two of its
five-part \texttt{Later} records.  All local choices in GCF12 can be
kept when moving from \(Z\) to \(Z'\).  In particular the original
\texttt{Is\allowbreak Local\allowbreak Family} for a chosen \(g\) transports to \(Z'\)
with that same \(g\), the same \(B,I,\psi,i_K,\lambda\), and
the transported level-zero forms.

The relevant dependency check is concrete.  The host, coefficient
field, its places and embeddings, ranks, word and class sets, carried
lattices, level-zero groups and characters are fields of the fixed
base and columns.  They do not depend on the later
\(W,V,S,N,G\).  The Frobenius choice \(h_P\), orbit positions,
local weights \(\lambda_{\rm rec}\), and coefficient embedding
\(\tau_{\rm rec}\) read only those fields.  Hence they too are
identical.  The level-zero form transport keeps each underlying
vector space, value function and evaluation map, and merely presents
its level and chart records over the new \(D\).
Every equation of \texttt{Is\allowbreak Local\allowbreak Family} therefore has the same two
sides after transport: the standard \(B\), its local frames,
\(i_K\mathrm{embRec}=\lambda\mathrm{embLL}\), all position
identities, the inverse away character, and \(g=\mathrm{gatOf}\).

For the full symmetric columns, the same carrier and form data define
\(\mathrm{colsRe}_{Z'}(C)\).  The endpoint and middle
\texttt{Reads} equations, residue-level identities, and actual value
functions are unchanged, so their \(\mathrm{StdC}\) and every
\(\mathrm{StdLvC}\) remain true.  Taking degree zero commutes
exactly with transport:
\[
  (\mathrm{colsRe}_{Z'}C).\mathrm{forms}(u).\mathrm{to0}
              =\mathrm{forms0Re}_{Z'}(C.\mathrm{forms}(u).\mathrm{to0}).
 \tag{GCF.37}
\]
When only the \texttt{gat} field is replaced by the chosen \(g\),
this is exactly the \(\mathrm{forms0Set}\) used in GCF12.
The value polynomials continue to have the same universal evaluation
formula; VE7's injectivity makes their equality literal, rather than
an unidentified linear equivalence.  Thus the common local choices
need not be repeated after selecting a row stage, division stage,
metric stage, or global group stage.
\hfill\(\square\)
% END GCF13

% BEGIN GCF14
\paragraph{GCF14. One actual \(G\) supplies the full original \(\mathrm{GIs2}\).}
\label{claim:GCF14}
For any original column-choice property \(P\), the literal statement
\[
 \mathrm{StepW}\ P\
  \left(\lambda K,B,f,T,w.\ \exists G:\mathrm{GJStage}(B,w.C),\
                                  \mathrm{GIs2}(K,w,G)\right)
 \tag{GCF.38}
\]
holds.  In particular it holds with \(P=\mathrm{Pinned}\), which
is the original endpoint
\[
 \texttt{exists\_\allowbreak GJStage\_\allowbreak GIs2\_\allowbreak pinned}.
\]
Its witnesses preserve the given coefficient field, embeddings,
column data and ambient period root.

Fix arbitrary objects quantified by \texttt{Step\allowbreak W}, including its
\(K.\mathrm{Over}\) and \texttt{Forced} guards, and a given
\(w:\mathrm{ColW}(B,f,T)\).  The conclusion does not use any
additional clause of \(P\).  The \texttt{Over} data give the
original \texttt{Has\allowbreak Letters}; LS22 and HF19 construct an initial
\(Z_0=(W_0,V_0,S_0,N_0,G_0)\) over these same \(B,w.C\).
TG12 gives its actual \(C_0:\mathrm{ColsTiedSym}\), \(\mathrm{StdC}\),
and all \(\mathrm{StdLvC}\).  Its companion identity is
\[
           L_{b,u}+(\bar\pi_u)L_{b,u}^{\#}=\Lambda_0
\]
by TG12, using the principal ideal generated by the actual
\(\bar\pi_u\).  Its \texttt{Joint0} is part of these columns,
and FC16 supplies the coefficient-field level-zero character with
the original additive-character values.  The local construction in
GCF12 therefore chooses one \(g\) and its exact rank identities
\(h_d\), at scale \(K.s\).

Choose, now and once,
\[
         G=\{G_0.\mathrm{withTrue}
                        \text{ with }\mathrm{gat}=g,
                                      \mathrm{gat\_\allowbreak d}=h_d\}.
 \tag{GCF.39}
\]
The \(\mathrm{withTrue}\) operation keeps its original \(cd\)
and all other stage data, uses the specified true volume, and sets
\(gx=\mathrm{trueGx}\).  GV19 proves \texttt{Vol\allowbreak Tie\allowbreak G} for it,
with the original Haar normalization and original true-volume
foundation retained there.  Replacing its \texttt{gat} field does
not change that volume identity.

Only after (GCF.39), let \(W,V,S,N\) be arbitrary later stages,
and put \(Z=(W,V,S,N,G)\).  GCF13 transports \(C_0\) and all its
standardness data to this \(Z\), and keeps the chosen local
\(I,\psi,B,i_K,\lambda,g\).  GCF12 then supplies
\((G_t,c,\mathcal G)\), the joint \texttt{Is\allowbreak Rec\allowbreak G}, and the exact
prime identity (GCF.35), with the transported forms and the same
level-zero datum.  GCF6 extends the symmetric columns using
\emph{this} \(\mathcal G.\mathrm{herm.g}\) as its prescribed
group-tie field.  Its level and local ties are therefore the two
entries of \(\mathcal G.\mathrm{herm}\), giving the exact equality
\[
 \mathcal G.\mathrm{herm}=F_u.\mathrm{to0},\qquad
 \mathcal G.\mathrm{forms}(u)=C.\mathrm{forms}(u).\mathrm{to0}.
 \tag{GCF.40}
\]
The \texttt{gpl} clause is (GCF.35) on those same ties.  Thus
\(C,\mathrm{StdC},\mathrm{StdLvC},F_u,\mathrm{StdDt},G_t,c,
\mathcal G\), (GCF.40), and \texttt{gpl} are every original
\texttt{Col\allowbreak T} field.

The \texttt{Col\allowbreak W} input supplies \texttt{Field\allowbreak OK},
\(\mathrm{ccK}=\mathcal O_K\), and the exact twelfth-root law for
\(w.I.J.\delta\).  GCF11 supplies \texttt{Omega\allowbreak Compatible} for
this \(Z\).  Adding it and the already constructed \texttt{Is\allowbreak Rec\allowbreak G}
to \texttt{Col\allowbreak T} gives the entire original
\texttt{Column\allowbreak Data\allowbreak Split\allowbreak Primes}(K,w,Z).  Finally (GCF.36) gives
\texttt{Gat\allowbreak Num}.  We have proved, with \(G\) preceding every
later variable,
\[
 \begin{gathered}
 \forall W,V,S,N,\qquad
       \mathrm{VolTieG}(w.DZ)\ \land\ \mathrm{GatNum}(K.s,w.DZ)\\
       {}\land\ \operatorname{Nonempty}
                  (\mathrm{ColumnDataSplitPrimes}(K,w,Z)).
 \end{gathered}                                                   \tag{GCF.41}
\]
This is exactly \(\mathrm{GIs2}\).  No Bessel nonvanishing,
row-count formula, theta-kernel identity or arbitrary-family
compatibility assumption was inserted into (GCF.38).
\hfill\(\square\)
% END GCF14

% BEGIN GCF15
\paragraph{GCF15. Arbitrary division-lattice changes preserve the comparison map.}
\label{claim:GCF15}
Let \(N_u\) be arbitrary \(\mathcal O_{D.\mathrm{LL}}\)-submodules
of the original division vector spaces and let \(h_u\in\mathbb R\)
be arbitrary.  Let \(D'=\mathrm{setDiv}(D,N,h)\), which changes
only the designated division-lattice and height fields.  Then the
actual recursive comparison satisfies
\[
 \mathrm{cmpTr}_{N,h}(\mathrm{latCmpC}(T,G_t,I,H))
   =\mathrm{latCmpC}(T',G_t',I',H').                  \tag{GCF.42}
\]
All symbols on the right are the original transports, not new
independently chosen forms or comparisons.  No finite-generation,
full-rank or height bound is imposed on the newly inserted \(N_u,h_u\).

The source \texttt{Cmp\allowbreak In} contains the coefficient embedding,
integrality equivalence, common \(B\), its standard-coordinate
identity and the local frame identities.  None of these mentions
the changed fields.  The coefficient polynomial rings, original
form vectors and the local lattice tuples are also unchanged.
Consequently the six \texttt{Cmp\allowbreak Hyp0P} assertions
\(\mathrm{ex},\mathrm{evInj},\mathrm{ptEx},\mathrm{ptAll},
\mathrm{chi},\mathrm{latM}\) are the same assertions on those
same elements.  Copying them proves the transported hypothesis.
The derived \texttt{Cmp\allowbreak Hyp} assertions of value existence and
uniqueness, positions, characters, linearity and integral span
transport for the same reason.

For an actual lattice tuple, the map on either side of (GCF.42)
is the coefficient embedding of the unique polynomial satisfying
the VE7 evaluation formula.  The formula and all its arguments are
unchanged, so the polynomials agree.  On a nonlattice tuple both
maps have the source's zero value.  The coefficient embeddings
themselves agree.  These observations prove equality of the full
\texttt{Lat\allowbreak Cmp} records, rather than only equality of their images
or ranks.

Suppose now \(\mathrm{IsRecG}(\mathrm{sc},D,G_t,c,H_0,F)\).
Take its original recursive witness over \(D_0\).  Replace \(D_0\)
by \(\mathrm{setDiv}(D_0,N,h)\), retaining the same levels,
local identifications, away characters, rank equality, matrix \(B\),
completion embeddings and position identities.  Copy its six-field
comparison hypothesis as above.  Equation (GCF.42) proves the
comparison entry of the required dependent equality; the other four
entries are the original transports.  This gives the literal
\[
 \mathrm{IsRecG}(\mathrm{sc},D',G_t',c',H_0',F').        \tag{GCF.43}
\]
It is the full original
\texttt{concrete\allowbreak Split\allowbreak Prime\allowbreak Automorphic\allowbreak Data\_\allowbreak division\allowbreak Lattice\allowbreak Change}.
\hfill\(\square\)
% END GCF15

% BEGIN GCF16
\paragraph{GCF16. The complete column compatibility survives that same change.}
\label{claim:GCF16}
The operation of GCF15 also transports the full original
\(\mathrm{GroupCompatibilityAll}\), \(\mathrm{ColT}\), and
\(\mathrm{ColumnDataSplitPrimes}\), keeping their actual
\(G_t,c\) as the specified \(\mathrm{gatTr},\mathrm{cmpTr}\).

The Hermitian spaces, local groups, split-place maps, coefficient
rings, class representatives and form evaluations do not depend on
the changed division submodules.  Therefore the five fields
\(\mathrm{herm},\mathrm{forms},\mathrm{cmp},\mathrm{ring},
\mathrm{gat}\) copy over directly; in the comparison field use
the equality of actual maps in GCF15.  The full Hermitian column
record copies its determinant map and all its laws, every level
and residue map, its volume and relative determinant.  Its forgotten
columns are still the transported original columns.  Every standard
reading, the Hermitian and form equalities, and \texttt{gpl} have
unchanged arguments.  These are all the \texttt{Col\allowbreak T} fields.

For \texttt{Column\allowbreak Data\allowbreak Split\allowbreak Primes}, (GCF.43) supplies its
\texttt{is\allowbreak Rec} field.  The \texttt{om} field also copies: the host,
complex embeddings, \(gx\), the period lattice, its coefficient
data and its one-dimensional different-ideal modules are all
unchanged.  This proves the exact
\texttt{column\allowbreak Local\allowbreak Data\allowbreak Transport\allowbreak Division}, including both extra
fields.  In particular later arithmetic choices of the division
lattice cannot force a fresh choice of a comparison or a fresh
coefficient field.  This statement concerns the named transports;
it is not a construction of every other field of a later
\texttt{Extra} record.
\hfill\(\square\)
% END GCF16

% BEGIN GCF17
\paragraph{GCF17. Equality with the given extra comparison is a joint equality.}
\label{claim:GCF17}
For every original supplied \(q:\mathrm{WPt}(\mathcal S,K)\),
the literal conclusion of \texttt{WPt.\allowbreak extra\allowbreak Tied\_\allowbreak col} holds:
\[
 \begin{gathered}
  \mathrm{ColsExtraTiedS}(q.x.\mathrm{col.full},q.O.\mathrm{ex}),\\
  \exists\mathcal G:\mathrm{GroupCompatibilityAll}
                (q.w.D(q.Z),q.O.\mathrm{gatT},q.O.\mathrm{ex.cmp}),\\
  \mathcal G=q.x.\mathrm{col.g},\quad
  \mathcal G.\mathrm{herm}=q.x.\mathrm{col.full.to0},\quad
  \mathcal G.\mathrm{forms}(u)=q.x.\mathrm{col.C.forms}(u).\mathrm{to0}.
 \end{gathered}                                                   \tag{GCF.44}
\]
The equality of \(\mathcal G\) is read after the original dependent
identifications.  The assertion is for the given \(q.O\), including
its supplied comparison and kernel.

The \texttt{Wit} in \(q\) contains the exact equality
\[
 (O.\mathrm{gatT},O.\mathrm{cmp},O.\mathrm{tw},O.\mathrm{ker})
       =(\mathrm{col}.G_t,\mathrm{col}.c,w.\mathrm{tw},\mathrm{ker}_0)
 \tag{GCF.45}
\]
in the dependent space in which the kernel is indexed by the first
three entries.  Substitute this equality once, simultaneously.  Use
\(\mathrm{col}.g\) as \(\mathcal G\); its existing Hermitian and
form equalities are now precisely the last two equations in
(GCF.44).  The \texttt{Cols\allowbreak Extra\allowbreak Tied\allowbreak S} level and error-datum fields
are \(\mathrm{Wit.lev}\) and \(\mathrm{Wit.errX}\), while its
prime compatibility is \(\mathrm{col.gpl}\).  Its two kernel
embedding clauses are the actual \(\mathrm{KerT.iw}\) and
\(\mathrm{KerT.iw\_\allowbreak pos}\), transported through (GCF.45).
Thus every field is on the same given kernel and comparison.

The resulting family implication has an exact boundary.  Suppose a
given family is eventually realized by such \(q\)'s with the
original joint equality of its \((D,O)\), for every family index
and every later admissible height witness.  Then (GCF.44) gives
\(\mathrm{GroupCompatibilityAll}(\mathrm{Fm}.D,
\mathrm{Fm}.\mathrm{gatT},\mathrm{Fm}.\mathrm{ex.cmp})\) at
those same points.  Taking the maximum of the realization threshold
and \(\mathrm{Fm}.x_1+1\) gives exactly
\(\mathrm{GroupCompatibilityFamily}\).  Likewise the original
\texttt{TSel\allowbreak F}, if supplied, provides its actual \texttt{Rec\allowbreak T}
at every point, whose \(\mathrm{TiedDataS.g}\) is already the
same compatibility record; the same maximum gives the original
\texttt{g\allowbreak Tied\allowbreak F\_\allowbreak of\_\allowbreak sel}.

GCF14 constructs the complete column input used by these witnesses,
and GCF16 preserves it through division-lattice choices.  Neither
of those conclusions alone supplies the other \texttt{Wit} fields:
the selected rows, kernel/theta and divisor ties, objects, gauge,
and the joint equality (GCF.45) must still be realized in the common
construction.  Their existence is an explicit remaining producer,
not a consequence of choosing separate compatible records at the
individual places.
\hfill\(\square\)
% END GCF17

% BEGIN GCF18
\paragraph{GCF18. The selected design class already pays for its top field.}
\label{claim:GCF18}
For the original class
\[
 D_3(C_0,\mathcal A)=\mathrm{Des}(\mathrm{WithChN}(\mathcal A),C_0),
\]
membership supplies \(\exists C_H\,\mathrm{ChosenE}(\mathrm{Fm},C_0,C_H)\).
That record contains the actual selected \(\mathrm{ChosenTT}\)
and one positive constant \(C_b\) for \(\mathrm{ChosenBad}\).
Hence it supplies \(\mathrm{TopsF}\) on the unchanged selected
subset.  GLG16 proves its full original \(\mathrm{Adm2F}\), with
the unchanged explicit threshold
\[
 x_0=\max\{x_a,x_b,\mathrm{Fm}.x_1+1,
                              (16C_b+2)^{3/\epsilon}\}.
 \tag{GCF.46}
\]
No new selected subset or independent probability experiment occurs.
In particular \(\mathrm{Adm2F}\) is no longer a separate producer
once membership in this selected design class is given.

If the original selector \(\mathrm{TSel}(\mathcal A)\) is
supplied, GCF17 gives the first \texttt{Closed} input required by
LGS13 on \(D_3(3,\mathcal A)\).  The preceding projection gives
the second \texttt{Closed} input for \(\mathrm{Adm2F}\).
These are class-level statements: they hold for every family in the
class and are available before the later family and height choices.
Take the pointwise maximum of their original coefficient bounds and
the finite bounds in LGS13.  Apply that claim on the same family.
This yields exactly the original
\[
 \mathrm{Closed}\bigl(D_3(3,\mathcal A),\mathrm{LinkG}\bigr),
 \tag{GCF.47}
\]
with all five \(\mathrm{LinkG}\) fields, its error
\(b_o=b^{x^{-\epsilon/2}}\), and its unchanged
\(C_G=4C_{\rm count}+1\).  The original
\texttt{link\allowbreak G\_\allowbreak frame} follows with its literal subsequent
\(D_3,\mathrm{TuSize},\mathrm{DGamSize}\) guards; it introduces
no stronger coefficient quantifier.

The antecedent \(\mathrm{TSel}(\mathcal A)\), or an exact eventual
realization sufficient for GCF17, is retained in (GCF.47).
The completed supplier in this chapter is the full earlier
\(\mathrm{GIs2}\), its determinant and period inputs, and its
literal comparison transport.  No whole \texttt{GLOBAL\_\allowbreak LOCAL\_\allowbreak GAIN}
closure, no other global quantitative leaf, and no complete RH or GRH
proof is asserted here.
\hfill\(\square\)
% END GCF18

\endgroup
\end{document}
